\documentclass[11pt,a4paper]{article}
\usepackage[margin=2.4cm]{geometry}
\usepackage{amsmath,amssymb,amsthm,mathtools}
\usepackage[protrusion=true,expansion=false]{microtype}
\usepackage{booktabs,enumitem}
\usepackage{placeins}
\usepackage{tikz}
\usetikzlibrary{arrows.meta,positioning,calc,fit,backgrounds}
\definecolor{tierspend}{RGB}{200,60,60}
\definecolor{tierfull}{RGB}{210,140,40}
\definecolor{tierin}{RGB}{60,130,180}
\definecolor{tierrecv}{RGB}{70,150,90}
\usepackage[colorlinks=true,linkcolor=blue!50!black,citecolor=blue!50!black,urlcolor=blue!50!black]{hyperref}

\emergencystretch=3em
\hbadness=10000
\hfuzz=2pt

\newtheorem{theorem}{Theorem}[section]
\newtheorem{lemma}[theorem]{Lemma}
\newtheorem{proposition}[theorem]{Proposition}
\newtheorem{corollary}[theorem]{Corollary}
\theoremstyle{definition}
\newtheorem{definition}[theorem]{Definition}
\newtheorem{construction}[theorem]{Construction}
\newtheorem{assumption}[theorem]{Assumption}
\newtheorem{remark}[theorem]{Remark}
\newtheorem{example}[theorem]{Example}

\newcommand{\F}{\mathbb{F}}
\newcommand{\Z}{\mathbb{Z}}
\newcommand{\N}{\mathbb{N}}
\newcommand{\G}{\mathbb{G}}
\newcommand{\E}{\mathbb{E}}
\newcommand{\PP}{\mathbb{P}}
\renewcommand{\Pr}{\mathrm{Pr}}
\newcommand{\Fp}{\F_p}
\newcommand{\Fq}{\F_q}
\newcommand{\ip}[2]{\langle #1,\, #2 \rangle}
\newcommand{\Comm}{\mathsf{Commit}}
\newcommand{\Open}{\mathsf{Open}}
\newcommand{\Setup}{\mathsf{Setup}}
\newcommand{\Verify}{\mathsf{Verify}}
\newcommand{\negl}{\mathsf{negl}}
\newcommand{\poly}{\mathsf{poly}}
\newcommand{\PPT}{\mathsf{PPT}}
\newcommand{\Adv}{\mathcal{A}}
\newcommand{\Prov}{\mathcal{P}}
\newcommand{\Ver}{\mathcal{V}}
\newcommand{\Sim}{\mathcal{S}}
\newcommand{\Ext}{\mathcal{E}}
\newcommand{\Rel}{\mathcal{R}}
\newcommand{\Lang}{\mathcal{L}}
\newcommand{\nfsym}{\mathsf{nf}}
\newcommand{\cmsym}{\mathsf{cm}}
\newcommand{\cmx}{\mathsf{cmx}}
\newcommand{\cvsym}{\mathsf{cv}}
\newcommand{\rk}{\mathsf{rk}}
\newcommand{\Pallas}{\ensuremath{\mathsf{Pallas}}}
\newcommand{\Vesta}{\ensuremath{\mathsf{Vesta}}}
\newcommand{\Ep}{\ensuremath{\mathcal{E}}}
\newcommand{\NoteCommit}{\ensuremath{\mathsf{NoteCommit}}}
\newcommand{\Sinsemilla}{\ensuremath{\mathsf{Sinsemilla}}}
\newcommand{\ShortCommit}{\ensuremath{\mathsf{ShortCommit}}}
\newcommand{\CommitIvk}{\ensuremath{\mathsf{Commit}^{\mathsf{ivk}}}}
\newcommand{\Extract}{\ensuremath{\mathsf{Extract}}}
\newcommand{\Acc}{\ensuremath{\mathsf{Acc}}}
\newcommand{\GroupHash}{\ensuremath{\mathsf{GroupHash}}}
\newcommand{\ask}{\ensuremath{\mathsf{ask}}}
\newcommand{\aksym}{\ensuremath{\mathsf{ak}}}
\newcommand{\nksym}{\ensuremath{\mathsf{nk}}}
\newcommand{\ivk}{\ensuremath{\mathsf{ivk}}}
\newcommand{\fvk}{\ensuremath{\mathsf{fvk}}}
\newcommand{\ovk}{\ensuremath{\mathsf{ovk}}}
\newcommand{\dk}{\ensuremath{\mathsf{dk}}}
\newcommand{\pkd}{\ensuremath{\mathsf{pk_d}}}
\newcommand{\gd}{\ensuremath{\mathsf{g_d}}}
\newcommand{\rcv}{\ensuremath{\mathsf{rcv}}}
\newcommand{\rcm}{\ensuremath{\mathsf{rcm}}}
\newcommand{\rivk}{\ensuremath{\mathsf{r_{ivk}}}}
\newcommand{\rsk}{\ensuremath{\mathsf{rsk}}}

\renewenvironment{abstract}{%
  \small\begin{center}{\bfseries\abstractname}\end{center}\medskip\noindent\ignorespaces
}{\par\medskip}

\title{\textbf{\Huge Crosslink Guide}\\[6pt]\large Trailing finality for Zcash: the construction, its formal models, and the wallet-visible contract}
\author{m@rek.onl}
\date{}
\begin{document}
\maketitle
\thispagestyle{empty}
\begin{abstract}
Assuming the Orchard protocol of the \emph{Ironwood Guide} and the wallet layer
of the \emph{Wallet Guide}, this volume of \emph{The Zcash Arboretum}
documents Crosslink, the trailing-finality construction targeted at Zcash:
the ebb-and-flow model and what its theorems actually prove, the construction
as the Trailing Finality Layer book specifies it, stake and delegation as far
as any source pins them down, and the wallet-visible contract two ledgers
create. The design is in motion and its sources disagree --- the
implementation drops mechanisms the book calls central, and the
machine-checked model formalises a neighbouring rule set --- so every claim
carries its classification, every divergence is stated rather than resolved,
and the questions a future ZIP must answer are listed as such.
\end{abstract}
\clearpage
\tableofcontents
\clearpage
\section{Scope, status, and sources}\label{sec:xl-scope}

Crosslink is a hybrid proof-of-work/BFT consensus construction for Zcash:
it retains the existing best-chain (PoW) protocol and adds a trailing
finality layer (TFL), a BFT protocol whose blocks and the best-chain
blocks commit to one another --- the eponymous cross-links. The current
version is Crosslink~2 (tfl-book @ \texttt{fe6e1d6},
\texttt{src/design/crosslink.md} lines 1--3;
\texttt{the-arguments-for-bounded-availability-and-finality-overrides.md}
line~9). The construction descends from the ebb-and-flow security model
and the Snap-and-Chat construction of Neu, Ta\c{s}, and Tse
(``Ebb-and-Flow Protocols: A Resolution of the Availability--Finality
Dilemma'', eprint 2020/1091; arXiv 2009.04987, v1 2020-09-10, v3
2021-02-04; published IEEE S\&P 2021). Snap-and-Chat's two-ledger
architecture Crosslink keeps; its show-stopping defect for Zcash ---
users cannot in general spend outputs of finalized transactions taken
from a rolled-back fork's snapshot, because chain validity in
$\Pi_{\mathrm{lc}}$ cannot depend on a node's local finalized ledger
--- is what the TFL book records as forcing the redesign: ``We are
forced to conclude that the chain $\mathsf{ch}_i^t$ must include the
information needed to calculate $\mathsf{LOG}_{\mathrm{fin}}$''
(tfl-book @ \texttt{fe6e1d6}, \texttt{notes-on-snap-and-chat.md} lines
73--107). No successor paper covering Crosslink itself has appeared
(arXiv listing checked 2026-07-15).

Nothing about Crosslink is deployed on any Zcash network and nothing is
consensus-specified: the ZIP repository at \texttt{6961098} (2026-07-09)
contains no Crosslink, TFL, or NU8 ZIP or draft, and the implementers'
own ZIP draft is a Google document intended ``to enter the formal ZIP
process as the prototype approaches maturity'' (grep of zcash/zips @
\texttt{6961098}; zebra-crosslink @ \texttt{6d02a1b},
\texttt{book/src/design.md} line~7). This volume is therefore a
design-stage volume in the sense the \emph{Tachyon Guide} established
(\S``Scope, sources, and method''): it fixes the current state of a
moving design against commit-pinned sources, classifies every claim by
rigor, and asserts nothing normatively. The design corpus splits three
ways --- a construction book that has been static since 2024-12-13, an
implementation whose design has since moved past the book in documented
ways, and a machine-checked model of one narrow slice --- and the first
duty of this section is to say which document is authoritative for what.

\subsection{The best-chain protocol as a black box}\label{sec:xl-blackbox}

This series does not document legacy consensus internals, and this
volume does not break that rule for proof of work. Equihash, difficulty
adjustment, and block relay never appear here; the best-chain protocol
enters only through the interface Crosslink's model assumes of any
$\Pi_{\mathrm{bc}}$ (tfl-book @ \texttt{fe6e1d6},
\texttt{construction.md} lines 245--281): a set of chains rooted at a
genesis block; a bc-block-validity rule depending only on the block and
its ancestors; block production hard unless selected in proportion to
resources; a score function strictly increasing along a chain
(accumulated work, for PoW); honest nodes adopting a highest-scoring
valid chain; and headers that commit to blocks, with header validity
(PoW and difficulty adjustment) checkable before block download. The
book's own assessment is that ``Typically, Bitcoin-derived best chain
protocols do not need much adaptation to fit into this model''
(\texttt{construction.md} line~281).

The safety assumptions Crosslink places on this black box are Prefix
Consistency at confirmation depth $\sigma$ (PSS2016's ``consistency'')
and Prefix Agreement at depth $\sigma$, both stated as unconditional
properties of executions; the probabilistic reasoning that justifies
them is deliberately separated out, with NTT2020 Theorem~2 supplying the
corresponding bound $e^{-\Omega(\sqrt{\sigma})}$ on rollback failure
(tfl-book @ \texttt{fe6e1d6}, \texttt{construction.md} lines 285--374,
351--357). We state these assumptions formally in the construction
section and use them as the interface; we never derive them from PoW
internals. Where the interface touches deployed reality we cite the
deployed code: Zebra bounds its best non-finalized chain at
$\texttt{MAX\_BLOCK\_REORG\_HEIGHT} = 1000$ blocks, and its own comment
concedes that ``deeper reorganisations are outside Zebra's rollback
window. This is a local-only node policy; it is not part of consensus''
(zebra @ \texttt{9f3166b},
\texttt{zebra-chain/src/parameters/constants.rs} lines 20--30) --- the
standing admission that today's Zcash has no consensus-level finality at
any depth.
%% NOTE (updated 2026-07-15): Crypto51 lists ZEC at 21 GH/s, 1-hour 51%
%% cost ~$34,271, NiceHash-able 1% --- the classic >100%-rentable claim
%% no longer holds. Pool concentration is historical: ViaBTC exceeded 51%
%% in Sept 2023 (Blockworks 2023-09-21), but by mid-2026 shares are ~32%
%% ViaBTC / ~29% Foundry / ~16% F2Pool, no pool above one half. Frame
%% both as historical; re-verify before print.

\subsection{The three bins, and a proved tag}\label{sec:xl-bins}

Design-stage prose invites the failure mode the deployed-protocol
volumes never face: text that reads as specification but rests on a
draft. Following the series convention and the \emph{Tachyon Guide}
(\S``The three-bin method''), every claim in this volume is exactly one
of \emph{specified} (a pinned artifact at specification rigor fixes it;
we cite file and commit), \emph{designed but unspecified} (an informal
source describes a concrete mechanism no specification fixes; we cite
source and date), or an \emph{open problem} (no source resolves it;
where the designers acknowledge the gap, we quote them). Crosslink adds
a dimension Tachyon lacked: parts of its corpus carry actual proofs. We
therefore additionally tag a claim \emph{proved} when a peer-reviewed
theorem or a machine-checked artifact establishes it, always together
with the scope at which it was established.

The inventory, developed across the rest of the volume, maps as
follows. Proved: the NTT2020 theorems for ebb-and-flow and Snap-and-Chat
(peer-reviewed, IEEE S\&P 2021), and the Quint invariant
\texttt{prefix\_inv} at its small checked scope
(\S\ref{sec:xl-sources}). Specified (pinned, verifiable artifacts only):
the Quint \texttt{isValid} model itself and the quoted prototype and
Zebra constants. Designed but unspecified: the Crosslink~2 construction
rules and the two Assured Finality safety theorems, with book-level
prose proofs --- one of which is textually defective
(\S\ref{sec:xl-sources}), leaving that theorem effectively unproven ---
plus the zebra-crosslink staking and tokenomics design, the BFT
subprotocol instantiation (tenderlink), active-set selection, slashing,
and the implementation's no-Stalled-Mode variant. Open: precise liveness bounds and the dependence on $L$,
completion of the $\mathsf{LOG}_{\mathrm{ba}}$ safety development,
production values of $\sigma$ and $L$, and NU8 inclusion and dates
(classification synthesized from the sources cited throughout this
section).

\subsection{The source landscape}\label{sec:xl-sources}

Table~\ref{tab:xl-sources} lists the sources, pinned. All repository
facts were verified firsthand at the stated commits; web sources were
fetched 2026-07-15.

\begin{table}[htbp]
\centering
\small
\begin{tabular}{@{}lll@{}}
\toprule
Artifact & Pin / date & Role and rigor \\
\midrule
tfl-book (ECC)
  & \texttt{fe6e1d6}, 2024-12-13 & construction; static since pin \\
Neu--Ta\c{s}--Tse
  & eprint 2020/1091; S\&P 2021 & peer-reviewed lineage \\
crosslink-spec (Informal Systems)
  & \texttt{8edc3e9}, 2025-07-23 & Quint model; proof of concept \\
quint-crosslink (zmanian)
  & empty upstream (zero refs)  & nothing to examine \\
zebra-crosslink (Shielded Labs/Eiger)
  & \texttt{6d02a1b}, 2026-04-21 & prototype + own design book \\
zcash/zips
  & \texttt{6961098}, 2026-07-09 & zero Crosslink/NU8 ZIPs \\
zcash/zebra
  & \texttt{9f3166b}             & deployed $\Pi_{\mathrm{bc}}$ facts \\
crosslink-protocol-guide
  & \texttt{b77d845}, 2026-03-26 & secondary; visual tour \\
Shielded Labs web/forum/Arborist
  & fetched 2026-07-15           & status only, dated \\
\bottomrule
\end{tabular}
\caption{Primary Crosslink sources, pinned. Verification date
2026-07-15.}
\label{tab:xl-sources}
\end{table}

\paragraph{The TFL book.} The construction source of record is the ECC
Trailing Finality Layer book, an mdBook formerly deployed at
\url{https://electric-coin-company.github.io/tfl-book/} (404 since at
least 2026-07-15, following the repository's transfer to
\texttt{github.com/daira/tfl-book}) and now rendered at
\url{https://daira.github.io/tfl-book/}; the local clone
%% full hash: fe6e1d6f403f62da46c64e8f5a7db3cb188ffae2
HEAD is \texttt{fe6e1d6}, dated 2024-12-13, the merge of PR \#162
``crosslink2'' by Daira-Emma Hopwood, with no later commits
(\texttt{git log}). Prehistory: a design sketch
was presented at Zcon4 by Nate Wilcox (``Interactive Design of a Zcash
Trailing Finality Layer''); book version 0.1.0 introduced Crosslink, and
0.2.0 updated the construction and security analysis to Crosslink~2
(\texttt{src/version-history.md}). Three rigor caveats govern every
citation of it. First, the book self-describes as ``an early and
incomplete protocol design proposal'', listing as missing the PoS
subprotocol selection, issuance and supply mechanics, transaction
semantics integration, the transition plan, ZIPs, and the security and
economic analyses (\texttt{src/introduction/status-and-next-steps.md}
lines 1--30). Second, it is internally inconsistent about its own
version: \texttt{book.toml} still titles the deployed book v0.1.0 while
\texttt{src/version-history.md} records 0.2.0 (``Crosslink 2'') as
current (\texttt{book.toml} line~9 vs
\texttt{src/version-history.md} lines 3--5). Third, and most
consequentially, its rigor is uneven across chapters: the construction
chapter is specification-grade prose, but the liveness argument carries
explicit TODOs (``make that more precise''; ``more detailed argument
needed, especially for the dependence on $L$''), and the safety section
is marked ``\$TODO\$ Not updated for Crosslink 2 below this point''
immediately after it, with the
$\mathsf{LOG}_{\mathrm{fin}}$/$\mathsf{LOG}_{\mathrm{ba}}$ proofs below
that marker still written in Snap-and-Chat/Crosslink-1 notation ---
including the sanitization machinery Crosslink~2 exists to remove
(\texttt{security-analysis.md} lines 5--50, 54, 139--150). Two further
defects are findings of this volume, developed in the security-analysis
section: the proof text under ``Safety Theorem: Final Agreement of
$\Pi_{\mathrm{bft}}$ implies Assured Finality'' is a verbatim duplicate
of the Prefix Agreement proof and never uses Final Agreement
(\texttt{security-analysis.md} lines 87--91 vs 73--77), and the book
itself concedes that ``some rules, e.g.\ the Linearity rule, only
contribute heuristically to security in the analysis so far'', with the
$\mathsf{LOG}_{\mathrm{ba}}$ development trailing off mid-sentence
(\texttt{security-analysis.md} lines 275--281, 257).
%% RE-VERIFIED 2026-07-15: upstream HEAD unchanged (fe6e1d6f, repo now at
%% github.com/daira/tfl-book; rendered at daira.github.io/tfl-book) and
%% the rendered Final Agreement theorem still carries the Prefix
%% Agreement proof. Decision still pending: reconstruct the intended
%% proof (via the Final Agreement/Linearity sketch at
%% construction.md:592) or only report the defect.

\paragraph{The machine-checked model.} The repository
\texttt{crosslink-spec} (Informal Systems) is a single squashed commit
%% full hash: 8edc3e941fb45b32b7b68e3204e8df8330bcd43a
\texttt{8edc3e9} (2025-07-23, PR \#4).
It formalizes exactly one slice of Crosslink: the \texttt{isValid} check
of a BFT proposal against local PoW and BFT block stores, written in
Quint as a literate source (\texttt{validity.md}, tangled to
\texttt{src/validity.qnt} via \texttt{lmt}; a transpiled
\texttt{validity.tla} is committed) (README lines 25--32;
\texttt{validity.md} lines 1--7). Its own README fixes its standing: ``a
proof of concept'', whose logic was ``inferred \ldots{} from an
insightful Youtube talk'' and ``might not be up-to-date with the current
protocol design ideas'' (README line~27). What was actually checked is
premium material at a stated scope: the invariant \texttt{prefix\_inv}
--- the PoW history defined by BFT finalization is a path in the PoW
block store --- survived random simulation of 10{,}000 traces of 20
steps and TLC model checking in about three minutes at scope
$\texttt{confirmation\_depth}=1$ with stores of at most 8 blocks, and a
three-block finalization witness trace was produced by both simulation
and TLC (README lines 61--93; \texttt{validity.md} lines 376--421,
464--481). The model diverges from the book in ways we record as
model-versus-specification discrepancies, not protocol facts: its
\texttt{linear()} requires a strict-ancestor snapshot where the book's
Linearity rule permits repeating the parent snapshot (and the book needs
that permission for BFT liveness; the Quint source itself comments
``this actually seems allowed. TODO: perhaps remove''); it models the
BFT store as a single linear list with no forks and omits signatures,
notarization proofs, epochs, voting, bft-last-final, and every
$\Pi_{\mathrm{bc}}$-side rule; and its \texttt{add\_pow} assigns block
hashes colliding with genesis, an artifact PR \#4 fixed only on the BFT
side (\texttt{validity.md} lines 173--176, 132--137, 244--247, 270--278,
433--440 vs tfl-book \texttt{construction.md} lines 573, 611--618,
113--193). The second Quint repository, quint-crosslink (remote
zmanian/quint-crosslink), is empty upstream as well as locally ---
\texttt{git ls-remote} returns zero refs (checked 2026-07-15); there is
nothing to examine.
%% RESOLVED 2026-07-15: the add_pow genesis-hash collision is confined
%% to the root --- the first add_pow in every trace mints {hash 1,
%% parent 1}; later hashes are unique. prefix_inv cannot be spuriously
%% satisfied (non-root hashes are unique; descent from the universal
%% root is vacuous), so the model-checking result stands; but
%% sigmaConfirmed(genesis, ., 1) can be witnessed by the phantom block,
%% so coverage at the verification scope differs slightly from the
%% intended semantics. Report as a model artifact, not a soundness gap.

\paragraph{The implementation and its fork of the book.} The prototype
is zebra-crosslink, Shielded Labs' fork of Zebra built together with
%% full hash: 6d02a1b80f896d08f923e39b2505f0565efb5787
Eiger; the local clone HEAD is \texttt{6d02a1b}, dated 2026-04-21, the
merge of PR \#282. It ships its own mdBook whose
\texttt{cl2-construction.md} and \texttt{cl2-security-analysis.md} are
declared ``almost direct paste'' copies of the TFL book chapters,
committed verbatim ``for the purposes of precisely committing to a
\ldots{} precise set of book contents for a security design review'',
with the explicit note that Daira-Emma Hopwood and Jack Grigg ``have not
reviewed and do not necessarily endorse'' the changes
(\texttt{book/src/design/cl2-construction.md} lines 3--7). The fork then
diverges from the book in two documented, load-bearing ways. It drops
Stalled Mode entirely --- ``The zebra-crosslink design in this book
diverges from this text by \emph{not} including `Stalled Mode' rules''
--- relying instead on stakers redelegating to prevent hazards such as
BFT stalls (\texttt{cl2-construction.md} lines 10--11); and it removes
the proposer's outer signature on bft-blocks, conceding that ``a direct
hash or copy of the most recent BFT finality certificate is insufficient
evidence to ensure that specific bitwise copy is canonical''
(\texttt{cl2-construction.md} lines 14--15). The fork's own warnings
bound what this volume may claim for it: ``These extensions and changes
from the Crosslink 2 construction may violate some of the assumptions in
the security proofs'' and ``We are using a custom-fit in-house BFT
protocol for prototyping. Its own security properties need to be more
thoroughly verified'' (\texttt{book/src/design.md} lines 11--13). That
in-house engine is tenderlink (\texttt{git.sr.ht/\textasciitilde
azmr/stuff}, rev \texttt{0e88509}), Shielded Labs' lightweight
Tendermint implementation that replaced the earlier Malachite (Informal
Systems) integration --- at the pinned commit \texttt{Cargo.lock}
contains no malachitebft crates and \texttt{src/malctx.rs} survives only
as dead code --- with ed25519-zebra finalizer keys and BLAKE3 block
hashes (\texttt{zebra-crosslink/Cargo.toml} lines 74--79;
\texttt{src/lib.rs}; \texttt{src/chain.rs} lines 224--226); the
prototype parameters are $\sigma = 3$ and finalization gap bound $7$,
with the source stating ``No verification has been done on the security
or performance of these parameters'' (\texttt{src/chain.rs} lines
219--221). The staking design (the ``nutshell'' chapter) self-describes
as a ``provisional sketch'' and matches Shielded Labs' published staking
design post of 2025-10-27 on Orchard-only staking with revealed amounts,
five-day epochs, one-epoch unbonding, and power-of-ten stake
quantization (\texttt{book/src/design/nutshell.md} lines 13--133;
\url{https://shieldedlabs.net/crosslink-updated-staking-design/},
accessed 2026-07-15).

\paragraph{Secondary sources.} The community ``A Visual Tour of Zcash
Crosslink'' (crosslink-protocol-guide @ \texttt{b77d845}, 2026-03-26,
MIT) is a short mermaid-diagram tour with a TFL lexicon; we use it for
orientation only, never as a claim source. Incidental ZIP references
exist --- draft ZIP 218 (``25-second Block Target Spacing'', created
2026-03-13) cites Crosslink as complementary, and
draft-ecc-onchain-accountable-voting cites the TFL book's Assured
Finality definition --- but neither specifies any part of Crosslink
(zips @ \texttt{6961098}, \texttt{zip-0218.md} lines 77, 742).

\subsection{Status: milestones, testnets, and the absence of
dates}\label{sec:xl-status}

Shielded Labs' public roadmap records five completed milestones: M1
``Placeholders and Processes'' (2025-03-21), M2 ``Disconnected BFT''
(2025-05-16), M3 ``PoW Consensus'' (2025-08-28), M4 ``Staking
Transaction Semantics'' (2026-01-28), and M5 ``Pre-Productionization,
Tech Debt, Final Refactoring'' (2026-04-15). The Productionization Phase
is stated as Q2~2026 -- Q2~2027, in progress, described as deploying
Crosslink on a testnet, iterating, ``and, if there is consensus,
integrating it into the Zcash protocol''; no NU8 or mainnet activation
date is given anywhere on it (shieldedlabs.net/roadmap/, accessed
2026-07-15).

Deployment practice runs ahead of specification. FeatureNet, a series of
incentivized testnets announced 2026-04-02, launched Season~1 on
2026-04-15: miners, finalizers, and stakers earn real ZEC, with 500~ZEC
allocated across seasons (Season~1: 25~ZEC, split 40\% miners / 40\%
stakers / 20\% Dev Fund), and the announcement states that ``any
decision to activate Crosslink on mainnet will follow Zcash's existing
governance process'' (forum.zcashcommunity.com t/55210, accessed
2026-07-15). As of the 2026-06-25 Arborist call, two FeatureNet
workshops had completed, ``using the current stalled network to
strengthen tooling, improve diagnostics, and better handle rollback and
fork scenarios rather than simply resetting the network''; Workshop~3
was tentatively set for 2026-07-23, and the team was ``prototyping a
user-driven slashing tool, drafting ZIPs for the Proof-of-Stake
subsystem, and implementing protocol improvements''
(zechub.substack.com, published 2026-06-27, accessed 2026-07-15).
%% RESOLVED 2026-07-15: primary source --- Shielded Labs in forum
%% t/55210 (2026-06-27): ``the experience hasn't been smooth with BFT
%% finality stalled and participants ending up on different chains''; a
%% user-led forking tool was targeted for mid-July to get finality
%% unstalled. The June-2026 stall was an unplanned BFT finality failure,
%% not a deliberate exercise; the stalled network was then retained as a
%% diagnostic environment (zechub, 2026-06-27).

On dates this volume is deliberately silent, because the primary sources
are. The community ``NU8 Decision Framework'' thread (opened 2026-02-11)
records that no governance framework or dates exist for
consensus-change inclusion, while crediting execution: ``Shielded Labs
has executed well. Milestones 1 through 4 landed on or ahead of
schedule'' (forum.zcashcommunity.com t/54670, accessed 2026-07-15). A
secondary analyst report claims NU8 ``may be activated as early as Q4
2026''; we found no primary statement supporting it and do not repeat it
as fact (messari.io, unverified secondary).

\subsection{Which document is normative for what}\label{sec:xl-normative}

Throughout this volume, ``Crosslink~2'' means the construction as
published in the TFL book at \texttt{fe6e1d6}: its model, rules, and
safety-theorem statements are the referent of every unqualified claim,
and the book is normative for the construction --- while its security
analysis is only partly updated to the construction it accompanies, as
recorded above. ``The zebra-crosslink design'' means the implementation
design at \texttt{6d02a1b}: its book is normative for what Shielded Labs
is building, including where that contradicts the TFL book (Stalled
Mode, the outer signature); where the two conflict we present both and
tag the divergence explicitly, since the book has been static since
2024-12-13 while the implementation design has moved. The Quint model is
normative for nothing; it is evidence, at its stated scope, that one
slice of the design is coherent, and a source of recorded discrepancies.
Design goals frame both documents but bind neither: the TFL book targets
expected time-to-finality below 30 minutes, cost-of-attack for a 1-hour
rollback not reduced versus today, halting cost above 24 hours of PoW
mining revenue, and unchanged wallet (lightwalletd, full-node API) and
GetBlockTemplate-miner interfaces (tfl-book @ \texttt{fe6e1d6},
\texttt{src/design/goals.md} lines 22, 28--36, 73). The 30-minute goal
is uncheckable today because no production $\sigma$ or $L$ exists in any
source.

\subsection{Relation to the companion volumes}\label{sec:xl-siblings}

The interface-stability goal above is why the deployed-stack volumes
remain accurate under Crosslink and why this volume cites rather than
re-derives them. Anchor semantics and reorg behaviour of the commitment
tree are those of the \emph{Ironwood Guide} (\S``Proof of ownership of
\emph{some} valid note: the commitment tree'', in particular ``Anchor
choice as the tree grows''); what a wallet counts as confirmed is the
\emph{Wallet Guide}'s ZIP-315 \texttt{ConfirmationsPolicy} and expiry
treatment (\S``Transaction construction'' and \S``Broadcast, expiry, and
the transaction lifecycle''), whose trusted/untrusted depths are exactly
the policy layer a finality layer would let wallets strengthen. The
header-chain commitment a finality certificate would sit beside, and its
unbuilt light-client bridge, are the \emph{FlyClient Guide}'s
(\S``Relation to the frontier''), which owns the composition pointer in
the converse direction. The
\emph{Tachyon Guide} already records the institutional relationship
between the two frontier efforts --- formally independent, with no
interaction analysis in either direction (\S``Relation to Crosslink and
the upgrade environment''); this volume owns the converse
cross-reference.
%% RESOLVED 2026-07-15: the caveat is in the \emph{Ironwood Guide}, not
%% the Voting Guide --- ironwood-guide.tex:1016-1018 (\S``Proof of
%% ownership of \emph{some} valid note: the commitment tree'', ``Anchor
%% choice as the tree grows''): ``anchoring at least $100$ blocks deep
%% (Zcash's reorg limit) avoids this entirely''. Drop the Voting Guide
%% cross-reference; the reconciliation with Zebra's 1000-block
%% MAX_BLOCK_REORG_HEIGHT (a local-only policy, not consensus) and the
%% cross-volume erratum flag already live in
%% \S\ref{sec:xl-node-interface}.

\section{The ebb-and-flow model}\label{sec:xl-model}

Crosslink attaches a BFT finality layer to Zcash's proof-of-work chain, and
every security claim it makes is phrased against a model published before
any Zcash design work began: the \emph{ebb-and-flow} formulation of Neu,
Nusret Tas, and Tse, ``Ebb-and-Flow Protocols: A Resolution of the
Availability-Finality Dilemma'' (arXiv:2009.04987; v1 2020-09-10, v2
2020-12-28, v3 2021-02-04; the abstract page notes ``Forthcoming in IEEE
Symposium on Security and Privacy 2021''; hereafter [NTT2020]). This section
documents that model: the impossibility it responds to, its execution
environments, its two-ledger security definition, the snap-and-chat
construction and what the paper actually proves about it, and --- because
the Crosslink books adopt the model only after amending it --- the specific
points where the Trailing Finality Layer book departs. Throughout, the
best-chain protocol $\Pi_{\mathrm{bc}}$ remains a black box, per this
volume's scope rule: we record what the model assumes of it (chain
structure, confirmation by depth, header validity), never its internals.

\paragraph{Sources.} The primary text is arXiv:2009.04987v3 (accessed
2026-07-15). The Trailing Finality Layer book (\texttt{tfl-book} @
\texttt{fe6e1d6f}) cites the IACR ePrint mirror 2020/1091 by page number;
every passage the book quotes was verified identical in arXiv v3, and we
cite section, definition, and theorem numbers of v3 rather than pages.
%% PARTIAL (2026-07-15): ePrint 2020/1091 last revised 2021-02-08, four
%% days after arXiv v3; both marked S&P 2021, so v3 numbering and the
%% book's ePrint pages refer to the same full version (the appendices
%% with the proofs exist only there). Remaining: check the IEEE
%% camera-ready's theorem numbering only if it is ever cited directly.
The Shielded Labs/Eiger book (\texttt{zebra-crosslink} @ \texttt{6d02a1b8},
2026-04-21) carries the Crosslink~2 construction and security analysis as
\texttt{cl2-construction.md} and
\texttt{analysis/cl2-security-analysis.md} under
\texttt{book/src/design/}; a diff against
\texttt{tfl-book} @ \texttt{fe6e1d6f} shows only 20 added admonition
lines (14 in the construction copy, 6 in the security-analysis copy) and
no other change, none touching
the [NTT2020] references, the ebb-and-flow discussion, Prefix Agreement, or
Assured Finality --- the model presentation is textually unchanged between
the two books (diff run 2026-07-15), so we cite the \texttt{tfl-book}
chapter files throughout. The community
\texttt{crosslink-protocol-guide} (@ \texttt{b77d8456}) contains no
occurrence of ``ebb'' at all and is not a source for this section (grep,
2026-07-15). The ZIPs repository (@ \texttt{69610984}) contains no
Crosslink or NU8 trailing-finality draft; its only mentions are ZIP~218
citing ``mechanisms such as Crosslink'' in the context of faster block
times, and \texttt{draft-ecc-onchain-accountable-voting} citing the
book's Assured Finality definition (\S\ref{sec:xl-counterparts}).

\begin{remark}[Binning]\label{rem:xl-model-bins}
Claims in this section about what [NTT2020] states and proves are
verifiable against the published text and carry no bin tag; within its
model the paper's theorems are settled scholarship, and we cite them by
number. Claims about Crosslink itself remain design-stage and carry the
series' three bins (\emph{specified}, \emph{designed but unspecified},
\emph{open problem}); no normative Crosslink artifact exists (no draft ZIP
in \texttt{zips} @ \texttt{69610984}), so nothing in this section reaches
the specified bin except pinned source code. Where the book disputes the
adequacy of the paper's model, the sources disagree about the model, not
about what the paper proves within it; we record such disputes as findings
with both citations.
%% TODO: editorial decision pending --- this draft bins P1-liveness of
%% snap-and-chat as ``proved, in the single-transition GST model'' with the
%% book's adequacy caveat visible (Remark~\ref{rem:xl-gst-critique});
%% the alternative is to downgrade it outright.
\end{remark}

\begin{remark}[Terminology]\label{rem:xl-model-terms}
The name \emph{ebb-and-flow} denotes the security model and goal
(Definition~\ref{def:xl-ebb-flow}); \emph{snap-and-chat} denotes the
construction proposed in the same paper to achieve that goal
(\S\ref{sec:xl-snap-and-chat}). The distinction is the book's own
admonition (\texttt{tfl-book} @ \texttt{fe6e1d6f},
\texttt{notes-on-snap-and-chat.md}, ``Terminology''), and it matters
because Crosslink adopts the model while replacing the construction. The
book further replaces the paper's ``longest chain protocol'' by
\emph{best-chain protocol} ($\Pi_{\mathrm{bc}}$): Bitcoin and Zcash follow
the consensus-valid chain with most accumulated work, not the longest, and
[NTT2020]'s own footnote~2 does not require a longest-chain protocol
anyway (\S\ref{sec:xl-snap-and-chat}). Finally, [NTT2020] uses ``depth''
for what Bitcoin and Zcash call ``height''; and \texttt{zcashd} counts the
tip at depth $1$ where the book counts it at depth $0$ (\texttt{tfl-book},
\texttt{construction.md} line 66). We write $\Pi_{\mathrm{lc}}$ when
reporting the paper and $\Pi_{\mathrm{bc}}$ when reporting the books.
Notation for chains: for a chain $c$ and $\sigma \in \N$, the truncation
$c^{\lceil\sigma}$ removes the last $\sigma$ blocks of $c$; the relation
$x \preceq y$ says $x$ is a prefix of $y$; and $x \sim y$ (``$x$ and $y$
agree'') says $x \preceq y$ or $y \preceq x$. Subscripts
($\preceq_{\mathrm{bc}}$, $\sim_{\mathrm{bft}}$) name the chain kind.
\end{remark}

\subsection{The availability--finality dilemma}\label{sec:xl-dilemma}

The paper opens from an impossibility. Its abstract states the position in
three sentences (arXiv:2009.04987v3, Abstract): ``The CAP theorem says that
no blockchain can be live under dynamic participation and safe under
temporary network partitions. To resolve this availability-finality
dilemma, we formulate a new class of flexible consensus protocols,
ebb-and-flow protocols, which support a full dynamically available ledger
in conjunction with a finalized prefix ledger. The finalized ledger falls
behind the full ledger when the network partitions but catches up when the
network heals.''

The impossibility itself is attributed, not proved, in [NTT2020]
(\S\,I-A): ``it is impossible for any protocol to be both safe under
network partition and live under dynamic participation: individual nodes
in the network cannot distinguish between the two scenarios to act
differently. This intuition is formalized in [3] and its connection to the
CAP theorem [17] was made precise recently in [18].'' Reference [3] is
Pass and Shi, ``The sleepy model of consensus'' (ASIACRYPT 2017);
reference [17] is Gilbert and Lynch's CAP theorem exposition (SIGACT News
33(2), 2002); reference [18] is Lewis-Pye and Roughgarden, ``Resource
Pools and the CAP Theorem'' (arXiv:2006.10698; v1 2020-06-18; no journal
note on the abstract page, accessed 2026-07-15), which the book cites as
[LR2020] (\texttt{tfl-book} @ \texttt{fe6e1d6f},
\texttt{the-arguments-for-bounded-availability\allowbreak
-and-finality-overrides.md}).

The resolution is a two-ledger abstraction. A \emph{conservative} client
trusts only the finalized ledger, which is a prefix of the longer,
dynamically available ledger believed by an \emph{aggressive} client;
``This ebb-and-flow property avoids a system-wide determination of
availability versus finality and instead leaves this decision to the
clients'' ([NTT2020] \S\,I-B). Zcash wallets already practice a
client-side version of this choice at the confirmation-depth layer: the
ZIP-315 \texttt{ConfirmationsPolicy} parameterizes, per wallet, how deep
an output must be before it is treated as spendable (\emph{Wallet Guide},
\S``Confirmations, trust, and spendability''). The ebb-and-flow model
lifts the same freedom to the level of which ledger a client reads.

\subsection{Execution model and environments}\label{sec:xl-execution}

The model of [NTT2020] \S\,III-A has five cornerstones: $n$ total nodes;
time proceeding in slots with synchronized clocks (footnote~4: bounded
clock offsets can be captured as part of the network delay); a public-key
infrastructure with unique cryptographic identities; a random oracle as
the source of randomness; and a probabilistic polynomial-time adversary.
Corruption is static and Byzantine: ``Before the protocol execution
starts, the adversary gets to corrupt (up to) $f$ nodes, then called
adversarial. Adversarial nodes surrender their internal state to the
adversary and can deviate from the protocol arbitrarily (Byzantine faults)
under the adversary's control. The remaining $(n - f)$ nodes are honest
and follow the protocol as specified.''

Dynamic participation is modelled by sleeping: ``The adversary chooses,
for every time slot and honest node, whether the node is awake or asleep
in that slot \dots\ An honest node that is asleep in a slot does not
execute the protocol in that slot, and messages that would have arrived in
that slot are queued and delivered in the first slot in which the node is
awake again. Adversarial nodes are always awake.'' The paper is explicit
that the permissioned setting and dynamic participation are ``two
orthogonal aspects of the environment'' ([NTT2020] \S\,III-A).

Partial synchrony is modelled with a single asynchrony-to-synchrony
transition: ``Although in reality, multiple such periods of (a-)synchrony
could alternate, we follow the long-standing practice in the BFT
literature and study only a single such transition'' ([NTT2020]
\S\,III-A). Two adversary-environment pairs then formalize the two
regimes.

\begin{definition}[P1 environment $(\mathcal{A}_1(\beta), \mathcal{Z}_1)$;
{[NTT2020]} \S\,III-A]\label{def:xl-env-p1}
The pair $(\mathcal{A}_1(\beta), \mathcal{Z}_1)$ formalizes a partially
synchronous network under dynamic participation, with respect to the
fraction $\beta$ of adversarial nodes: the adversary $\mathcal{A}_1$
corrupts $f = \beta n$ nodes. Before a global stabilization time
$\mathsf{GST}$, the adversary can delay network messages arbitrarily;
after $\mathsf{GST}$, it must deliver all messages sent between honest
nodes within $\Delta$ slots. Before a global awake time $\mathsf{GAT}$,
the adversary determines which honest nodes are awake or asleep and when;
after $\mathsf{GAT}$, all honest nodes are awake. Both $\mathsf{GST}$ and
$\mathsf{GAT}$ are chosen by the adversary, are unknown to the honest
nodes, and can be causal functions of the randomness in the protocol.
\end{definition}

\begin{definition}[P2 environment $(\mathcal{A}_2(\beta), \mathcal{Z}_2)$;
{[NTT2020]} \S\,III-A]\label{def:xl-env-p2}
The pair $(\mathcal{A}_2(\beta), \mathcal{Z}_2)$ formalizes a synchronous
network under dynamic participation, with respect to a bound $\beta$ on
the fraction of awake nodes that are adversarial: at all times the
adversary must deliver all messages sent between honest nodes within
$\Delta$ slots; at all times it determines which honest nodes are awake or
asleep and when, subject to the constraint that at most a fraction $\beta$
of awake nodes are adversarial and at least one honest node is awake.
\end{definition}

The global awake time exists because without it the model would collapse
into the impossibility (footnote~5): ``Without slightly restricting
dynamic participation via a GAT after which all nodes are awake, this
adversary would fall under the CAP theorem so that no secure protocol
against it can exist. In many applications it is realistic that every now
and then there is a period in which all nodes are awake.'' The two
classical regimes are recovered at the extremes ([NTT2020] \S\,I-D): ``If
GAT = 0, then the environment is the classical partially synchronous
network, and the ledger $\mathsf{LOG}_{\mathrm{fin}}$ has the optimal
resilience achievable in that environment. On the other hand, if GST = 0
and GAT = $\infty$, then the environment is a synchronous network with
dynamic participation, and the ledger $\mathsf{LOG}_{\mathrm{da}}$ has the
optimal resilience achievable in that environment.''

\begin{remark}[Finding: the single-transition idealization]
\label{rem:xl-gst-critique}
The book argues at length that the DLS1988 $\mathsf{GST}$ formalization
was intended for one-shot consensus and ``cannot correctly be applied to
liveness properties of non-terminating protocols'', and that [NTT2020]'s
acknowledgement quoted above (``long-standing practice'') ``is not
adequate: \dots\ it is not valid in general to infer that properties
holding for the first transition to synchrony also apply to subsequent
transitions'' (\texttt{tfl-book} @ \texttt{fe6e1d6f},
\texttt{construction.md} lines 13--41; textually unchanged in
\texttt{zebra-crosslink} @ \texttt{6d02a1b8},
\texttt{cl2-construction.md}). This is the stated reason Crosslink~2 does
not take synchrony or partial synchrony as an implicit assumption of its
communication model. The dispute concerns the model's adequacy for
non-terminating liveness, not the correctness of the paper's proofs
within the model. Classification: the critique and the resulting
communication-model choice are \emph{designed but unspecified} (book
prose, no normative artifact).
\end{remark}

\subsection{Security definitions}\label{sec:xl-security-defs}

\begin{definition}[Ledger security; {[NTT2020]} Definition 1]
\label{def:xl-ledger-security}
Let $T_{\mathrm{confirm}}$ be a polynomial function of the security
parameter $\sigma$. A state machine replication protocol $\Pi$ outputting
a ledger $\mathsf{LOG}$ is \emph{secure after time $T$}, with transaction
confirmation time $T_{\mathrm{confirm}}$, if $\mathsf{LOG}$ satisfies:
\begin{itemize}
\item \emph{Safety:} for any two times $t \ge t' \ge T$ and any two honest
  nodes $i$ and $j$ awake at times $t$ and $t'$ respectively, either
  $\mathsf{LOG}_i^t \preceq \mathsf{LOG}_j^{t'}$ or
  $\mathsf{LOG}_j^{t'} \preceq \mathsf{LOG}_i^t$.
\item \emph{Liveness:} if a transaction is received by an awake honest
  node at some time $t \ge T$, then for any time
  $t' \ge t + T_{\mathrm{confirm}}$ and honest node $j$ awake at time
  $t'$, the transaction is included in $\mathsf{LOG}_j^{t'}$.
\end{itemize}
A protocol safe (live) after $T = 0$ is called simply safe (live). Here
$\mathsf{LOG}_i^t$ denotes the ledger in the view of node $i$ at time $t$;
for longest-chain protocols, $\sigma$ is the confirmation delay for
transactions.
\end{definition}

\begin{definition}[Flexible protocol; {[NTT2020]} Definition 2]
\label{def:xl-flexible}
A \emph{flexible protocol} is a pair of state machine replication
protocols $(\Pi_1, \Pi_2)$, where $\Pi_1$ and $\Pi_2$ have the same input
transactions $\mathsf{txs}$ and output ledgers $\mathsf{LOG}_1$ and
$\mathsf{LOG}_2$, respectively.
\end{definition}

\begin{definition}[$(\beta_1, \beta_2)$-secure ebb-and-flow protocol;
{[NTT2020]} Definition 3]\label{def:xl-ebb-flow}
A $(\beta_1, \beta_2)$-secure ebb-and-flow protocol is a flexible protocol
$(\Pi_{\mathrm{da}}, \Pi_{\mathrm{fin}})$ outputting an available ledger
$\mathsf{LOG}_{\mathrm{da}}$ and a finalized ledger
$\mathsf{LOG}_{\mathrm{fin}}$, such that, for security parameter
$\sigma := T_{\mathrm{confirm}}$:
\begin{enumerate}
\item \emph{P1 (Finality):} under
  $(\mathcal{A}_1(\beta_1), \mathcal{Z}_1)$, the ledger
  $\mathsf{LOG}_{\mathrm{fin}}$ is safe at all times, and there exists a
  constant $C$ such that $\mathsf{LOG}_{\mathrm{fin}}$ is live after time
  $C(\max\{\mathsf{GST}, \mathsf{GAT}\} + \sigma)$, except with
  probability $\negl(\sigma)$.
\item \emph{P2 (Dynamic Availability):} under
  $(\mathcal{A}_2(\beta_2), \mathcal{Z}_2)$, the ledger
  $\mathsf{LOG}_{\mathrm{da}}$ is secure, except with probability
  $\negl(\sigma)$.
\item \emph{Prefix:} for any honest node $i$ and time $t$,
  $\mathsf{LOG}_{\mathrm{fin},i}^t$ is a prefix of
  $\mathsf{LOG}_{\mathrm{da},i}^t$.
\end{enumerate}
Here $\negl(\cdot)$ decays faster than all polynomials: for all $c > 0$
there exists $\sigma_0$ such that $\negl(\sigma) < \sigma^{-c}$ for all
$\sigma > \sigma_0$.
\end{definition}

\paragraph{Optimality.} The resilience constants are sharp in the paper's
model (\S\,III-C): ``Observe that no ebb-and-flow protocol can tolerate a
Byzantine adversary $\mathcal{A}_1(\beta_1)$ with $\beta_1 \ge 1/3$ in a
partially synchronous network. Similarly, no ebb-and-flow protocol can
tolerate a Byzantine adversary $\mathcal{A}_2(\beta_2)$ with
$\beta_2 \ge 1/2$ in a synchronous network. Hence the security of
$\Pi_{\mathrm{sac}}$ implies that it is optimally resilient. We denote the
worst-case adversary-environments as
$(\mathcal{A}_1^*, \mathcal{Z}_1) := (\mathcal{A}_1(1/3), \mathcal{Z}_1)$
and
$(\mathcal{A}_2^*, \mathcal{Z}_2) := (\mathcal{A}_2(1/2), \mathcal{Z}_2)$.''
%% RESOLVED 2026-07-15: [NTT2020] SIII-A corrupts ``(up to) f'' nodes
%% with f = beta*n, and Lemma 1 (quorum intersection, from [CS2020,
%% Lemma 14]) requires strictly fewer than n/3 non-honest voters; the
%% worst-case notation A1*(1/3) is the boundary family with strictness
%% carried by ``up to'' and by the underlying lemmas. Quote the
%% resilience constants as sharp suprema (secure for beta1 < 1/3,
%% beta2 < 1/2), matching the book's explicitly strict one-third bound
%% (construction.md:138-141).

\paragraph{Relation to flexible BFT.} Ebb-and-flow goes beyond the
flexible BFT of Malkhi, Nayak, and Ren ([NTT2020] reference [20]) in two
ways ([NTT2020] \S\,I-E): dynamic participation is brought in as a client
belief, and prefix consistency between the two ledgers is required ``in
all circumstances'' rather than only when both clients' assumptions hold.
Where flexible BFT supports a tradeoff for $f$ between $n/3$ and $n/2$,
``The snap-and-chat protocol achieves $(n/2, n/3)$, simultaneously
optimal. No tradeoff is necessary'' (Fig.~3 caption). The remaining gap to
Blum, Katz, and Loss ([NTT2020] reference [27]), a non-flexible protocol
with an optimal synchrony/asynchrony tradeoff, ``can be interpreted as the
value of flexibility''.

\subsection{The snap-and-chat construction}\label{sec:xl-snap-and-chat}

\begin{construction}[Snap-and-chat; {[NTT2020]} \S\S\,I-D, III-B]
\label{def:xl-sac}
Nodes run two off-the-shelf protocols in parallel: a dynamically available
protocol $\Pi_{\mathrm{lc}}$ (footnote~2: ``Longest chain protocols are
representative members of this class of protocols, hence the notation
$\Pi_{\mathrm{lc}}$, but this class includes many other protocols as
well'') and a partially synchronous BFT protocol $\Pi_{\mathrm{bft}}$.
Each node $i$ takes snapshots of its confirmed $\Pi_{\mathrm{lc}}$ ledger
$\mathsf{ch}_i^t$ and inputs them into $\Pi_{\mathrm{bft}}$; the output
$\mathsf{Ch}_i^t$ of $\Pi_{\mathrm{bft}}$ is an ordered list --- a chain
of chains --- of snapshots, where a BFT block carries as payload only a
reference $B.\mathsf{ch}$ to an LC block identifying the snapshot. The
ledgers are extracted as
\[
\mathsf{LOG}_{\mathrm{fin},i}^t :=
  \mathsf{sanitize}(\mathsf{flatten}(\mathsf{Ch}_i^t)),
\qquad
\mathsf{LOG}_{\mathrm{da},i}^t :=
  \mathsf{sanitize}(\mathsf{flatten}(\mathsf{Ch}_i^t)
    \,\Vert\, \mathsf{ch}_i^t),
\]
where $\mathsf{flatten}$ concatenates the referenced chains of LC blocks
as ordered, $\Vert$ is concatenation, and $\mathsf{sanitize}$ removes all
but the first occurrence of each block ([NTT2020] \S\,III-B3, Fig.~5
caption). At transaction level (footnote~6): ``Formally,
$\mathsf{LOG}_{\mathrm{da}}$ and $\mathsf{LOG}_{\mathrm{fin}}$ are now
represented as sequences of LC blocks. Proper transactions ledgers are
readily obtained by concatenating the transactions contained in the blocks
and removing duplicate and invalid transactions (sanitization).''
\end{construction}

The coupling from $\Pi_{\mathrm{lc}}$ into $\Pi_{\mathrm{bft}}$ is a
boycott: ``in the $\Pi_{\mathrm{bft}}$ sub-protocol, each honest node
boycotts the finalization of snapshots that are not confirmed in
$\Pi_{\mathrm{lc}}$ in its view'' (\S\,I-D). Concretely, Streamlet's
voting rule is extended --- ``An honest node only votes for a proposed BFT
block $B$ if it views $B.\mathsf{ch}$ as confirmed'' --- and this is the
only change (line 19 of Algorithm~1; ``Sleepy is applied unaltered and the
modification required for Streamlet is minor''). ``In addition,
$\mathsf{ch}_i^t$ is used as side information in $\Pi_{\mathrm{bft}}$ to
boycott the finalization of invalid snapshots proposed by the adversary''
(Fig.~5 caption).

\paragraph{Instantiation.} The analyzed instance takes
$\Pi_{\mathrm{lc}}$ to be Sleepy consensus (Pass--Shi; permissioned
longest chain, where ``blocks of a certain depth on the longest chain are
confirmed'') and $\Pi_{\mathrm{bft}}$ to be Streamlet (epochs with
pseudo-random leaders; a block is notarized when at least two-thirds of
nodes vote for it; ``Out of three adjacent notarized blocks from
consecutive epochs, the middle one gets finalized along with its
prefix'') ([NTT2020] \S\,III-B). Section III-D extends the analysis to
HotStuff and PBFT as $\Pi_{\mathrm{bft}}$ with minor modifications, and
Theorem~\ref{thm:xl-sac-security} carries over.

\begin{theorem}[Snap-and-chat security; {[NTT2020]} Theorem 1]
\label{thm:xl-sac-security}
The protocol $\Pi_{\mathrm{sac}}$ is a $(1/3, 1/2)$-secure ebb-and-flow
protocol. (Proved under the static-corruption, synchronized-slot, PKI plus
random-oracle model of \S\ref{sec:xl-execution}: P1 under
$(\mathcal{A}_1^*(1/3), \mathcal{Z}_1)$, P2 under
$(\mathcal{A}_2^*(1/2), \mathcal{Z}_2)$, Prefix by construction; Appendix
D proves the same with HotStuff in place of Streamlet.)
\end{theorem}

\paragraph{Proof structure.} We record what is proved and how ([NTT2020]
\S\,III-C, Appendix B, dependency diagrams Figures 7 and 8, boxes 1--8),
because the book's departures target specific joints of this argument.
P1 safety is Lemma~1 --- safety of $\Pi_{\mathrm{bft}}$ by quorum
intersection, unaffected by the voting-rule change --- plus the
observation that ledger extraction preserves safety. P1 liveness is
Lemma~4, which needs Theorem~\ref{thm:xl-lc-after-gst} and Lemmas 2--3
(Streamlet liveness, each carrying the highlighted extra condition that
``the leaders' proposals reference LC blocks that are viewed as confirmed
by all honest nodes''). P2 is the security of $\Pi_{\mathrm{lc}}$ under
$(\mathcal{A}_2^*, \mathcal{Z}_2)$, taken from Pass--Shi (their Theorem 3
and Lemma 1), plus Lemma~\ref{lem:xl-ntt-l5}. Prefix holds by
construction.

\begin{theorem}[Longest-chain security after
$\max\{\mathsf{GST},\mathsf{GAT}\}$; {[NTT2020]} Theorem 2]
\label{thm:xl-lc-after-gst}
For all $p < (n - 2f) / (2 \Delta n (n - f))$, there exists a constant
$C > 0$ such that for any $\mathsf{GST}$ and $\mathsf{GAT}$ specified by
$(\mathcal{A}_1^*, \mathcal{Z}_1)$, the protocol $\Pi_{\mathrm{lc}}(p)$ is
secure after $C(\max\{\mathsf{GST}, \mathsf{GAT}\} + \sigma)$, with
transaction confirmation time $T_{\mathrm{confirm}} = \sigma$, except with
probability $e^{-\Omega(\sqrt{\sigma})}$. (Here $p$ is the probability
that a given node produces a block in a given slot; the constant $C$
depends on $p$, $n$, $f$, and $\Delta$ (footnote 9); footnote 10 improves
the error to $A_\varepsilon\, e^{-a_\varepsilon \sigma^{1-\varepsilon}}$
for any $\varepsilon > 0$ via the DKT+2020 recursive bootstrapping
argument.)
\end{theorem}

The proof of Theorem~\ref{thm:xl-lc-after-gst} extends Pass--Shi pivots to
\emph{GST-strong pivots} ([NTT2020] Definition 4 and eq.~(8)): a slot
$t \ge \max\{\mathsf{GST}, \mathsf{GAT}\}$ such that for any $t_a \le t
\le t_b$, the convergence opportunities counted only after
$\max\{t_a, \mathsf{GST}, \mathsf{GAT}\}$ outnumber the adversarial slots
in $[t_a, t_b]$ --- only post-transition convergence opportunities count,
because their useful properties fail under asynchrony or while honest
nodes may sleep.

\begin{lemma}[Consistency of $\mathsf{LOG}_{\mathrm{fin}}$ with
$\Pi_{\mathrm{lc}}$ under P2; {[NTT2020]} Lemma 5]\label{lem:xl-ntt-l5}
The ledger $\mathsf{LOG}_{\mathrm{fin}}$ is a safe prefix of the output of
$\Pi_{\mathrm{lc}}$ in the view of every honest node at all times under
$(\mathcal{A}_2^*, \mathcal{Z}_2)$, except with probability
$e^{-\Omega(\sqrt{\sigma})}$. The key proof step: for a BFT block to
become final in the view of an honest node under
$(\mathcal{A}_2^*, \mathcal{Z}_2)$, at least one vote from an honest node
is required, and honest nodes only vote for a BFT block if they view the
referenced LC block as confirmed; this holds even if the final BFT blocks
conflict in the output of $\Pi_{\mathrm{bft}}$.
\end{lemma}

\paragraph{Why the composition is safe.} Figure 9 of [NTT2020] presents
the trichotomy. (a) When both subprotocols are safe, $\mathsf{ch}$ and
$\mathsf{Ch}$ do not fork. (b) When $\Pi_{\mathrm{lc}}$ is unsafe,
``snapshots taken by different nodes or at different times can conflict.
However, $\Pi_{\mathrm{bft}}$ is still safe and thus orders these
snapshots linearly. Any transactions invalidated by conflicts are
sanitized during ledger extraction. As a result,
$\mathsf{LOG}_{\mathrm{fin}}$ remains safe.'' (c) When
$\Pi_{\mathrm{bft}}$ is unsafe, in a synchronous network with
$n/3 < f < n/2$, ``finalization of a snapshot requires at least one honest
vote, and thus only valid snapshots become finalized. Since finalized
snapshots are consistent, $\mathsf{LOG}_{\mathrm{fin}}$ is consistent with
$\mathsf{LOG}_{\mathrm{lc}}$.''

\paragraph{Simulation evidence.} The paper simulated Sleepy-plus-Streamlet
with $n = 100$ nodes, $f = 25$ adversarial, $\Delta = 1\,$s, Sleepy rate
$\lambda = 0.1\,\mathrm{s}^{-1}$ (per-node $\lambda_0 =
10^{-3}\,\mathrm{s}^{-1}$), confirmation depth $k = 20$, and Streamlet
epoch $5\,$s ([NTT2020] \S\,IV). Under dynamic participation,
$\mathsf{LOG}_{\mathrm{da}}$ grows steadily while
$\mathsf{LOG}_{\mathrm{fin}}$ stalls whenever fewer than a two-thirds
quorum (67 nodes) is awake, catching up afterwards; under intermittent
partitions (honest nodes split $2(n-f)/3$ and $(n-f)/3$), BFT finalization
stalls, the parts' $\mathsf{LOG}_{\mathrm{da}}$ drift apart, and on
reunion nodes converge on the longer $\mathsf{LOG}_{\mathrm{da}}$ while
$\mathsf{LOG}_{\mathrm{fin}}$ catches up. Simulation code:
\url{https://github.com/tse-group/ebb-and-flow} (footnote 7). The paper
offers no implementation against a deployed chain; this matters in
\S\ref{sec:xl-book-findings}.

\subsection{The proof-of-work setting}\label{sec:xl-pow-setting}

Section V-B of [NTT2020] adapts the model to the shape Crosslink actually
instantiates: nodes come in two flavors, \emph{miners} (quantified by hash
rate, powering the PoW longest chain) and \emph{validators} (unique
cryptographic identities, providing finality). The informal theorem for
this setting reads: in a network environment where (1) communication is
asynchronous until a global stabilization time $\mathsf{GST}$ after which
it becomes synchronous, (2) validators are always awake, and (3) miners
can sleep and wake at any time --- (P1, Finality) the ledger
$\mathsf{LOG}_{\mathrm{fin}}$ is safe at all times and live after
$\mathsf{GST}$, if fewer than 33\% of validators are adversarial, fewer
than 50\% of awake hash rate is adversarial, and awake hash rate is
bounded away from zero; (P2, Dynamic Availability) if $\mathsf{GST} = 0$,
the ledger $\mathsf{LOG}_{\mathrm{da}}$ is safe and live at all times, if
fewer than 66\% of validators are adversarial, fewer than 50\% of awake
hash rate is adversarial, and awake hash rate is bounded away from zero
([NTT2020] \S\,V-B).

The paper explains why the requirements exceed the permissioned case:
under P1, ``fewer than 50\% of awake hash rate have to be adversarial, and
awake hash rate has to be bounded away from zero, since otherwise
$\Pi_{\mathrm{lc}}$ might not be live and there would be no input for
$\Pi_{\mathrm{bft}}$ to finalize''; under P2, ``fewer than 66\% of
validators have to be adversarial, since otherwise adversarial validators
could finalize $\Pi_{\mathrm{lc}}$ blocks that are not on the longest
chain, which would later enter the prefix of
$\mathsf{LOG}_{\mathrm{da}}$'' ([NTT2020] \S\,V-B). The paper closes the
section with an open question, which we record verbatim: ``It remains an
open question whether these additional requirements are fundamental or a
limitation of the snap-and-chat construction.'' Classification: \emph{open
problem}, the paper's own.

\paragraph{The deployed interface today.} On the black-box side of the
boundary, the only datum this volume records about the deployed
$\Pi_{\mathrm{bc}}$ is its rollback bound: Zebra bounds its best
non-finalized chain by \texttt{MAX\_BLOCK\_REORG\_HEIGHT} ($1000$ blocks),
whose documentation states ``This is a local-only node policy; it is not
part of consensus. The window is sized as a defence-in-depth measure
against sustained consensus splits''; Zebra finalizes its oldest blocks
once the chain grows past this height (\texttt{zebra} @ \texttt{9f3166b9},
\texttt{zebra-chain/src/parameters/constants.rs:20--30},
\texttt{zebra-state/src/constants.rs:16--19}). Nothing in today's Zcash
consensus provides finality; the constant is a node policy, which is
precisely the gap Crosslink addresses.
%% TODO: reconcile with the Ironwood Guide's anchor caveat, which cites
%% ``anchoring at least 100 blocks deep (Zcash's reorg limit)''
%% (\emph{Ironwood Guide}, \S``Proof of ownership of \emph{some} valid
%% note: the commitment tree''); the 100-block figure is zcashd lineage,
%% the 1000-block constant is Zebra's. State both or neither, with crate
%% cites, after checking zcashd.

\subsection{Findings: where the book amends the paper}
\label{sec:xl-book-findings}

The book accepts the ebb-and-flow model as the starting point and then
records defects and gaps that shape Crosslink. We reproduce each finding
with both citations; all book material below is \emph{designed but
unspecified} (prose in \texttt{tfl-book} @ \texttt{fe6e1d6f}, textually
unchanged in \texttt{zebra-crosslink} @ \texttt{6d02a1b8}; no normative
artifact).

\paragraph{Scope of the one-honest-vote argument.} The trichotomy case (c)
above ``is technically correct but has to be interpreted with care. It
only applies when the number of malicious nodes $f$ is such that
$n/3 < f < n/2$. What we are trying to do with Crosslink is to ensure that
a similar conclusion holds even if $\Pi_{\mathrm{bft}}$ is completely
subverted, i.e.\ the adversary has 100\% of validators (but only
${<}\,50\%$ of $\Pi_{\mathrm{bc}}$ hash rate)''
(\texttt{notes-on-snap-and-chat.md}, ``Comment on security
assumptions''). Crosslink's safety target is therefore strictly stronger
than what Theorem~\ref{thm:xl-sac-security} case (c) delivers.

\paragraph{A hidden safety assumption in Lemma 5.} The proof of
Lemma~\ref{lem:xl-ntt-l5} assumes that ``since $\Pi_{\mathrm{lc}}$ is safe
\dots, the fact that one honest node sees that snapshot as confirmed
implies that every honest node sees the same snapshot as confirmed''
([NTT2020] Lemma 5 proof). The book replies that ``while this may be a
reasonable assumption to make for parts of the security analysis, in
practice we should always require any adversary to do the relevant amount
of Proof-of-Work to construct block headers that are plausibly confirmed''
(\texttt{notes-on-snap-and-chat.md}, admonition quoting ePrint p.~9).

\paragraph{A subtlety in sanitization.} Two readings of ledger extraction
--- concatenate transactions then sanitize, versus concatenate blocks,
deduplicate blocks, then sanitize --- ``are equivalent in the setting
considered by [NTT2020], but the argument for their equivalence is not
obvious''; the equivalence requires that the \emph{only} reasons for
contextual invalidity in $\Pi_{\mathrm{lc}}$ are double-spends and missing
inputs, and ``strictly speaking Snap-and-Chat should require of
$\Pi_{\mathrm{lc}}$ that there is no such other reason. This does not seem
to be explicitly stated anywhere in [NTT2020]''
(\texttt{notes-on-snap-and-chat.md}, ``Subtlety in the definition of
sanitization''). Zcash violates the premise: transaction expiry can be
handled (the height-bounded validity of ZIP-203; \emph{Wallet Guide},
\S``Expiry (ZIP-203)''), but missing anchors break the argument, because
a transaction invalid for a missing anchor can become valid later
(anchors are the tree roots developed in the \emph{Ironwood Guide},
\S``Proof of ownership of \emph{some} valid note: the commitment tree'').
%% RESOLVED 2026-07-15: read arXiv:2010.10447v1 (2020-10-20;
%% accountability + light clients). Its LC-block reference to the
%% last-final BFT block with a monotonicity condition anticipates the
%% Extension rule, but only as light-client metadata; sanitization-based
%% extraction is kept and chain validity never depends on the finalized
%% ledger (its own footnote 2 concedes conflicts with future prefixes).
%% Both findings stand; note the LC-reference similarity when discussing
%% the Extension rule's lineage. No NTT erratum found.

\paragraph{Spending finalized outputs.} In snap-and-chat, transactions in
$\mathsf{ch}_i^t$ must be able to spend outputs of finalized transactions
that are not in $\mathsf{ch}_i^t$'s own history --- they may come from
another fork's snapshot. Since consensus validity of the tip must be a
deterministic function of the chain itself, $\Pi_{\mathrm{lc}}$ ``must
include the information needed to calculate
$\mathsf{LOG}_{\mathrm{fin},i}^s$''. The book's assessment: ``The authors
of [NTT2020] probably just missed this: the paper only has evidence that
they simulated their construction, rather than implementing it for
Bitcoin or any other concrete block chain as $\Pi_{\mathrm{lc}}$''
(\texttt{notes-on-snap-and-chat.md}, ``Spending finalized outputs'' and
``Finalization Availability''; compare the simulation-only evidence in
\S\ref{sec:xl-snap-and-chat}). The repair adds a property the paper does
not have --- \emph{finalization availability} --- via a commitment
$\mathsf{LF}(H)$ in each block header to the last final BFT block known to
the block producer, governed by the following rule.

\begin{definition}[Extension rule; \texttt{tfl-book} @ \texttt{fe6e1d6f},
\texttt{notes-on-snap-and-chat.md}]\label{def:xl-extension}
If $H$ is not the genesis block header, then the final BFT block committed
to by $H$ either descends from or equals the final BFT block committed to
by $H$'s parent:
\[
\mathsf{LF}(H^{\lceil 1}) \preceq_{\mathrm{bft}} \mathsf{LF}(H).
\]
\end{definition}

The book states that the Extension rule ``will be preserved into
Crosslink~2'', and that it yields, by construction and regardless of
$\Pi_{\mathrm{bft}}$: if node $i$ observes no rollback deeper than
$\sigma$ in $\Pi_{\mathrm{lc}}$, then $i$ observes no instability in
$\mathsf{LOG}_{\mathrm{fin},i}$, even if $\Pi_{\mathrm{bft}}$'s
assumptions are violated (\texttt{notes-on-snap-and-chat.md},
``Finalization Availability'').

\subsection{From dynamic availability to bounded availability}
\label{sec:xl-bounded-availability}

The book's one model-level departure from ebb-and-flow concerns what
happens during a long partition. Since partition cannot be prevented, the
CAP theorem --- ``that loosely speaking is made precise by [LR2020]'' ---
implies that any finality-providing protocol may stall for as long as the
partition lasts; strict dynamic availability then makes the finalization
gap grow without bound (\texttt{tfl-book} @ \texttt{fe6e1d6f},
\texttt{the-arguments-for-bounded-availability\allowbreak
-and-finality-overrides.md}). The book argues that spending under an
unbounded gap is unacceptable, and replaces strict dynamic availability
with \emph{bounded availability} --- a weakening of dynamic availability
in the sense of [DKT2020, arXiv:2010.08154] --- plus a \emph{Stalled Mode}
(coinbase-only blocks) and explicitly governed \emph{finality overrides}.
Block production keeps running during a stall rather than halting, because
of tail-thrashing attacks: during a stall, an adversary with 10\% of hash
power can build a 10-block fork in 100 block times, and the $k$-block tail
can ``thrash'' between chains (same chapter).

\begin{definition}[Finality depth and the Stalled Mode rule;
\texttt{tfl-book} @ \texttt{fe6e1d6f},
\texttt{notes-on-snap-and-chat.md}]\label{def:xl-finality-depth}
For a block header $H$, let
\begin{align*}
\mathsf{tailhead}(H) &:=
  \mathsf{lastcommonancestor}(\mathsf{snapshot}(\mathsf{LF}(H)), H), \\
\mathsf{finalitydepth}(H) &:=
  \mathsf{height}(H) - \mathsf{height}(\mathsf{tailhead}(H)).
\end{align*}
Consensus rule (Stalled Mode variant): for every $\Pi_{\mathrm{lc}}$ block
$H$, either $\mathsf{finalitydepth}(H) \le L$, or $H$ follows the Stalled
Mode restrictions (for Zcash: coinbase-only blocks), where $L$ is the
finality gap bound.
\end{definition}

\begin{definition}[Bounded hazard-freeness; \texttt{tfl-book} @
\texttt{fe6e1d6f}]
\label{def:xl-bounded-hazard}
For a finality gap bound of $L$ blocks: there is never observed to be,
for any node $i$ at time $t$, a more-available ledger
$\mathsf{LOG}_{\mathrm{da},i}^t$ containing a hazardous transaction that
originates in a block $H$ of $\mathsf{ch}_i^t$ with
$\mathsf{finalitydepth}(H) > L$
(\texttt{notes-on-snap-and-chat.md}, Definition ``Bounded
hazard-freeness'').
\end{definition}

Classification: the bounded-availability argument, Stalled Mode, and
finality overrides are \emph{designed but unspecified} --- book prose with
no draft ZIP (\texttt{zips} @ \texttt{69610984}). The governance shape of
finality overrides in particular has no normative text.

\subsection{Crosslink 2's counterpart properties}
\label{sec:xl-counterparts}

The book restates the model's goals in a different proof style before
replacing the construction: security arguments take the form ``in an
execution where safety properties are not violated, undesirable thing
cannot happen'', with probabilistic reasoning separated out
(\texttt{tfl-book} @ \texttt{fe6e1d6f}, \texttt{construction.md} lines
342--359). On the relation to the paper: [NTT2020] Theorem 2
``essentially says that under certain conditions Prefix Agreement holds
except with probability $e^{-\Omega(\sqrt{\sigma})}$''; the bound is not
tight (footnote 10's bootstrapping via DKT+2020 \S\,4.2 brings it to
$A_\varepsilon\, e^{-a_\varepsilon \sigma^{1-\varepsilon}}$); and in the
proofs of Lemma 5 and Theorem 1 this probability ``just passes through
the proof from the premisses to the conclusions, because the argument is
not probabilistic''. The counterpart properties, in the book's
execution-based style (the star is the book's protocol-name wildcard:
$\Pi_{\star\mathrm{bc}}$ ranges over
$\{\Pi_{\mathrm{origbc}}, \Pi_{\mathrm{bc}}\}$, likewise for bft, so
each property is stated once for the original and the adapted
subprotocol alike):

\begin{definition}[Prefix Consistency; \texttt{tfl-book} @
\texttt{fe6e1d6f}, \texttt{construction.md}]\label{def:xl-prefix-consistency-ntt}
An execution of $\Pi_{\star\mathrm{bc}}$ has \emph{Prefix Consistency}
at confirmation depth $\sigma$ iff for all times $t \le u$ and all nodes
$i$, $j$ (potentially the same) such that $i$ is honest at time $t$ and
$j$ is honest at time $u$,
\[
(\mathsf{ch}_i^t)^{\lceil\sigma} \preceq_{\mathrm{bc}} \mathsf{ch}_j^u.
\]
(Taken from PSS2016's ``consistency''.)
\end{definition}

\begin{definition}[Prefix Agreement; \texttt{tfl-book} @
\texttt{fe6e1d6f}, \texttt{security-analysis.md}]
\label{def:xl-prefix-agreement-ntt}
An execution of $\Pi_{\star\mathrm{bc}}$ has \emph{Prefix Agreement} at
confirmation depth $\sigma$ iff for all times $t$, $u$ and all nodes $i$,
$j$ (potentially the same) such that $i$ is honest at time $t$ and $j$ is
honest at time $u$,
\[
(\mathsf{ch}_i^t)^{\lceil\sigma} \sim_{\mathrm{bc}}
(\mathsf{ch}_j^u)^{\lceil\sigma}.
\]
\end{definition}

\begin{definition}[Final Agreement; \texttt{tfl-book} @
\texttt{fe6e1d6f}, \texttt{security-analysis.md}]
\label{def:xl-final-agreement-ntt}
An execution of $\Pi_{\star\mathrm{bft}}$ has \emph{Final Agreement}
iff for all $\star$bft-valid blocks $C$ in honest view at time $t$ and
$C'$ in honest view at time $t'$,
$\mathsf{bftlastfinal}(C) \sim_{\star\mathrm{bft}}
\mathsf{bftlastfinal}(C')$.
\end{definition}

\begin{definition}[Assured Finality; \texttt{tfl-book} @
\texttt{fe6e1d6f}, \texttt{construction.md} line 505]
\label{def:xl-assured-finality-ntt}
An execution of Crosslink~2 has \emph{Assured Finality} iff for all times
$t$, $u$ and all nodes $i$, $j$ (potentially the same) such that $i$ is
honest at time $t$ and $j$ is honest at time $u$,
$\mathsf{fin}_i^t \sim_{\mathrm{bc}} \mathsf{fin}_j^u$.
\end{definition}

\begin{theorem}[Safety from either subprotocol; \texttt{tfl-book} @
\texttt{fe6e1d6f}, \texttt{security-analysis.md}]
\label{thm:xl-assured-either}
Prefix Agreement of $\Pi_{\mathrm{bc}}$ at confirmation depth $\sigma$
implies Assured Finality; and, independently, Final Agreement of
$\Pi_{\mathrm{bft}}$ implies Assured Finality. Safety of Crosslink~2's
main goal therefore holds under reasonable assumptions about
\emph{either} subprotocol.
\end{theorem}

The book notes that its $\mathsf{LOG}_{\mathrm{fin}}$ safety proof from
Prefix Agreement ``does not depend on any safety property of
$\Pi_{\mathrm{bft}}$, and is an elementary proof. ([NTT2020, Theorem 2]
is a much more complicated proof that takes nearly 3 pages, not counting
the reliance on results from [PS2017].)'' This is the book's answer to
the one-honest-vote scope finding of \S\ref{sec:xl-book-findings}: the
target regime includes a completely subverted $\Pi_{\mathrm{bft}}$.

\begin{theorem}[Ledger prefix property; \texttt{tfl-book} @
\texttt{fe6e1d6f}, \texttt{construction.md} lines 522--536]
\label{thm:xl-ledger-prefix-ntt}
For any node $i$ honest at time $t$, and any confirmation depth $\mu$,
$\mathsf{fin}_i^t \preceq (\mathsf{ba}_\mu)_i^t$, where
$(\mathsf{ba}_\mu)_i^t = (\mathsf{ch}_i^t)^{\lceil\mu}$ if
$\mathsf{fin}_i^t \preceq (\mathsf{ch}_i^t)^{\lceil\mu}$, and
$\mathsf{fin}_i^t$ otherwise. This is the analogue of the Prefix property
of Definition~\ref{def:xl-ebb-flow}, and it holds by construction of
$(\mathsf{ba}_\mu)_i^t$.
\end{theorem}

The definitions above are shapes for comparison with [NTT2020]; the
constructions behind them are canonical in this volume's construction
section --- $\mathsf{fin}$ and $\mathsf{ba}_{\mu}$ in
\S\ref{sec:xl-ledgers}, the block structure in \S\ref{sec:xl-structure},
and $\operatorname{bft\text{-}last\text{-}final}$ with the validity
rules in \S\ref{sec:xl-bft} and \S\ref{sec:xl-bc}.

\paragraph{Prefix Consistency and partitions.} When Prefix Consistency is
assumed of PoW protocols, it implicitly rules out partitions of the honest
network longer than about $\sigma$ block times; GKL2020 ``proves'' it only
because the required no-partition assumption is implicit in the
(partially) synchronous communication model. BFT protocols, by contrast,
can be safe under full asynchrony --- which is why Crosslink~2's
communication model makes no implicit synchrony assumption, ``or else we
could end up with a protocol with weaker safety properties than
$\Pi_{\mathrm{origbft}}$ alone'' (\texttt{tfl-book} @ \texttt{fe6e1d6f},
\texttt{construction.md} lines 317--333; compare
Remark~\ref{rem:xl-gst-critique}).

\paragraph{What Crosslink 2 does not inherit.} The model section must not
imply that Crosslink~2 inherits [NTT2020]'s P2 as proved. The book itself
notes that its ``$\mathsf{LOG}_{\mathrm{ba}}$ does not roll back past the
agreed $\mathsf{LOG}_{\mathrm{fin}}$'' theorem ``is not as strong as we
would like'', and of the bound $L$ on non-Stalled-Mode blocks after the
finalization point, ``we have also not proven that so far''
(\texttt{security-analysis.md}). Classification: Crosslink~2's dynamic
availability analogue is an \emph{open problem} in the book's own
accounting; its safety counterparts
(Theorem~\ref{thm:xl-assured-either}) are \emph{designed but
unspecified}, proved in book prose but fixed by no normative artifact.
The Assured Finality definition already circulates beyond the book: the
ZIPs repository's \texttt{draft-ecc-onchain-accountable-voting} cites it
(\texttt{zips} @ \texttt{69610984}, line 169; the draft's placement in
the voting landscape is documented in the \emph{Voting Guide},
\S``Lineage'').

\subsection{Context: Gasper, Casper FFG, and the successor line}
\label{sec:xl-model-context}

The paper's motivating negative result is a balancing attack on Gasper
(Ethereum's then-candidate consensus) in the standard synchronous model
with adversarial --- not stochastic --- network delay, showing Gasper is
not live and that its block proposal mechanism is rendered unsafe; for
large $n$, any positive adversarial fraction $f/n$ suffices, with the
first opportune epoch after about $n/f$ epochs and the per-epoch launch
probability approximately $\beta$ ([NTT2020] \S\S\,II, V-A, Appendix A).
Appendix E recapitulates the bouncing attack on Casper FFG and concludes
that the ``bidirectional interaction'' between a finality gadget and its
blockchain ``is intricate to reason about and a gateway for liveness
attacks'', in contrast to the ``completely decoupled'' confirmation and
finalization of snap-and-chat.

A finding follows for this volume: the book acknowledges Appendix E's
warning while Crosslink~2 itself ``involves a bidirectional dependence
between $\Pi_{\mathrm{bft}}$ and $\Pi_{\mathrm{bc}}$'' --- the book takes
the bouncing attack as a lesson to be addressed, not a prohibition
(\texttt{tfl-book} @ \texttt{fe6e1d6f}, \texttt{security-analysis.md}
lines 15--20). The Extension rule of Definition~\ref{def:xl-extension} and
the finality-depth rule of Definition~\ref{def:xl-finality-depth} are both
instances of that dependence.

\paragraph{The successor paper.} The same authors identified a second
dilemma: no protocol can provide both accountable safety and liveness
under dynamic participation. Their resolution, \emph{accountability
gadgets}, checkpoints a longest-chain protocol using any BFT protocol as
a black box (``The Availability-Accountability Dilemma and its Resolution
via Accountability Gadgets'', arXiv:2105.06075; latest v3 2022-05-18;
accessed 2026-07-15). We record it here for attribution and for the
outlook; whether Crosslink's design draws on it is a question for the
design-lineage section, not this one.

\paragraph{Deployment status.} As of this writing, Crosslink is
implementation-stage, not consensus-stage: Shielded Labs completed
Crosslink Milestone 5 on 2026-04-15 and pivoted to long-running
``seasonal'' FeatureNets; an incentivized FeatureNet paying real ZEC
rewards had run for roughly six weeks as of the 2026-05-28 Arborist call;
hardening is slated for 2026, followed by productionization (ZIP
finalization and security audits), with mainnet inclusion pending
community approval. Sources, all accessed 2026-07-15:
\begin{itemize}[leftmargin=*]\small
\item \url{https://zechub.substack.com/p/arborist-call-zcash-protocol-updates-ad8}
  (2026-05-28);
\item \url{https://zechub.substack.com/p/arborist-call-zcash-protocol-updates-330}
  (2026-04-30);
\item
  \url{https://forum.zcashcommunity.com/t/shielded-labs-crosslink-deployment-updates/49706};
\item \url{https://shieldedlabs.net/crosslink-roadmap-q1-2025/}.
\end{itemize}
The sibling volume's view of the upgrade landscape is the
\emph{Tachyon Guide}, \S``Relation to Crosslink and the upgrade
environment''.

\subsection{Machine-checked status}\label{sec:xl-machine-checked}

One machine-checked artifact exists, and its scope must be stated
precisely to avoid overclaiming. The Informal Systems Quint specification
(\texttt{crosslink-spec} @ \texttt{8edc3e94}) is a proof-of-concept model
of \emph{only} the $\mathsf{isValid}$ check for a Crosslink BFT proposal
against a single node's PoW and BFT block stores --- the Tail Confirmation
rule ($\sigma$-block tail) and the Linearity rule --- with the PoW and BFT
chains modelled abstractly (actions \texttt{add\_pow} and
\texttt{try\_add\_bft}). The README states that the logic was ``inferred
\dots\ from an insightful Youtube talk'' and ``might not be up-to-date
with the current protocol design ideas'' (\texttt{README.md}). It does
\emph{not} formalize the ebb-and-flow ledgers
$\mathsf{LOG}_{\mathrm{fin}}$/$\mathsf{LOG}_{\mathrm{da}}$ or the
P1/P2/Prefix properties of Definition~\ref{def:xl-ebb-flow}: the
machine-checked coverage of the ebb-and-flow model itself is nil.

\paragraph{What is checked.} As formalized
(\texttt{validity.md}, ``Validity Function''), the check is
$\mathsf{isValid}(\mathit{bft\_store}, \mathit{pow\_store}, p)$ requiring:
the store contains $p.\mathit{block}$; the proposal context equals
\texttt{bft\_store.last()}; \textsf{sigmaConfirmed} --- a fork-free tail
of at least $\sigma$ PoW blocks, some element of which has the proposed
block as its parent, and whose parents all lie within the tail; and
\textsf{linear}. The
invariant \texttt{prefix\_inv}, requiring for all consecutive indices $i$
\begin{center}
\texttt{pow\_store.isDescendant(bft\_store[i].pow\_hash,
bft\_store[i+1].pow\_hash)}
\end{center}
and glossed ``The PoW history defined by finalization is a path in the
PoW blockstore'', passed random
simulation (10000 traces of 20 steps, no violation) and TLC model
checking after transpiling to TLA+ on the verification scope (stores of
up to 8 blocks, $\mathtt{confirmation\_depth} = 1$, about three minutes).
The property \texttt{finalization\_witness} generates traces with 3 BFT
blocks, and \texttt{run finalizationTest} is an assertion-checked
concrete execution. Files: \texttt{src/validity.qnt} (tangled from
\texttt{validity.md} via \texttt{lmt}), \texttt{src/validity.tla},
\texttt{src/validity.cfg}; instances: verification with
$\mathtt{confirmation\_depth} = 1$ and store limit 8, simulation with
$\mathtt{confirmation\_depth} = 2$ and no limit (\texttt{README.md}
\S\S\,2--3; \texttt{validity.md}, ``Invariants''/``Tests'', instances
block lines 463--482).

\paragraph{Finding: divergences from the book.} Function \textsf{linear}
requires
\begin{center}
\texttt{work\_store.isDescendant(context.pow\_hash, b.hash)}\\
\texttt{and not (context.pow\_hash == b.hash)},
\end{center}
the second conjunct carrying the specification's own comment
``\texttt{//~this actually seems allowed. TODO: perhaps remove}'' ---
that is, it is \emph{stricter} than the book's Linearity
rule $\mathsf{snapshot}(B^{\lceil 1}_{\mathrm{bft}})
\preceq_{\mathrm{bc}} \mathsf{snapshot}(B)$, which permits equality.
Additionally, $\mathsf{isValid}$ requires
$p.\mathit{bftContext} = \texttt{bft\_store.last()}$ (proposals only on
the current BFT tip), an abstraction not present as such in the book
(\texttt{crosslink-spec} @ \texttt{8edc3e94}, \texttt{validity.md} lines
131--176; \texttt{tfl-book} @ \texttt{fe6e1d6f},
\texttt{construction.md}, Linearity rule).
%% TODO: resolve with the spec authors which semantics is intended for
%% Crosslink 2 (strict ancestor vs preceq; tip-pinned proposals) before
%% citing the Quint model as confirming the book.

The repository \texttt{quint-crosslink}
(\texttt{github.com/zmanian/quint-crosslink}) contains no commits locally,
and \texttt{git ls-remote} against the upstream returns no refs --- the
repository is empty upstream as well (checked 2026-07-15). The Quint
ground truth is \texttt{crosslink-spec} in its entirety.
%% TODO: search the informalsystems and ShieldedLabs organizations for
%% further Quint models before finalizing the ``what is machine-checked''
%% claim.

\paragraph{Classification.} The Quint artifact itself is
\emph{specified} in the only sense available before a protocol
specification: verifiable, machine-checked source at a pinned commit. The
claim it checks (\texttt{prefix\_inv}) is far weaker than the model
properties of this section; every ebb-and-flow-level property of
Crosslink~2 remains \emph{designed but unspecified} or \emph{open}, per
\S\ref{sec:xl-counterparts}.

\section{The Crosslink construction}\label{sec:xl-construction}

Crosslink couples Zcash's proof-of-work best-chain protocol to a
Byzantine-fault-tolerant (BFT) protocol so that transactions become final
some time after they first appear in PoW blocks --- the property the
Trailing Finality Layer book names \emph{Trailing Finality}, and whose
target is \emph{Assured Finality}: the assurance that transactions cannot
be reverted \emph{by that protocol}. The book deliberately eschews
``absolute finality'', because out-of-band forks can always revert
anything (tfl-book @ \texttt{fe6e1d6}, \texttt{src/terminology.md}:11--13,
45). The two-ledger shape --- a conservative finalized ledger beside an
available one --- descends from the ebb-and-flow paradigm of Neu, Tas, and
Tse, ``Ebb-and-Flow Protocols: A Resolution of the Availability-Finality
Dilemma'' (arXiv 2009.04987, v3 of 2021-02-04; IEEE S\&P 2021; accessed
2026-07-15).

The current version of the construction is \emph{Crosslink~2}: the book's
construction chapter defines it as such, and version-history entry 0.2.0
records the update of both the construction and its security analysis from
Crosslink~1 (\texttt{construction.md}:1--3;
\texttt{src/design/crosslink.md}:3; \texttt{src/version-history.md}).
%% TODO: the rendered tfl-book site --- now at daira.github.io/tfl-book
%% after the repository's transfer to github.com/daira/tfl-book; the ECC
%% github.io URL 404s (2026-07-15) --- still shows "v0.1.0" in page
%% titles while the repository version history says 0.2.0; flag when
%% citing the rendered site rather than the repository.
Throughout this section, a citation of the form \texttt{construction.md}:$n$
means \texttt{src/design/crosslink/construction.md} in the ECC Trailing
Finality Layer book at commit \texttt{fe6e1d6}; other files of that
repository are cited by path at the same commit.

Crosslink~2 modifies a BFT protocol $\Pi_{\mathrm{origbft}}$ and a
best-chain protocol $\Pi_{\mathrm{origbc}}$ into two coupled protocols
$\Pi_{\mathrm{bft}}$ and $\Pi_{\mathrm{bc}}$, by adding structural
elements, changing validity rules, and changing honest-node behaviour. A
Crosslink~2 node must participate in both $\Pi_{\mathrm{bft}}$ and
$\Pi_{\mathrm{bc}}$; acting as a bft-proposer, a bft-validator, or a
bc-block-producer is optional, but all such actors are assumed to be
Crosslink~2 nodes (\texttt{construction.md}:103--111).

The stakes for the rest of the Zcash stack are concrete. Today
confirmation depth is policy, not property: wallets gate spendability on
heuristic depths (\emph{Wallet Guide}, \S``Confirmations, trust, and
spendability'') and recover from rollbacks by truncate-and-rescan
(\emph{Wallet Guide}, \S``Reorganisations: truncate and rescan''); an
Orchard anchor is only as settled as the chain suffix beneath it
(\emph{Ironwood Guide}, \S``Proof of ownership of \emph{some} valid note:
the commitment tree''); and the coin-weighted vote fixes its electorate at
a snapshot height that nothing today prevents from reorganising
(\emph{Voting Guide}, \S``The snapshot roots: posted, never verified'').
Crosslink~2 is the construction intended to replace those heuristics with
a protocol property.

\paragraph{Classification.} Every rule in this section is \emph{designed
but unspecified}: the book is the only source at design rigor, and no
Crosslink or trailing-finality ZIP exists --- as of zips @
\texttt{6961098} (2026-07-09) the sole mention in the ZIP corpus is a
citation footnote in ZIP~218 (``25-second Block Target Spacing'', Status:
Draft, created 2026-03-13). Two artifact classes attach \emph{specified}
status, in the pre-spec sense of verifiable pinned sources, to fragments
of the material: a machine-checked Quint model of one validity check
(\S\ref{sec:xl-quint}) and prototype code (\S\ref{sec:xl-divergences}) ---
and both diverge from the book text in ways we record. Per-subsection
Classification paragraphs refine this default.

\subsection{Model and notation}\label{sec:xl-constr-model}

The communication model takes neither synchrony nor partial synchrony as
an implicit assumption: unless otherwise specified, messages can be
arbitrarily delayed or dropped; a message is received at most once, and is
nonmalleably authenticated. The book argues at length against applying the
Global Stabilization Time formalization of partial synchrony (Dwork,
Lynch, and Stockmeyer, 1988) to liveness of non-terminating block-chain
protocols (\texttt{construction.md}:13--41). Where a subprotocol's own
analysis needs timing assumptions --- as Streamlet's liveness analysis
does below --- those assumptions are imported explicitly, not ambiently.

Throughout, subscript $\star$ ranges over the two chain flavours
$\mathrm{bc}$ and $\mathrm{bft}$; a block determines the chain of its
ancestors with itself as tip, and we pass between block and chain freely.

\begin{definition}[Chain notation]\label{def:xl-notation}
Let $C$ be a $\star$-chain and $k \in \N$.
\begin{enumerate}[label=(\roman*)]
\item \emph{Truncation.} The chain $C \lceil^{k}_{\star}$ is $C$ with the
  last $k$ blocks pruned, except that if $\mathrm{len}(C) \le k$ the
  result is the genesis chain consisting only of
  $\mathsf{Origin}_{\star}$. The block at depth $k$ in $C$ is the tip of
  $C \lceil^{k}_{\star}$. For a bft-proposal $P$ with parent block $B$ and
  $k \ge 1$, define $P \lceil^{k}_{\mathrm{bft}} \coloneqq B
  \lceil^{k-1}_{\mathrm{bft}}$
  (\texttt{construction.md}:61--63, 121).
\item \emph{Prefix and agreement.} We write $B \preceq_{\star} C$ iff the
  chain with tip $B$ is a prefix of the chain with tip $C$ (including $B =
  C$). Blocks $B$ and $C$ \emph{agree}, written $B \sim_{\star} C$, iff $B
  \preceq_{\star} C$ or $C \preceq_{\star} B$; they \emph{conflict}
  otherwise (\texttt{construction.md}:69--75).
\item \emph{Linearity.} A time series $S : I \to
  \text{$\star$-block}$ is $\star$-linear iff $t \le u$ implies $S(t)
  \preceq_{\star} S(u)$ (\texttt{construction.md}:69--75).
\end{enumerate}
\end{definition}

\begin{lemma}[Linear prefix]\label{lem:xl-linear-prefix}
If $A \preceq_{\star} C$ and $B \preceq_{\star} C$, then $A \sim_{\star}
B$.
\end{lemma}
\begin{proof}
The chain of ancestors of $C$ is $\star$-linear, and blocks $A$ and $B$
both lie on that chain (\texttt{construction.md}:77--82).
\end{proof}

\begin{definition}[Agreement on a view]\label{def:xl-agreement-view}
An execution of a protocol $\Pi$ has \emph{Agreement on the view} $V :
\mathrm{Node} \times \mathrm{Time} \to \text{$\star$-chain}$ iff for all
times $t$, $u$ and all $\Pi$ nodes $i$, $j$ (potentially the same) such
that $i$ is honest at time $t$ and $j$ is honest at time $u$, we have
$V_i^t \sim_{\star} V_j^u$ (\texttt{construction.md}:96--99).
\end{definition}

\begin{remark}[Depth conventions]\label{rem:xl-depth}
The book's ``depth'' differs from both neighbouring conventions: NTT2020
uses \emph{depth} for what the book calls height, and the book's count
differs by one from the zcashd confirmation-depth convention --- the tip
sits at depth $0$, not $1$ (\texttt{construction.md}:65--67; arXiv
2009.04987). The confirmation-depth defaults quoted from the
\emph{Wallet Guide} in \S\ref{sec:xl-parameters} use the wallet
convention.
\end{remark}

\subsection{The two black boxes}\label{sec:xl-blackboxes}

Consistent with the scope of this series, both subprotocols enter the
construction only through their interfaces. We record what Crosslink~2
assumes of each; we do not derive the internals of either --- in
particular, PoW block production remains a black box.

\paragraph{Requirements on $\Pi_{\mathrm{origbc}}$.} Block producers are
selected by a resource-proportional random process. A function
$\mathsf{score} : \text{bc-valid-chain} \to \mathrm{Score}$ is fixed, with
$<$ a strict total order on $\mathrm{Score}$ and, for any non-genesis
bc-valid-chain $H$, $\mathsf{score}(H \lceil^{1}_{\mathrm{bc}}) <
\mathsf{score}(H)$; for PoW instantiations the score is accumulated work.
Honest nodes choose a highest-scoring valid chain as their best chain,
with any tie-breaking rule. Headers commit to blocks, and header validity
is necessary but not sufficient for block validity
(\texttt{construction.md}:245--290, 255--257). The model also fixes a
predicate $\mathsf{iscoinbaseonlyblock} : \text{bc-block} \to
\mathrm{boolean}$, true iff the block has exactly one transaction and that
transaction is a coinbase transaction --- one that only distributes newly
issued funds and has no inputs (\texttt{construction.md}:269--271).

The two chain-quality properties the analysis consumes are stated as
unconditional properties of executions; the probabilistic reasoning that
makes them plausible for a PoW protocol is deliberately separated from the
construction (\texttt{construction.md}:245--290, 337--366).

\begin{definition}[Prefix Consistency]\label{def:xl-prefix-consistency}
An execution of $\Pi_{\mathrm{bc}}$ has \emph{Prefix Consistency} at
confirmation depth $\sigma$ iff for all times $t \le u$ and all nodes $i$,
$j$ (potentially the same) such that $i$ is honest at time $t$ and $j$ is
honest at time $u$, we have $\mathsf{ch}_i^t \lceil^{\sigma}_{\mathrm{bc}}
\preceq_{\mathrm{bc}} \mathsf{ch}_j^u$ (\texttt{construction.md}:287--290;
the property is PSS2016 \S3.3 ``consistency'').
\end{definition}

\begin{definition}[Prefix Agreement]\label{def:xl-prefix-agreement}
An execution of $\Pi_{\mathrm{bc}}$ has \emph{Prefix Agreement} at
confirmation depth $\sigma$ iff it has Agreement on the view $(i, t)
\mapsto \mathsf{ch}_i^t \lceil^{\sigma}_{\mathrm{bc}}$
(\texttt{construction.md}:337--340).
\end{definition}

\paragraph{Requirements on $\Pi_{\mathrm{origbft}}$.} The BFT protocol
proceeds in epochs with proposals, ballots, and a \emph{notarization
rule}. The notarization rule must require at least a two-thirds absolute
supermajority of the epoch's voting units to have been cast for a
proposal $P$, and may require other conditions; a notarization proof can
be obtained only with knowledge of ballots collectively satisfying the
rule (\texttt{construction.md}:129).

A \emph{$\star$bft-block} consists of $(P, \mathsf{proof}_P)$ re-signed by
the same proposer using a strongly unforgeable signature scheme. It is
bft-block-valid iff $P$ is bft-proposal-valid, $\mathsf{proof}_P$ validly
proves that the notarization rule was satisfied by ballots cast for $P$,
and the proposer's outer signature on $(P, \mathsf{proof}_P)$ is valid. A
bft-proposal's parent reference hashes the \emph{entire} parent bft-block:
proposal, proof, and outer signature (\texttt{construction.md}:156--161).
%% TODO: book-internal TODO at construction.md:168 --- check the use of
%% the unique honestly-submitted block per epoch in the liveness analysis;
%% the outer signature is what makes that uniqueness enforceable.

A voting unit is cast \emph{non-honestly} iff it is cast other than by its
holder (for example after key compromise), it is double-cast for distinct
proposals, or its holder should not have cast that vote according to the
honest-voting conditions in its view (\texttt{construction.md}:131--154).

\begin{definition}[One-third bound on non-honest
voting]\label{def:xl-one-third}
An execution of $\Pi_{\mathrm{bft}}$ has the \emph{one-third bound on
non-honest voting} property iff for every epoch, strictly fewer than one
third of the total voting units for that epoch are \emph{ever} cast
non-honestly (\texttt{construction.md}:138--141).
\end{definition}

The bound quantifies over ballots cast at \emph{any} time in the entire
execution. Consequently the validator keys of honest nodes must remain
secret indefinitely, and rotated-out keys must be securely deleted
(\texttt{construction.md}:131--154; tfl-book issue \#140).

Finality inside the BFT protocol is a function of context, never an
objective predicate: it is correct to talk about the last final block
\emph{of a chain}, but not to call any bft-block objectively bft-final
(\texttt{construction.md}:171--187).

\begin{definition}[Required properties of
$\operatorname{bft\text{-}last\text{-}final}$]\label{def:xl-last-final}
Function $\operatorname{bft\text{-}last\text{-}final} :
\text{bft-block} \to \text{bft-block-or-}\bot$ outputs, for a
bft-block-valid context $C$, the last ancestor of $C$ that is final in the
context of $C$. Required:
$\operatorname{bft\text{-}last\text{-}final}(C) = \bot$ iff $C$ is not
bft-block-valid; if $C$ is bft-block-valid, then
$\operatorname{bft\text{-}last\text{-}final}(C) \preceq_{\mathrm{bft}} C$
(hence is itself bft-block-valid), and for all bft-valid-blocks $D$ with
$C \preceq_{\mathrm{bft}} D$,
$\operatorname{bft\text{-}last\text{-}final}(C) \preceq_{\mathrm{bft}}
\operatorname{bft\text{-}last\text{-}final}(D)$. Also
$\operatorname{bft\text{-}last\text{-}final}(\mathsf{Origin}_{\mathrm{bft}})
= \mathsf{Origin}_{\mathrm{bft}}$ (\texttt{construction.md}:177--183).
\end{definition}

\begin{definition}[Final Agreement]\label{def:xl-final-agreement}
An execution of $\Pi_{\mathrm{bft}}$ has \emph{Final Agreement} iff for
all bft-valid blocks $C$ in honest view at time $t$ and $C'$ in honest
view at time $t'$, we have $\operatorname{bft\text{-}last\text{-}final}(C)
\sim_{\mathrm{bft}} \operatorname{bft\text{-}last\text{-}final}(C')$. A
block is \emph{in honest view} if a party observes it at some time at
which that party is honest (\texttt{construction.md}:200--206).
\end{definition}

\paragraph{Worked instance: adapting Streamlet.} The book's running
$\Pi_{\mathrm{origbft}}$ is an adapted Streamlet. Its notion of
``notarized'' is identified with bft-block-validity by having the
proposer aggregate a two-thirds-supermajority proof into a submitted
block, making notarization objectively and feasibly verifiable, so that
valid chains consist of notarized blocks; the liveness timing analysis
then assumes a notarized block is received within two epochs, and
progress needs five consecutive epochs with honest leaders. Streamlet's
finality rule --- three adjacent blocks with consecutive epoch numbers
finalize the chain up to the second --- becomes
$\operatorname{origbft\text{-}last\text{-}final}$: for an
origbft-valid-block $C$,
$\operatorname{origbft\text{-}last\text{-}final}(C)$ is the last
origbft-valid-block $B \preceq_{\mathrm{origbft}} C$ such that either $B =
\mathsf{Origin}_{\mathrm{origbft}}$ or $B$ is the second block of a group
of three adjacent blocks with consecutive epoch numbers
(\texttt{construction.md}:210--243, 224, 234--241).

\paragraph{Classification.} \emph{Designed but unspecified} throughout;
the requirements are the book's model choices, argued but not proven
against any deployed BFT protocol. A specification would need to fix the
concrete $\Pi_{\mathrm{origbft}}$, its voting-unit assignment, and the
notarization proof format.

\subsection{Structural additions}\label{sec:xl-structure}

Crosslink~2 adds three structural elements
(\texttt{construction.md}:417--421):
\begin{enumerate}[label=(\arabic*)]
\item each bc-header has, in addition to the origbc-header fields, a
  $\mathsf{context\_bft}$ field that commits to a bft-block;
\item each bft-proposal has a $\mathsf{headers\_bc}$ field containing a
  sequence of exactly $\sigma$ bc-headers, zero-indexed and deepest
  first;
\item each non-genesis bft-block likewise has a $\mathsf{headers\_bc}$
  field with exactly $\sigma$ bc-headers; the genesis bft-block has
  $\mathsf{headers\_bc} = \mathsf{null}$.
\end{enumerate}

\begin{definition}[Snapshot]\label{def:xl-snapshot}
For a bft-block or bft-proposal $B$:
\[
\mathsf{snapshot}(B) \coloneqq
\begin{cases}
\mathsf{Origin}_{\mathrm{bc}} & \text{if } B.\mathsf{headers\_bc} =
  \mathsf{null},\\
B.\mathsf{headers\_bc}[0] \lceil^{1}_{\mathrm{bc}} & \text{otherwise}
\end{cases}
\]
(\texttt{construction.md}:423--430). The snapshot is the parent of the
deepest supplied header, so the $\sigma$ headers sit directly on top of
it.
\end{definition}

\begin{definition}[$\mathsf{LF}$ and the finalization
candidate]\label{def:xl-lf-candidate}
For a bc-block $H$:
\begin{align*}
\mathsf{LF}(H) &\coloneqq
  \operatorname{bft\text{-}last\text{-}final}(H.\mathsf{context\_bft}),\\
\mathsf{candidate}(H) &\coloneqq
  \operatorname{last\text{-}common\text{-}ancestor}\bigl(
    \mathsf{snapshot}(\mathsf{LF}(H)),\; H \lceil^{\sigma}_{\mathrm{bc}}
  \bigr).
\end{align*}
When $H$ is the tip of a node's bc-best-chain, $\mathsf{candidate}(H)$
gives the candidate finalization point, subject to the local no-rollback
condition of $\mathsf{updatefin}$
(Construction~\ref{def:xl-updatefin};
\texttt{construction.md}:431--438).
\end{definition}

\paragraph{Why full headers.} The $\mathsf{headers\_bc}$ field
demonstrates to any party seeing a proposal or block that the snapshot had
been confirmed when the proposal was made. Full headers, rather than
hashes, let a validator check Proofs-of-Work and difficulty adjustment
locally \emph{before} downloading blocks, limiting denial-of-service;
nodes may trim headers overlapping with ancestor bft-blocks. The genesis
special case --- $\mathsf{headers\_bc} = \mathsf{null}$ with the snapshot
patched to $\mathsf{Origin}_{\mathrm{bc}}$ --- avoids a hash cycle between
the two genesis blocks (\texttt{construction.md}:440--454).

\paragraph{Why only $\mathsf{context\_bft}$.} No separate
$\mathsf{final\_bft}$ field is needed in bc-headers: the context block
alone determines its last final ancestor, and
$\operatorname{bft\text{-}last\text{-}final}$ is efficiently computable
for typical BFT protocols (\texttt{construction.md}:456--460).

\paragraph{Classification.} \emph{Designed but unspecified}. A
specification would need to pin the commitment encoding of
$\mathsf{context\_bft}$, the header wire format, and the trimming rule.

\subsection{The BFT side: $\Pi_{\mathrm{bft}}$}\label{sec:xl-bft}

\begin{definition}[$\Pi_{\mathrm{bft}}$ validity]\label{def:xl-bft-validity}
The genesis rule is that $\mathsf{Origin}_{\mathrm{bft}}$ is
bft-block-valid (\texttt{construction.md}:568). A bft-proposal
(respectively, non-genesis bft-block) $B$ is bft-proposal-valid
(respectively, bft-block-valid) iff \emph{all} of:
\begin{description}
\item[Inherited origbft rules.] The corresponding
  origbft-proposal-validity (respectively, origbft-block-validity) rules
  hold for $B$; these are assumed to include the \emph{Parent rule} that
  $B$'s parent is bft-valid.
\item[Linearity rule.] $\mathsf{snapshot}(B \lceil^{1}_{\mathrm{bft}})
  \preceq_{\mathrm{bc}} \mathsf{snapshot}(B)$.
\item[Tail Confirmation rule.] The headers $B.\mathsf{headers\_bc}$ form
  the $\sigma$-block tail of a bc-valid-chain.
\end{description}
(\texttt{construction.md}:570--577.)
\end{definition}

Finality-in-context is unchanged from $\Pi_{\mathrm{origbft}}$:
$\operatorname{bft\text{-}last\text{-}final}$ is defined the same way as
$\operatorname{origbft\text{-}last\text{-}final}$
(\texttt{construction.md}:621--623).

\paragraph{Objective validity versus honest voting.} The validity rules
are separated from the honest voting condition on purpose: even a
proposal voted for by $100\%$ of validators is not bft-proposal-valid
unless it satisfies the rules, and a purportedly valid bft-block carrying
a valid notarization proof is not recognized by \emph{any} honest node if
it fails the other bft-block-validity rules. This separation is essential
to keeping the finalized chain safe against a complete break of
$\Pi_{\mathrm{bft}}$ --- even of the validator signature algorithm --- as
long as $\Pi_{\mathrm{bc}}$ remains safe
(\texttt{construction.md}:581--585).

\paragraph{The Linearity rule.} Within a bft-valid-chain, the referenced
bc-snapshots cannot roll back. The rule replaces and implies Crosslink~1's
Increasing Score rule; it prevents attacks that propose bc-chains forked
from much earlier (lower-difficulty) blocks; it bounds the disruption to
the bounded-available ledger even if $\Pi_{\mathrm{bft}}$ is subverted,
because the finalized chain \emph{is} a $\Pi_{\mathrm{bc}}$ chain; and it
eliminates the Snap-and-Chat/Crosslink~1 ``sanitization'' step entirely.
The rule is relative to the parent bft-block's snapshot even when that
parent is not yet final in the current context
(\texttt{construction.md}:587--609).
%% TODO: book-internal TODO at construction.md:608 --- desideratum on PoS
%% leader selection (PoSAT suggested); record if the book resolves it.

The rule deliberately permits keeping the \emph{same} snapshot as the
parent. The allowance is required to preserve liveness of
$\Pi_{\mathrm{bft}}$ relative to $\Pi_{\mathrm{origbft}}$: honest
proposers may need to propose at a minimum rate --- Streamlet, for
example, needs three notarized blocks in consecutive epochs and assumes
proposals from honest leaders each epoch. The alternative of null
proposals was rejected because it would require specifying null-proposal
rules in $\Pi_{\mathrm{origbft}}$ (\texttt{construction.md}:611--619).
The machine-checked model of \S\ref{sec:xl-quint} diverges from the book
exactly here.

\begin{construction}[Honest proposal]\label{def:xl-honest-proposal}
An honest proposer chooses $P.\mathsf{headers\_bc}$ as the $\sigma$-block
tail of its bc-best-chain if that is consistent with the Linearity rule,
and otherwise sets $P.\mathsf{headers\_bc}$ to the same
$\mathsf{headers\_bc}$ as $P$'s parent bft-block. It makes no proposals
until its bc-best-chain is at least $\sigma + 1$ blocks long, since
otherwise $\mathsf{headers\_bc}$ cannot be constructed; at exactly $\sigma
+ 1$ blocks the snapshot is $\mathsf{Origin}_{\mathrm{bc}}$, which
satisfies Linearity because $\mathsf{Origin}_{\mathrm{bft}}$'s snapshot is
$\mathsf{Origin}_{\mathrm{bc}}$ (\texttt{construction.md}:625--637).
\end{construction}

\begin{construction}[Honest voting]\label{def:xl-honest-voting}
A validator first updates its view of both subprotocols with the
bc-headers in $P.\mathsf{headers\_bc}$, downloading and validating the
referenced bc-blocks, and recursively any unseen bft-chains referenced by
their $\mathsf{context\_bft}$ fields. The validation order ---
supermajority checks first, then PoW header checks, before block download
--- bounds denial-of-service, and the Linearity rule ensures only one
bc-chain need be validated per bft-chain. The validator votes for $P$
\emph{only if}:
\begin{description}
\item[Valid proposal criterion.] Proposal $P$ is bft-proposal-valid.
\item[Confirmed best-chain criterion.] The chain $\mathsf{snapshot}(P)$ is
  part of the validator's bc-best-chain at a bc-confirmation-depth of at
  least $\sigma$.
\end{description}
(\texttt{construction.md}:639--657.)
\end{construction}

The Confirmed best-chain criterion is subtle. It ensures that
$\mathsf{snapshot}(P)$ is bc-block-valid with $\sigma$ bc-block-valid
blocks after it in the validator's best chain, but
$P.\mathsf{headers\_bc}[\sigma - 1]$ need \emph{not} be in the
validator's best chain --- the chain may fork after
$\mathsf{snapshot}(P)$. Proposal validity must remain objective. Versus
Snap-and-Chat: that construction had the confirmed-snapshot voting
condition, but neither shipped headers with proposals nor made objective
confirmation a validity matter; shipping headers improves the
finalization-latency/security trade-off
(\texttt{construction.md}:659--678).

\paragraph{Classification.} \emph{Designed but unspecified}. The
prototype does not yet enforce this subsection's rules on its voting path
(\S\ref{sec:xl-divergences}).

\subsection{The best-chain side: $\Pi_{\mathrm{bc}}$}\label{sec:xl-bc}

All Crosslink~2 consensus-rule changes to $\Pi_{\mathrm{bc}}$ are
non-contextual; modulo them, contextual validity is the same as in
$\Pi_{\mathrm{origbc}}$. The bc-transaction consensus rules are entirely
unchanged: no reordering, no sanitization; parent bc-chains and heights
are preserved (\texttt{construction.md}:267, 701--703).

\begin{definition}[$\Pi_{\mathrm{bc}}$ validity]\label{def:xl-bc-validity}
For the genesis bc-block we must have
$\mathsf{Origin}_{\mathrm{bc}}.\mathsf{context\_bft} =
\mathsf{Origin}_{\mathrm{bft}}$, and therefore
$\mathsf{LF}(\mathsf{Origin}_{\mathrm{bc}}) =
\mathsf{Origin}_{\mathrm{bft}}$ (\texttt{construction.md}:684). A
bc-block $H$ is bc-block-valid iff \emph{all} of:
\begin{description}
\item[Inherited origbc rules.] Block $H$ satisfies the corresponding
  origbc-block-validity rules.
\item[Valid context rule.] $H.\mathsf{context\_bft}$ is bft-block-valid.
\item[Extension rule.] $\mathsf{LF}(H \lceil^{1}_{\mathrm{bc}})
  \preceq_{\mathrm{bft}} \mathsf{LF}(H)$.
\item[Last Final Snapshot rule.] $\mathsf{snapshot}(\mathsf{LF}(H))
  \preceq_{\mathrm{bc}} H$.
\item[Finality depth rule.] Define
  $\operatorname{finality\text{-}depth}(H) \coloneqq \mathrm{height}(H) -
  \mathrm{height}(\mathsf{snapshot}(\mathsf{LF}(H)))$; then either
  $\operatorname{finality\text{-}depth}(H) \le L$ or
  $\operatorname{is\text{-}stalled\text{-}block}(H)$.
\end{description}
(\texttt{construction.md}:686--693; the community protocol guide restates
the Last Final Snapshot rule verbatim, crosslink-protocol-guide @
\texttt{b77d845}, \texttt{src/tfl-lexicon.md}:3--5.)
\end{definition}

\paragraph{Stalled Mode.} The finality gap at time $t$ is, roughly, the
number of bc-blocks not yet finalized; the Finality depth rule makes this
precise as the objective $\operatorname{finality\text{-}depth}$ of a
block. If the gap exceeds the threshold $L$, that likely signals an
exceptional or emergency condition in which accepting spends into new
bc-blocks is undesirable. The Stalled Mode policy is modelled by the
predicate $\operatorname{is\text{-}stalled\text{-}block} :
\text{bc-block} \to \mathrm{boolean}$; a block satisfying it is a
\emph{stalled block} (\texttt{construction.md}:401--413). The node-local
prose notion (``blocks not yet finalized at time $t$'') and the objective
per-block notion differ in corner cases just after a reorg
(\texttt{construction.md}:695--699).
%% TODO: pick one canonical finality-gap definition for this guide and
%% flag the post-reorg corner case where the two notions differ.

Two cautions. First, a bc-block producer is only constrained to produce
stalled blocks while its view of the finalization point is not advancing;
an adversary that subverts the BFT protocol \emph{without} stalling
finalization is never constrained by Stalled Mode
(\texttt{construction.md}:401--413). Second, the book warns that an
honest bc-block-producer must \emph{not} use information from the BFT
protocol beyond the specified consensus rules to decide which
bc-valid-chain to follow; additional constraints could let an adversary
that subverts $\Pi_{\mathrm{bft}}$ influence the bc-best-chain in ways
outside the safety argument (\texttt{construction.md}:715--717).

\begin{construction}[Honest block production]\label{def:xl-honest-production}
The producer must follow the $\Pi_{\mathrm{bc}}$ validity rules,
including producing a stalled block when the Finality depth rule requires
it. To choose $H.\mathsf{context\_bft}$ it considers the set
\[
\{\, T : T \text{ is bft-block-valid and }
  \mathsf{LF}(H \lceil^{1}_{\mathrm{bc}}) \preceq_{\mathrm{bft}}
  \operatorname{bft\text{-}last\text{-}final}(T) \,\},
\]
chooses one of the longest such chains $C$ --- breaking ties first by
maximizing the score of the snapshot of
$\operatorname{bft\text{-}last\text{-}final}(C)$, then by smallest hash
--- and sets $H.\mathsf{context\_bft} = C$.
Rationale: choosing a longest chain follows Streamlet's
build-on-longest-notarized; the score tie-break inhibits an adversary
from finalizing a lesser-scoring chain
(\texttt{construction.md}:705--731).
\end{construction}

\paragraph{Classification.} \emph{Designed but unspecified}; the
deployment track has, at the pinned commit, dropped the Stalled Mode
rules and does not enforce the Extension, Last Final Snapshot, or
Finality depth rules in bc-block validity (\S\ref{sec:xl-divergences}).

\subsection{The two ledgers}\label{sec:xl-ledgers}

Each node $i$ derives two chains from its views: a locally finalized
chain $\mathsf{fin}_i^t$ and a locally bounded-available chain
$(\mathsf{ba}_{\mu})_i^t$, for a per-node confirmation depth $\mu$. The
book's macros render the task names \textsf{localfin} and
\textsf{localba}$_{\mu}$ as $\mathsf{fin}$ and $\mathsf{ba}$, and
\textsf{snapshotlf}$(H)$ as
$\mathsf{snapshot}(\mathsf{LF}(H))$ (tfl-book @ \texttt{fe6e1d6},
\texttt{macros.txt}:29, 63--64); we use the rendered forms.

\paragraph{The finalized chain.} The locally finalized chain starts at
$\mathsf{fin}_i^0 = \mathsf{Origin}_{\mathrm{bc}}$ and must not be
exposed to clients until the node has synced
(\texttt{construction.md}:462--488).

\begin{construction}[Finalization update
($\mathsf{updatefin}$)]\label{def:xl-updatefin}
On a best-chain update at time $t$, with previous state
$\mathsf{fin}_i^s$, node $i$ computes $N \coloneqq
\mathsf{candidate}(\mathsf{ch}_i^t)$ and:
\begin{enumerate}[label=(\roman*)]
\item if $N \succeq_{\mathrm{bc}} \mathsf{fin}_i^s$, sets
  $\mathsf{fin}_i^t \leftarrow N$;
\item otherwise sets $\mathsf{fin}_i^t \leftarrow \mathsf{fin}_i^s$, and
  if additionally $N \not\preceq_{\mathrm{bc}} \mathsf{fin}_i^s$ --- that
  is, $N$ conflicts with $\mathsf{fin}_i^s$ --- records a
  \emph{finalization safety hazard}, possible only if the global safety
  assumptions are violated. The hazard record includes $\mathsf{ch}_i^t$
  and the $\mathsf{fin}$ update history since the last update that was an
  ancestor of $N$.
\end{enumerate}
(\texttt{construction.md}:475--486.)
\end{construction}

\begin{definition}[Local finalization linearity]\label{def:xl-lfl}
Node $i$ has \emph{Local finalization linearity} up to time $t$ iff the
time series of bc-blocks $\mathsf{fin}_i^{r \le t}$ is bc-linear
(\texttt{construction.md}:466--469).
\end{definition}

The update rule guarantees Local finalization linearity \emph{by
construction} (\texttt{construction.md}:462--488). Local state is needed
because on a reorg --- even one shorter than $\sigma$ --- the value
$\mathsf{candidate}(\mathsf{ch}_i^t)$ may temporarily roll back; if Final
Agreement holds, the new snapshot is an ancestor of the old and will roll
forward along the same path, but applications must not see the temporary
rollback (\texttt{construction.md}:497--499). This is precisely the
service today's applications lack: wallets repair rollbacks after the
fact by truncate-and-rescan (\emph{Wallet Guide}, \S``Reorganisations:
truncate and rescan''), and anchor validity tracks the unfinalized chain
(\emph{Ironwood Guide}, \S``Proof of ownership of \emph{some} valid note:
the commitment tree'').

\begin{lemma}[Local fin-depth]\label{lem:xl-fin-depth}
In any execution of Crosslink~2, for any node $i$ honest at time $t$,
there exists a time $r \le t$ such that $\mathsf{fin}_i^t
\preceq_{\mathrm{bc}} \mathsf{ch}_i^r \lceil^{\sigma}_{\mathrm{bc}}$.
\end{lemma}
\begin{proof}
By $\mathsf{candidate}(\mathsf{ch}_i^u) \preceq_{\mathrm{bc}}
\mathsf{ch}_i^u \lceil^{\sigma}_{\mathrm{bc}}$ applied at the last time
$\mathsf{fin}$ changed, or at genesis
(\texttt{construction.md}:490--495).
\end{proof}

\begin{definition}[Assured Finality]\label{def:xl-assured-finality}
An execution of Crosslink~2 has \emph{Assured Finality} iff for all times
$t$, $u$ and all nodes $i$, $j$ (potentially the same) such that $i$ is
honest at time $t$ and $j$ is honest at time $u$, we have $\mathsf{fin}_i^t
\sim_{\mathrm{bc}} \mathsf{fin}_j^u$ (\texttt{construction.md}:504--506).
\end{definition}

Assured Finality is the main safety goal of Crosslink~2, intended to hold
under reasonable assumptions about \emph{either} $\Pi_{\mathrm{origbft}}$
\emph{or} $\Pi_{\mathrm{origbc}}$. Its restriction to $i = j$ is
equivalent to Local finalization linearity, so Assured Finality implies
Local finalization linearity for all honest nodes; the converse
implication does not hold --- the local property is necessary but not
sufficient (\texttt{construction.md}:473, 501--508).

\paragraph{Why $\mathsf{candidate}(H)$, not
$\mathsf{snapshot}(\mathsf{LF}(H))$.} The Last Final Snapshot rule
guarantees $\mathsf{snapshot}(\mathsf{LF}(H)) \preceq_{\mathrm{bc}} H$,
but \emph{not} that its depth relative to $H$ is at least $\sigma$. Using
$\mathsf{candidate}(H)$ ensures the finalization point is at least
$\sigma$-confirmed, which the proof of Assured Finality needs:
$\mathsf{candidate}(H) \preceq_{\mathrm{bc}} H
\lceil^{\sigma}_{\mathrm{bc}}$ combines with the Local fin-depth lemma
and Prefix Agreement. An open book TODO offers the alternative of
strengthening the Last Final Snapshot rule to
$\mathsf{snapshot}(\mathsf{LF}(H)) \preceq_{\mathrm{bc}} H
\lceil^{\sigma}_{\mathrm{bc}}$ (\texttt{construction.md}:510--518).
%% TODO: designed-but-unspecified fork in the design (construction.md:517):
%% candidate(H) as specified versus a strengthened Last Final Snapshot
%% rule; present candidate(H) as normative-for-the-book, track the
%% alternative.

\begin{definition}[Locally bounded-available chain]\label{def:xl-ba}
For a per-node confirmation depth $\mu$:
\[
(\mathsf{ba}_{\mu})_i^t =
\begin{cases}
\mathsf{ch}_i^t \lceil^{\mu}_{\mathrm{bc}} &
  \text{if } \mathsf{fin}_i^t \preceq_{\mathrm{bc}} \mathsf{ch}_i^t
  \lceil^{\mu}_{\mathrm{bc}},\\
\mathsf{fin}_i^t & \text{otherwise}
\end{cases}
\]
(\texttt{construction.md}:522--527).
\end{definition}

The ``otherwise'' arm is needed because after $\mathsf{fin}$ switches
forks under Zcash's damped difficulty adjustment
(\texttt{PoWAveragingWindow} $= 17$ blocks, \texttt{PoWMedianBlockSpan}
$= 11$ blocks, quoted by the book from the Zcash protocol specification
\S\,diffadjustment), the truncated best chain $\mathsf{ch}_i^t
\lceil^{\mu}_{\mathrm{bc}}$ might not descend from $\mathsf{fin}_i^t$
even with all safety assumptions satisfied; the definition also makes the
Ledger prefix property trivial to prove. The security goal for
$\mathsf{ba}_{\mu}$ is Agreement on the view $\mathsf{ba}_{\mu}$
(Definition~\ref{def:xl-agreement-view};
\texttt{construction.md}:545--555, 552).

\begin{theorem}[Ledger prefix property]\label{thm:xl-ledger-prefix}
For any node $i$ that is honest at time $t$, and any confirmation depth
$\mu$, $\mathsf{fin}_i^t \preceq_{\mathrm{bc}} (\mathsf{ba}_{\mu})_i^t$.
\end{theorem}
\begin{proof}
By construction of $(\mathsf{ba}_{\mu})_i^t$
(\texttt{construction.md}:531--536).
\end{proof}

\begin{lemma}[Local ba-depth]\label{lem:xl-ba-depth}
In any execution of Crosslink~2, for any confirmation depth $\mu \le
\sigma$ and any node $i$ honest at time $t$, there exists a time $r \le
t$ such that $(\mathsf{ba}_{\mu})_i^t \preceq_{\mathrm{bc}}
\mathsf{ch}_i^r \lceil^{\mu}_{\mathrm{bc}}$.
\end{lemma}
\begin{proof}
Either $(\mathsf{ba}_{\mu})_i^t = \mathsf{fin}_i^t$, and the Local
fin-depth lemma applies with $\mu \le \sigma$; or $(\mathsf{ba}_{\mu})_i^t
= \mathsf{ch}_i^t \lceil^{\mu}_{\mathrm{bc}}$, trivially with $r = t$
(\texttt{construction.md}:538--543).
\end{proof}

\paragraph{Stall, do not finalize unsafely.} In normal operation
$\mathsf{fin}_i^t$ lags only a few blocks behind
$(\mathsf{ba}_{\sigma})_i^t$, unless the chain enters Stalled Mode. With
$\mu = \sigma$, the choice between following $\mathsf{fin}$ and following
$\mathsf{ba}$ is a choice of stall behaviour: stall immediately
($\mathsf{fin}$), or keep processing user transactions until $L$ blocks
have passed ($\mathsf{ba}$) (\texttt{construction.md}:415).

\paragraph{Syncing and checkpoints.} Implementations should bake a
checkpointed bft-block $B_{\mathrm{checkpoint}}$ into each release; node
$i$ exposes $\mathsf{fin}_i^t$ and $(\mathsf{ba}_{\mu})_i^t$ only once
synced, that is, once $B_{\mathrm{checkpoint}} \preceq_{\mathrm{bft}}
\mathsf{LF}(\mathsf{ch}_i^t)$,
$\mathsf{snapshot}(B_{\mathrm{checkpoint}}) \preceq_{\mathrm{bc}}
\mathsf{fin}_i^t$, and the timestamp of $\mathsf{fin}_i^t$ is within a
threshold of the current time (\texttt{construction.md}:557--562).

The chapter closes: ``At this point we have completed the definition of
Crosslink 2. In Security Analysis of Crosslink 2, we will prove it
secure.'' The structural rules above are the complete definition; all
proofs beyond the in-chapter lemmas live in
\texttt{security-analysis.md} (\texttt{construction.md}:735); this
volume's account of that analysis is \S\ref{sec:xl-security}.
%% TODO: book-internal TODO at construction.md:86 (possibly-unused
%% ledger-prefix notation).

\paragraph{Classification.} \emph{Designed but unspecified}, with the
candidate-versus-strengthened-rule fork explicitly open in the book.

\subsection{Parameters}\label{sec:xl-parameters}

The construction has three parameters and fixes production values for
none of them.
\begin{itemize}
\item The bc-confirmation-depth $\sigma \in \N^{+}$
  (\texttt{construction.md}:391). No production value is fixed anywhere
  in the sources.
\item The finalization gap bound $L \in \N^{+}$, significantly greater
  than $\sigma$; ``In practice, $L$ should be at least $2\sigma$''
  (\texttt{construction.md}:391, 405).
\item The per-node depth $\mu$ with $0 < \mu \le \sigma$; the default and
  recommendation is $\mu = \sigma$. A smaller $\mu$ is at the node's own
  risk and does not affect $\mathsf{fin}$
  (\texttt{construction.md}:395--399).
\end{itemize}
The prototype ships \texttt{PROTOTYPE\_PARAMETERS} with $\sigma = 3$ and
$L = 7$ --- satisfying $L \ge 2\sigma$ --- under the explicit caveat
``No verification has been done on the security or performance of these
parameters'' (zebra-crosslink @ \texttt{6d02a1b},
\texttt{zebra-crosslink/src/chain.rs}:216--222).
%% TODO: no production sigma or L exists in any source; revisit when a
%% ZIP or the productionized implementation fixes them.

For calibration: today's wallet policy defaults are $3$ confirmations for
trusted and $10$ for untrusted outputs under the ZIP-315
\texttt{ConfirmationsPolicy} (\emph{Wallet Guide}, \S``Confirmations,
trust, and spendability'') --- a per-wallet risk heuristic, where
$\sigma$ is a consensus-visible input to finalization. On the node side,
the $k$-deep interface constant of the deployed best-chain protocol is
\texttt{MAX\_BLOCK\_REORG\_HEIGHT} $= 1000$ in mainline Zebra, documented
as ``a local-only node policy; it is not part of consensus'' and sized as
defence-in-depth against sustained consensus splits (zebra @
\texttt{9f3166b}, \texttt{zebra-chain/src/parameters/constants.rs}:30);
the zebra-crosslink fork still uses the older value $99 =
\texttt{MIN\_TRANSPARENT\_COINBASE\_MATURITY} - 1$ (zebra-crosslink @
\texttt{6d02a1b}, \texttt{zebra-state/src/constants.rs}:31,
\texttt{zebra-chain/src/transparent.rs}:52). How that constant will
relate to $\sigma$ is undetermined.

\paragraph{Classification.} The parameter constraints are \emph{designed
but unspecified}; the concrete values are an \emph{open problem}. The
quoted constants are \emph{specified} in code at the pinned commits.

\subsection{The machine-checked fragment}\label{sec:xl-quint}

The Informal Systems Quint specification (crosslink-spec @
\texttt{8edc3e9}) is an explicit proof of concept covering \emph{only}
the \texttt{isValid} check for BFT proposals --- Tail Confirmation plus
Linearity --- over abstract single-node PoW and BFT block stores. Its
README states the logic was inferred mostly from a YouTube talk and
``might not be up-to-date with the current protocol design ideas''. The
Quint file \texttt{src/validity.qnt} is tangled from
\texttt{validity.md} via the \texttt{lmt} tool (README.md:25--32;
validity.md:1--7).

The checked predicate is:
\begin{quote}\ttfamily\small
pure def isValid(bft\_store, pow\_store, p) = all \{\\
\hspace*{1.5em}pow\_store.contains(p.block),\\
\hspace*{1.5em}p.bftContext == bft\_store.last(),\\
\hspace*{1.5em}sigmaConfirmed(p.block, p.confirmation\_tail,
confirmation\_depth),\\
\hspace*{1.5em}linear(p.block, p.bftContext, pow\_store)\\
\}
\end{quote}
where \texttt{sigmaConfirmed(b, t, sigma)} requires a fork-free
confirmation tail of size \emph{at least} $\sigma$ whose parents chain
back to \texttt{b}, and \texttt{linear(b, context, work\_store)} requires
\texttt{work\_store.isDescendant(context.pow\_hash, b.hash)} \emph{and}
\texttt{not(context.pow\_hash == b.hash)} (crosslink-spec @
\texttt{8edc3e9}, \texttt{validity.md}:131--138, 150--158, 171--176).

\paragraph{Machine-checked results.} The invariant \texttt{prefix\_inv}
--- ``the PoW history defined by BFT finalization is a path in the PoW
block store'', formalized for all adjacent indices $i$ as
\begin{quote}\ttfamily\small
pow\_store.isDescendant(bft\_store[i].pow\_hash,
bft\_store[i+1].pow\_hash)
\end{quote}
--- found no
violation over $10000$ random traces of $20$ steps, and was model checked
by TLC (via \texttt{quint compile -{}-target tlaplus}) in about three
minutes on the verification instance (\texttt{confirmation\_depth} $= 1$,
at most $8$ blocks per store; the simulation instance uses
\texttt{confirmation\_depth} $= 2$, unbounded). The reachability witness
\texttt{finalization\_witness} ($3$ BFT blocks stored) was found by both
the simulator and TLC, and the run \texttt{finalizationTest} passes under
\texttt{quint test} (README.md:39--89; validity.md:376--461, 405--410,
463--482).

\paragraph{Divergences from the book.} The model checks a rule set that
neighbours, but does not match, the book's:
\begin{itemize}
\item Its \texttt{linear} predicate \emph{forbids} keeping the same
  snapshot as the parent, via the conjunct
  \texttt{not(context.pow\_hash == b.hash)},
  whereas the book's Linearity rule allows equality and argues the
  allowance is necessary for $\Pi_{\mathrm{bft}}$ liveness
  (\S\ref{sec:xl-bft}). The spec itself carries the comment ``this
  actually seems allowed. TODO: perhaps remove''
  (\texttt{validity.md}:171--176 versus
  \texttt{construction.md}:573, 611--619).
\item Its \texttt{sigmaConfirmed} accepts a tail of size $\ge \sigma$
  where the book requires exactly $\sigma$ headers; its proposal carries
  the snapshot block itself plus an \emph{unordered} set of tail blocks,
  against the book's ordered deepest-first header sequence with
  $\mathsf{snapshot}(B) = B.\mathsf{headers\_bc}[0]
  \lceil^{1}_{\mathrm{bc}}$; and its \texttt{isValid} requires
  \texttt{p.bftContext == bft\_store.last()}, modelling the BFT side as a
  single forkless linear chain, unlike the book's set of bft-chains
  (\texttt{validity.md}:65--96, 131--158 versus
  \texttt{construction.md}:417--436).
\item Not modelled at all: $\mathsf{context\_bft}$ in bc-blocks, the
  Extension, Last Final Snapshot, and Finality depth rules, Stalled Mode,
  $\operatorname{bft\text{-}last\text{-}final}$, voting and notarization
  proofs, and signatures (crosslink-spec @ \texttt{8edc3e9} README.md,
  validity.md).
\end{itemize}

\paragraph{Classification.} The results above are \emph{specified} in
the strongest pre-spec sense available --- machine-checked at a pinned
commit --- but of a \emph{neighbouring} rule set: stricter than the book
on snapshot equality, looser on tail size and BFT forks, and silent on
the whole $\Pi_{\mathrm{bc}}$ side. We therefore cite the Quint model as
a proof-of-concept formalization, never as verification of the book's
construction.

\subsection{Divergences in the deployment track}\label{sec:xl-divergences}

The Shielded Labs/Eiger implementation (zebra-crosslink @
\texttt{6d02a1b}) carries a byte-for-byte copy of the book's construction
chapter, prepended with three warning admonitions: the copy is ``an
almost direct paste'' committing to a precise text for a security design
review, not reviewed or endorsed by Daira-Emma Hopwood or Jack Grigg; the
zebra-crosslink design diverges by \emph{not} including the Stalled Mode
rules, on the rationale that stakers redelegating prevents hazards such
as BFT stalls; and it does not use the outer signature, trading a hazard
many participants could leverage for one only the current proposer could,
judged not worth it --- with the recorded consequence that a direct hash
or copy of the most recent BFT finality certificate is insufficient
evidence that a specific bitwise copy is canonical
(\texttt{book/src/design/cl2-construction.md}:3--17, diffed against
\texttt{construction.md} at \texttt{fe6e1d6}).
%% RESOLVED at 6d02a1b (2026-07-15): the only trace of contingency mode
%% is the GH #172 RPC deliverable (deliverables.md:12); no rule, no
%% code, and finalization_gap_bound is never enforced. The L-bound is
%% dropped at the design level, with a contingency-mode RPC surface
%% still promised; the M3-era coinbase-only mode (forum, 2025-08-27) is
%% superseded by the cl2 admonition. Re-check only against
%% post-productionization commits.

At the pinned commit, the further construction-level divergences are:
\begin{itemize}
\item \emph{Fat pointer for $\mathsf{context\_bft}$.} The fork adds
  \texttt{fat\_pointer\_to\_bft\_block} to the PoW block header: a bundle
  of Ed25519 signed precommit votes --- a $44$-byte vote template whose
  first $32$ bytes are the BLAKE3 hash of the target BFT block, plus
  $96$-byte entries of public key and vote signature --- constructed from
  the deciding round's signed precommit votes (tenderlink's
  \texttt{FatPointerToBftBlock3}, converted at
  \texttt{zebra-crosslink/src/lib.rs}:1846--1860); the earlier Malachite
  commit-certificate conversion survives only in the undeclared
  dead-code module \texttt{src/malctx.rs}. It replaces both the book's
  plain $\mathsf{context\_bft}$ commitment and the
  notarization-proof/outer-signature structure with a self-certifying
  pointer (\texttt{zebra-chain/src/block/header.rs}:109--133, 195--201;
  \texttt{zebra-crosslink/src/lib.rs}:1811--2052).
\item \emph{Immediate finality instead of
  $\operatorname{bft\text{-}last\text{-}final}$.} The BFT engine is
  Tendermint-style: the in-house \texttt{tenderlink} package
  (\texttt{git.sr.ht/\textasciitilde azmr/stuff} rev \texttt{0e88509}),
  driven directly from \texttt{zebra-crosslink/src/lib.rs}. Each
  tenderlink-decided BFT block is treated as immediately final: the
  handler (installed as \texttt{tenderlink::ClosureToPushDecidedBlock},
  \texttt{lib.rs}:1234) sets the latest final block from the
  decided block's headers and issues
  \texttt{zebra\_state::Request::CrosslinkFinalizeBlock} so the PoW state
  finalizes that ancestor --- the implementation's $\mathsf{fin}$
  extraction. With no BFT-final block known, the node falls back to the
  PoW block at \texttt{MAX\_BLOCK\_REORG\_HEIGHT} depth
  (\texttt{zebra-crosslink/Cargo.toml}:74--79;
  \texttt{zebra-crosslink/src/lib.rs}:235--330, 499--601).
\item \emph{Validity rules mostly unimplemented.} The Rust
  \texttt{BftBlock} carries the fields \texttt{version},
  \texttt{height}, \texttt{previous\_block\_fat\_ptr},
  \texttt{finalization\_candidate\_height}, and \texttt{headers}; its
  sole constructor \texttt{try\_from} validates only
  $\texttt{headers.len()} = \sigma$ and logs ``not yet implemented: all
  the documented validations''; deserialization rejects more than $2048$
  headers. The vote-path validation checks only that the new block's
  previous fat pointer hashes to the current tip's, and that the node
  knows the PoW block referenced by \texttt{headers.first()}. The
  Linearity rule, Tail Confirmation validation, and the Confirmed
  best-chain criterion are \emph{not} enforced on this path; the
  decided-block handler additionally verifies fat-pointer signatures,
  with a TODO to check the public keys against the roster
  (\texttt{zebra-crosslink/src/chain.rs}:61--151, 99--104;
  \texttt{zebra-crosslink/src/lib.rs}:522--541, 790--822).
\item \emph{Header-order inconsistency.} The in-memory header order is
  documented as ``reversed from the specification'' (built tip-first),
  and the accessors disagree at this commit:
  \texttt{finalization\_candidate()} returns \texttt{headers.last()}
  while the decided-block handler computes the newly final PoW hash from
  \texttt{headers.first()}; the consistency assertion against
  \texttt{finalization\_candidate\_height} is commented out, and header
  fetching carries the comment ``TODO: do we want these in chain order or
  walk-back order'' (\texttt{chain.rs}:52--54, 124--126;
  \texttt{lib.rs}:543--546, 809, near :432).
  %% TODO: pin down which end of the header vector is the snapshot
  %% before this guide describes the deployed fin extraction.
\end{itemize}

Beyond the book's construction, the deployment track adds a
mechanism-design layer --- finalization rewards paid only in blocks that
advance finality, a staking transaction format with the Crosslink Staking
Orchard Restriction, staking epochs, a roster in ledger state, and the
General Ledger State Design Invariant that all ledger state, including
the PoS/BFT roster and finality status, is computable entirely from PoW
blocks and transactions, with finality status the sole
``height-eventual'' state (\texttt{book/src/design/nutshell.md}:23--128).
%% TODO: the mechanism-design layer and the five named security
%% properties (book/src/design/security-properties.md) belong to later
%% sections of this guide; cross-reference once drafted.

\paragraph{Classification.} The code facts above are \emph{specified} in
the pinned-commit sense, but the track is moving: Milestone~5
(``Pre-Productionization'') completed 2026-04-15, productionization is
ongoing with no confirmed mainnet activation date, and the first
independent security and mechanism-design review is complete
(forum.zcashcommunity.com/t/shielded-labs-crosslink-deployment-updates/49706;
shieldedlabs.net/roadmap; x.com/ShieldedLabs status
\texttt{2019547550009946528}; all accessed 2026-07-15). Every
implementation claim in this subsection is pinned to commit
\texttt{6d02a1b} and must be re-verified against a post-productionization
commit before it is treated as the deployed construction.
%% NOTE: 6d02a1b already post-dates M5 and pins Tenderlink; re-verify
%% only against future productionization-phase commits.

%% Crosslink Guide, section: Stake, delegation, and slashing. Body only;
%% assumes the shared preamble of the series (ironwood-guide.tex lines 1-90).
%%
%% Source pins: tfl-book @ fe6e1d6f (2024-12-13); zebra-crosslink @ 6d02a1b8
%% (2026-04-21; shallow clone, depth 1 -- per-file history unavailable);
%% crosslink-spec @ 8edc3e94 (2025-07-23); crosslink-protocol-guide @
%% b77d8456 (2026-03-26); zips @ 69610984 (2026-07-09). Web sources fetched
%% 2026-07-15. The local quint-crosslink clone is empty (branch main has no
%% commits); see the TODO at the end of this section.

\section{Stake, delegation, and slashing}\label{sec:xl-stake}

Crosslink's finality vote is stake-weighted, so the construction owes
answers to the questions every proof-of-stake system faces: how ZEC becomes
voting weight, how a wallet delegates without surrendering custody, and
what punishes a finalizer that misbehaves. This section collects those
answers --- and their absences. The division of labour among the sources is
unusual and governs every citation below: the ECC design book abstains from
the whole topic, listing ``PoS subprotocol selection'' and ``Issuance and
supply mechanics, such as how much ZEC stakers may earn'' among its
``Major Missing Components'' (tfl-book,
\texttt{src/introduction/status-and-next-steps.md} at \texttt{fe6e1d6f}),
so every concrete stake, delegation, and slashing decision lives with
Shielded Labs: the zebra-crosslink design book (below, ``the implementation
book'') and two dated posts\footnote{Shielded Labs, ``Crosslink: Updated
Staking Design'',
\url{https://shieldedlabs.net/crosslink-updated-staking-design/} (forum
post 2025-10-27, blog dated 2025-11-05); ``Crosslink Early Tokenomics
Design'', forum thread t/52154 (2025-09-12); deployment updates, forum
thread t/49706; forum threads under
\url{https://forum.zcashcommunity.com/t/}. All fetched 2026-07-15.}.

Nothing in this section reaches the \emph{specified} bin. No Crosslink or
staking ZIP exists: at \texttt{zcash/zips} @ \texttt{69610984}
(2026-07-09) the only Crosslink references are a citation footnote in
ZIP~218 and a tfl-book reference in
\texttt{draft-ecc-onchain-accountable-voting}; the ``IO Finalizer'' and
``Spend Finalizer'' of ZIP~374 are unrelated PCZT roles. The implementation
book states ``We will be refining a ZIP Draft [Google Doc link] which will
enter the formal ZIP process as the prototype approaches maturity''
(\texttt{book/src/design.md} at \texttt{6d02a1b8}). The machine-checked
corpus is equally silent: the Informal Systems Quint model covers only the
proposal-validity function (Tail Confirmation and Linearity rules) and
models no stake, roster, voting power, delegation, bonding, or slashing
(crosslink-spec @ \texttt{8edc3e94}, \texttt{validity.md},
\texttt{src/validity.qnt}).
%% TODO: replace the prose bin references with \S\ref to this volume's
%% scope-and-method section once drafted; the three-bin discipline
%% originates in the \emph{Tachyon Guide}, \S``The three-bin method''.
Everything below is therefore \emph{designed but unspecified} at best, or
implementation-defined prototype behaviour, or an \emph{open problem}; we
keep the bin visible throughout. The implementation book itself warns that
its extensions ``may violate some of the assumptions in the security
proofs'' of the Crosslink~2 construction and calls its staking chapter ``a
provisional sketch'' that ``will eventually be entirely redundant with...
a canonical ZIP'' (\texttt{book/src/design.md},
\texttt{book/src/design/nutshell.md} at \texttt{6d02a1b8}).

\subsection{Stake in the abstract construction}\label{sec:xl-stake-abstract}

The Crosslink~2 construction treats stake as an abstraction: ``For each
epoch, there is a fixed number of voting units distributed between the
star-bft-nodes, which they use to vote for a star-bft-proposal''
(tfl-book, \texttt{src/design/crosslink/construction.md} at
\texttt{fe6e1d6f}; the implementation book mirrors the passage verbatim in
\texttt{book/src/design/cl2-construction.md}). The notarization rule for a
proposal $P$ ``must require at least a two-thirds absolute supermajority
of voting units in $P$'s epoch to have been cast for $P$'', and may
require more (ibid.). How ZEC-denominated stake maps onto voting units is
delegated to the unselected PoS subprotocol --- precisely the ``Major
Missing Component'' above.

The adversary bound is stated over the same abstraction. Notation follows
the book: $\Pi_{\mathsf{bft}}$ is the BFT subprotocol as adapted by
Crosslink, $\Pi_{\mathsf{origbft}}$ its underlying engine; the book's
$\star$bft is a wildcard ranging over both.
%% TODO: harmonize $\Pi$ notation with this volume's construction section.

\begin{definition}[One-third bound on non-honest voting]
\label{def:xl-nonhonest}
An execution of $\Pi_{\mathsf{bft}}$ has the \emph{one-third bound on
non-honest voting} property iff for every epoch, \emph{strictly} fewer
than one third of the total voting units for that epoch are ever cast
non-honestly. A voting unit is cast non-honestly iff: it is cast other
than by the holder of the unit (due to key compromise or any flaw in the
voting protocol, for example); or it is double-cast; or the holder,
following the conditions for honest voting in $\Pi_{\star\mathsf{bft}}$
according to its view, should not have cast that vote (tfl-book,
\texttt{src/design/crosslink/construction.md} at \texttt{fe6e1d6f},
near-verbatim).
\end{definition}

The quantifier ``ever'' is load-bearing. The book's own security caveat
spells out the consequence: the bound counts ballots cast at \emph{any}
time in the execution, including votes cast with a compromised key for a
long-finished epoch; ``Therefore, validator keys of honest nodes must
remain secret indefinitely. Whenever a key is rotated, the old key must be
securely deleted'', with further discussion deferred to tfl-book issue
\#140 (ibid.). We return to what this demands of finalizer key management
in \S\ref{sec:xl-stake-keys}.

The one commitment the ECC book does make about delegation is a goal, not
a mechanism: ``The protocol should enable users of shielded mobile wallets
to delegate ZEC to PoS consensus providers and earn a return on that ZEC
coming via ZEC issuance or fees. Doing this may expose users to a risk of
loss of delegated ZEC (such as through `slashing fees'). The protocol must
guarantee that PoS consensus providers have no discretionary control over
such delegated funds (including that they cannot steal those funds)''
(tfl-book, \texttt{src/design/goals.md} at \texttt{fe6e1d6f}). Note the
tension recorded for \S\ref{sec:xl-stake-slashing}: the goal contemplates
``slashing fees'' as a user risk; the deployed design forswears them.

\paragraph{Classification.} \emph{Designed but unspecified}: the voting-unit
model and Definition~\ref{def:xl-nonhonest} come from a design book with
security arguments, not a specification. The mapping from stake to voting
units --- proportional weighting as in the prototype
(\S\ref{sec:xl-stake-prototype}) versus discretized units, and how the
per-epoch roster snapshot is taken --- is an \emph{open problem}: no design
document fixes it.

\subsection{Positions, the roster, and the active set}
\label{sec:xl-stake-roster}

The canonical staking statement is one sentence: ``Staking in Crosslink
takes place exclusively within the Orchard pool'' (Shielded Labs, Updated
Staking Design, 2025-10-27/11-05). Participants lock ZEC to delegate stake
to a chosen finalizer and earn rewards; Penumbra-style fully shielded
delegations were considered and deferred to future versions, prioritizing
time-to-market (ibid.). The 2025-10/11 design predates NU6.3's pool
terminology: read ``the Orchard pool'' there as the Orchard-protocol
shielded pool --- since NU6.3, the Ironwood pool. That pool is the one this
series constructs in the \emph{Ironwood Guide} (\S``Representation of a
note: notes and note commitments''); nothing below alters its note
machinery.

\begin{definition}[Staking position, roster, stake weight, active set]
\label{def:xl-roster}
The \emph{Roster} is a Ledger State table tracking staking positions
created by stakers and attached to specific finalizers. Each
\emph{staking position} comprises at least: a target finalizer
verification key, a staked ZEC amount, and an unstaking/redelegation
verification key. The finalizer verification key ``serves as the sole
on-chain reference to a specific finalizer''; an entity may produce any
number of finalizer keys. Unstaking/redelegation keys ``are unique to the
position itself (and not linked on-chain to a user's other potential
staking positions or other activity)''. The sum of outstanding staking
positions designating a specific finalizer verification key at a given
block height is the \emph{stake weight} of that finalizer, and equals its
BFT voting weight. At a given height the top $K$ (``currently'' $K = 100$)
finalizers by stake weight are the \emph{active set} (implementation book,
\texttt{book/src/design/nutshell.md} at \texttt{6d02a1b8}).
\end{definition}

The roster obeys the implementation book's General Ledger State Design
Invariant: ``All changes to Ledger State can be computed entirely from
(PoW) blocks, transactions, and data transitively referred to by such.
This includes all existing Ledger State in mainnet Zcash today, the PoS
and BFT roster state (see below), and the finality status of blocks and
transactions'' (ibid.). Roster and active set are thus height-immediate
--- a function of the chain up to the height in question. Finality status
is the sole height-eventual Ledger State: a property of height $H$
computable only from a later attesting block; all nodes seeing the
attesting block agree, and if the attesting block rolls back, the protocol
guarantees an alternative block will eventually attest the same candidate
(ibid.). Wallets today approximate the missing finality signal with
ZIP-315 confirmation depths (\emph{Wallet Guide}, \S``Transaction
construction''); Crosslink turns it into consensus state.

\paragraph{Classification.} \emph{Designed but unspecified} throughout.
The active-set bound is explicitly provisional (``currently'' 100), the
book's own TODO list flags ``active set selection'' as unresolved
(\texttt{book/src/design/nutshell.md} at \texttt{6d02a1b8}), and the
prototype applies no top-$K$ truncation at all
(\S\ref{sec:xl-stake-prototype}).

\subsection{Staking actions and draft consensus rules}
\label{sec:xl-stake-actions}

The implementation book extends the transaction format with an optional
staking-action field carrying four abstract actions
(\texttt{book/src/design/nutshell.md} at \texttt{6d02a1b8}):

\begin{itemize}
\item \emph{stake} --- creates a staking position assigning a transparent
  ZEC amount, which must be a power of 10, to a finalizer, and includes a
  signature verification key authorizing later redelegate/unstake actions;
\item \emph{redelegate} --- identifies a position, a new finalizer, and a
  signature valid under the position's published verification key;
\item \emph{unbond} --- moves a position into the \emph{unbonding} state:
  redelegations become invalid and the amount stops contributing to
  staking weight for finalizers or the staker;
\item \emph{claim} --- removes an unbonding-state position and contributes
  the ZEC to the chain value pool balance, which --- by the Orchard
  restriction below --- may only go to an Orchard destination and/or fees.
\end{itemize}

Value flows must balance per the Zcash protocol specification's chain
value pool accounting (\S4.17); the Orchard side of that balance is the
value-commitment machinery of the \emph{Ironwood Guide} (\S``Conservation
of value: value commitments and the binding signature''). The draft
consensus rules, verbatim from the implementation book
(\texttt{book/src/design/nutshell.md} at \texttt{6d02a1b8}):

\begin{enumerate}
\item \emph{Orchard restriction} (context-free validity check): ``A
  transaction which contains any \emph{staking actions} must not contain
  any other fields contributing to or withdrawing from the Chain Value
  Pool Balances \emph{except} Orchard actions, explicit transaction fees,
  and/or explicit ZEC burning fields.''
\item \emph{Staking epochs}: ``Once Crosslink activates at block height
  \texttt{H\_crosslink\_activation}, \emph{staking epochs} begin, which
  are continguous [sic] spans of block heights with two phases:
  \emph{staking day} and the \emph{locked phase}.'' The staking-day length
  is a protocol constant of roughly 24 hours of blocks; the locked-phase
  length is a constant of roughly $6 \times 24$ hours, ``so that
  \emph{staking day} is approximately one day per week.''
\item \emph{Locked staking actions restriction}: ``If a transaction
  includes any staking actions \emph{except for} redelegate, and that
  transaction is in a block height in a \emph{locked phase} of the
  \emph{staking epoch}, then that block is invalid.''
\item \emph{Unbonding delay}: ``If a transaction is [sic] block height
  $H$ includes any \emph{claim} staking actions which refer to
  \emph{unbonding state} positions which have not existed in the unbonding
  state for at least one full staking epoch, that block is invalid.'' The
  source notes the implication that the same staking day cannot include
  both an \emph{unbond} and a \emph{claim} for one position.
\item \emph{Amount quantization}: ``If a transaction includes any staking
  actions with an amount which is not a power of 10 ZEC (or ZAT), that
  transaction is invalid.'' The book's TODO list clarifies the intent that
  the restriction apply only to the \emph{stake} (create-position) action.
\end{enumerate}

The quantization is a privacy device. The blog states it plainly:
``Staking transactions must be made in exponential increments of 10
(e.g., 1, 10, 100, 1,000, 10,000, etc.)''; ``This approach prevents
observers from easily inferring individual staking patterns, preserving
privacy while maintaining accountability''; and, on the other side of the
trade, ``For security and transparency, staked amounts are revealed so
the community can verify the total ZEC staked and its distribution across
finalizers'', with ``Transactions... processed in fixed 5-day epochs,
grouping all commitments into a single batch'' (Updated Staking Design,
2025-10-27/11-05). Quantisation of shielded weight has precedent in this
series: the shielded voting protocol quantises note value into
$0.125$-ZEC ballots for the same linkability reason (\emph{Voting Guide},
\S``Ballots and eligibility'').

Two gaps are visible already at this level. First, the \emph{stake}
action assigns ``a transparent ZEC amount'' while staking is exclusively
Orchard-pool: how the revealed, quantized amount balances against Orchard
value fields under \S4.17 is unspecified, and the prototype sidesteps it
entirely --- a transaction carrying a staking action bypasses the
\texttt{has\_inputs\_and\_outputs} consensus check with a ``TODO: real
staking transactions with inputs/outputs''
(\texttt{zebra-consensus/src/transaction/check.rs} at \texttt{6d02a1b8}).
Second, the phrase ``a power of 10 ZEC (or ZAT)'' leaves the bucket range
ambiguous: whether sub-1-ZEC buckets ($10^k$ zatoshi) are permitted, and
hence the minimum stake, is unspecified; the blog's examples start at
1~ZEC. How staking actions pay fees, and whether a transaction may carry
several, are open TODOs in both book and implementation
(\texttt{book/src/design/nutshell.md} at \texttt{6d02a1b8}).

The delegation-safety goals for the first deployment are recorded as
issues: ``Delegating to a finalizer does not enable the finalizer to
steal your funds'' and delegation ``does not leak information that links
the user's action to other information about them, such as their IP
address, their other ZEC holdings that they are choosing not to stake, or
their previous or future transactions'' (GH \#124, implementation book,
\texttt{book/src/design/scoping.md} at \texttt{6d02a1b8}). The explicit
first-deployment trade-offs cut the other way: \emph{not} ``supporting a
large number of finalizers so that any user who wants to run their own
finalizer can do so'', no support for casual users running finalizers,
and no liquid staking derivatives --- the last excluded in both the
tfl-book goals and the blog's V1 (ibid.; tfl-book,
\texttt{src/design/goals.md} at \texttt{fe6e1d6f}).

\paragraph{Classification.} \emph{Designed but unspecified}: the actions
and rules above are draft consensus text in a self-described provisional
sketch. The transaction format (value balancing), fee payment, bucket
range, and action multiplicity are \emph{open problems} acknowledged in
the sources.

\subsection{Epochs and parameter drift}\label{sec:xl-stake-params}

The staking-time parameters have changed at every publication and the
sources now disagree. Table~\ref{tab:xl-stake-params} pins each statement
to its date.

\begin{table}[htbp]
\centering
\small
\begin{tabular}{@{}p{3.6cm}p{3.6cm}p{2.9cm}p{4.1cm}@{}}
\toprule
Source, date & Epoch & Unbonding & Also fixed \\
\midrule
Early Tokenomics Design (forum t/52154, 2025-09-12)
  & 10-day staking / batching windows
  & 30 days
  & top-100 roster; 40/40 issuance split; no protocol slashing;
    finalizer commissions expected \\
Updated Staking Design (forum 2025-10-27; blog 2025-11-05)
  & fixed 5-day epochs, one batch per epoch
  & 1 epoch; funds available in 5--10 days
  & Orchard-pool exclusivity; power-of-10 amounts \\
Mechanism-design audit proposal (2025-12-11)
  & 5-day epochs
  & 5--10 days
  & 40/40 split; ``no protocol-level slashing in V1'' \\
Implementation book (\texttt{6d02a1b8}, snapshot 2026-04-21)
  & staking day ($\sim$1 day) $+$ locked phase ($\sim$6 days)
  & $\geq 1$ full staking epoch
  & redelegate exempt from the locked phase \\
\bottomrule
\end{tabular}
\caption{Staking-parameter statements by source and date. Sources:
forum t/52154 (2025-09-12); Shielded Labs blog (2025-10-27/11-05);
\texttt{book/src/design/analysis/history/2025-12-11-mech-design/proposal.md}
and \texttt{book/src/design/nutshell.md}, both at \texttt{6d02a1b8}. Web
sources fetched 2026-07-15.}
\label{tab:xl-stake-params}
\end{table}

The blog shortened the September windows ``citing liquidity and
accessibility''; the implementation book --- which post-dates the blog ---
lengthens the effective cycle again to approximately one staking day per
week. Which parameterization is current is unresolved. The divergence is
not cosmetic on one point: the book's locked-phase rule exempts
\emph{redelegate}, making redelegation valid at any height, while the
blog has participants ``reallocate their stake to another finalizer
during these windows'', implying batched redelegation. Redelegation is
the sole protocol-level punishment mechanism
(\S\ref{sec:xl-stake-slashing}), so its latency is security-relevant.

\paragraph{Classification.} Each row is \emph{designed but unspecified}
with a date; the current values are an \emph{open problem}. A ZIP must
pin epoch length, staking-day length, and unbonding duration in blocks,
and must fix redelegation latency explicitly.

\subsection{Rewards and issuance}\label{sec:xl-stake-rewards}

The implementation book routes staking rewards through the coinbase, gated
on finality: consensus rules require the coinbase transaction to transfer
the Crosslink-reward portion of new issuance \emph{only} in blocks that
advance finality. Per-block Crosslink rewards otherwise accumulate as
\emph{pending finalization rewards}, and a finality-updating block must
distribute the entire pending amount to stakers and finalizers, computed
against the active set as of the most recent final block prior to the
update (\texttt{book/src/design/nutshell.md} at \texttt{6d02a1b8}). With
$c := \texttt{COMMISSION\_FEE\_RATE} = 0.1$ a fixed protocol parameter,
the slices of the pending amount $S_{\mathsf{total}}$ are
%
\begin{equation}\label{eq:xl-reward-slices}
S_{\mathsf{finalizer}}
  = c \cdot \frac{W_{\mathsf{finalizer}}}{W_{\mathsf{total}}}
\qquad\text{(active-set finalizers only)},
\qquad
S_{\mathsf{staker}}
  = (1 - c) \cdot \frac{W_{\mathsf{staker}}}{W_{\mathsf{total}}},
\end{equation}
%
where $W_{\mathsf{participant}}$ is that participant's total stake weight
and $1 - c = 0.9$ is called the staker reward rate. The staker weight
$W_{\mathsf{staker}}$ counts all of a staker's positions regardless of
whether they are delegated to an active-set finalizer; the sum of all
slices is $S_{\mathsf{total}}$ or less, ``with the difference coming from
stake assigned to finalizers outside the active set'' (ibid.).

The issuance side is thinner. The Early Tokenomics post fixes ``the block
reward should be split equally (40/40) between miners and finalizers''
(forum t/52154, 2025-09-12), and the mechanism-design audit proposal
confirms the ``40/40 issuance split'' as an audit assumption
(\texttt{book/src/design/analysis/history/2025-12-11-mech-design/proposal.md}
at \texttt{6d02a1b8}). The destination of the remaining 20\% is not stated
in any source we located.
%% TODO: presumably the ZIP-1015 lockbox/dev-fund allocation; verify
%% before asserting -- no Crosslink source says so.
How Crosslink rewards interact with the issuance schedule, halvings, and
the 21M cap is exactly the ``Issuance and supply mechanics'' component the
tfl-book lists as missing (tfl-book,
\texttt{src/introduction/status-and-next-steps.md} at \texttt{fe6e1d6f}).
The commission model is also unsettled: the book fixes $c = 0.1$ for all
active-set finalizers, but the September post said ``We also expect
finalizers to charge commissions'', suggesting per-finalizer market rates
(forum t/52154, 2025-09-12).

One UX goal collides with the quantization rule. Goal GH \#158 asks for
compounding returns without user or wallet action (implementation book,
\texttt{book/src/design/scoping.md} at \texttt{6d02a1b8}), but a reward
paid to an Orchard destination cannot silently increase a power-of-10
position. No source resolves the collision; the book's own TODO list
flags ``rewards distribution via coinbase'' and ``quantization of amounts
/ time'' as unresolved (\texttt{book/src/design/nutshell.md} at
\texttt{6d02a1b8}).

\paragraph{Classification.} \emph{Designed but unspecified}:
Equation~\eqref{eq:xl-reward-slices} and the pending-rewards rule are
draft consensus text. Issuance interaction, the 20\% remainder, the
commission model, reward eligibility across the unbond--claim lifecycle,
and the compounding mechanism are \emph{open problems}. The prototype
implements none of this (\S\ref{sec:xl-stake-prototype}).

\subsection{Slashing, and its deliberate absence}
\label{sec:xl-stake-slashing}

The design commitment is verbatim and unambiguous: ``There will be no
automatic, protocol-level slashing in this design, so the power lies
fully in the stakers' hands to reward and censure node operators''
(forum t/52154, 2025-09-12); the audit proposal restates the constraint
as ``no protocol-level slashing in V1''
(\texttt{book/src/design/analysis/history/2025-12-11-mech-design/}
\texttt{proposal.md} at \texttt{6d02a1b8}).

The ECC book never specified slashing either; it appears only as
discussion. A ``Potential changes'' passage observes that ``In most PoS
protocols, the requirement to have a minimum amount of stake in order to
make a proposal acts as a gatekeeping filter on proposals, and
potentially allows parties that make invalid proposals to be slashed'',
that stake-holding validators ``can be slashed'', and that ``there is no
technical reason why the ability to make a bft-proposal has to be
gatekept by a stake requirement --- given the situation of Zcash in which
we already have a mining infrastructure'' (tfl-book,
\texttt{src/design/crosslink/potential-changes.md} at \texttt{fe6e1d6f}).
The closest the book comes to accountability machinery is a sketch in its
questions file: a \emph{bft-final-vee} --- two final bft-blocks with the
same parent --- as objective rollback evidence; the observation that more
than $1/3$ of stake voting for conflicting blocks in one epoch proves the
adversary-stake assumption was violated; and the note that such evidence
``can be posted in a transaction to the bc-chain'', optionally latching
Stalled Mode until manual intervention (tfl-book,
\texttt{src/design/crosslink/questions.md} at \texttt{fe6e1d6f}). None of
this is adopted as a rule.

What replaces slashing is a chain of accountability between roles, stated
in the implementation book's Five Component Model
(\texttt{book/src/design/five-component-model.md} at \texttt{6d02a1b8}).
Stakers police finalizers: it is ``the job of the stakers to hold the
finalisers accountable by detecting and punishing misbehavior... by
redelegating stake away from ill-behaved finalizers to well-behaved
finalizers''. Users police stakers, ultimately via ``an emergency
user-led hard fork (which is what happened in the STEEM/HIVE hard
fork)''. The same file bounds finalizer power: finalizers ``cannot
unilaterally compromise liveness or censorship-resistance (even if an
attacker controls $>2/3$ of the stake-weighted votes)''; ``A plurality of
the finalisers (controlling $>1/3$ of the stake-weighted votes) could
execute stall attacks''; and finalizers ``cannot unilaterally execute
rollback attacks even though they can prevent rollback attacks''.

The chain has a load-bearing dependency on censorship-resistance of block
production. A collusion controlling both $>1/2$ hashpower and $>2/3$
stake could violate finality and execute rollback and unavailability
attacks; and ``Censorship-resistance is necessary for stakers to stake,
unstake, and redelegate, which means censorship-resistance is necessary
for stakers to be able to perform their duty of holding the finalizers
accountable. Therefore, if an attacker controls $>1/2$ of the hashpower
(and is thereby able to censor), then the attacker can prevent the
stakers from holding the finalizers accountable'' (ibid.; the file ends
mid-list on the ``$>2/3$ stake and $<1/2$ hashpower'' scenario --- the
document is incomplete). The reliance on stakers is so central that the
zebra-crosslink design drops the tfl-book's Stalled Mode entirely: ``This
simplification to the protocol rules follows the rationale that we are
already relying heavily on the stakers to guard safe operation of the
protocol by redelegating appropriately to prevent hazards like BFT
stalls'' (\texttt{book/src/design/cl2-construction.md} at
\texttt{6d02a1b8}).

Shielded Labs frames the substitution quantitatively as
Penalty-for-Failure-to-Defend (PFD): penalties are assessed per property
($\mathrm{PFD}_{\mathrm{liveness}}$ versus
$\mathrm{PFD}_{\mathrm{finality}}$); Tendermint-style slashing yields a
``financial value gauge'' PFD, while PoW --- and Crosslink's
slashing-free PoS --- yields a \emph{rate} PFD, the loss of future
revenue rate; the two types ``cannot be directly compared''
(\texttt{book/src/design/analysis/pffd.md} at \texttt{6d02a1b8}, a paste
of GH issue \#264). Redelegation-based punishment is thus a rate penalty
standing in for the gauge penalty of slashing. For redelegation to work,
stakers must be able to observe misbehaviour: Crosslink finality ``is
accountable because the chain contains explicit voting records which
indicate each participating finalizer's on-chain behavior''
(\texttt{book/src/design/terminology.md} at \texttt{6d02a1b8}), and goal
GH \#214 asks that ``Skilled users can easily collect records of
finalizer behavior and performance''
(\texttt{book/src/design/scoping.md} at \texttt{6d02a1b8}) ---
accountability is informational, feeding redelegation, not punitive. The
terminology echoes the accountability line of the Ebb-and-Flow
literature: the construction cites Neu--Tas--Tse [NTT2020, arXiv
2009.04987] as background (tfl-book,
\texttt{src/design/crosslink/construction.md} at \texttt{fe6e1d6f}), and
the same authors' accountability-gadgets paper (arXiv 2105.06075,
Financial Cryptography 2022) is the natural successor for this topic;
NTT2020 remains at v3 (2021-02-04) as of 2026-07-15.

The backstop of the chain --- the user-led emergency hard fork --- has a
drafted scope. The governance proposal limits legitimate emergency
intervention to layer-1 security failures, explicitly including
``situations in which would-be miners, finalizers, or stakers are being
systematically censored or prevented from participating in block
production, finalization, or staking''; it excludes ``staking or
unbonding parameters'' changes from emergency scope; and it states that
``Control over a large share of hashpower, finalization power, or stake
does not itself constitute a violation''
(\texttt{book/src/design/userled-emergency-hardforks.md} at
\texttt{6d02a1b8}).

\paragraph{Classification.} The absence of slashing is \emph{designed but
unspecified} --- a dated, verbatim design commitment confirmed by the
audit constraints. The accountability chain and PFD framing are design
rationale at the same bin. The tfl-book's evidence sketches
(bft-final-vee, conflicting-vote proofs) are an \emph{open problem}:
whether future versions add evidence-based slashing, and what a ZIP must
say about accountability records, is unresolved --- recall that the
tfl-book's own delegation goal budgets for ``slashing fees''
(\S\ref{sec:xl-stake-abstract}).
%% RESOLVED 2026-07-15: findings published at
%% nikete.com/crosslink-zebra-audit (verified); repo holds only the
%% pre-registration. Delegation/penalty consequences are covered in
%% \S\ref{sec:xl-econ}; no design changes adopting the recommendations
%% are visible at 6d02a1b.

\subsection{Finalizer keys}\label{sec:xl-stake-keys}

Definition~\ref{def:xl-roster} makes the finalizer verification key the
sole on-chain reference to a finalizer, and makes position keys
deliberately unlinkable to a user's other activity. Everything else about
finalizer key management is prototype fact or gap. The implementation
uses single ed25519 keys per finalizer (ed25519-zebra; ``ed25519 public
key of the finalizer''), and votes are plain ed25519 signatures: the vote
wire format is 76 bytes --- a 32-byte ed25519 finalizer public key, a
32-byte BLAKE3 value hash (all-zero encoding Nil), an 8-byte height, and
a 4-byte round whose most significant bit flags a commit --- with a
signed vote appending a 64-byte ed25519 signature over those 76 bytes
(\texttt{zebra-crosslink/src/lib.rs}, \texttt{MalVote}, lines
1811--2052, and \texttt{zebra-chain/src/block/header.rs}, lines
109--133, at \texttt{6d02a1b8}). No FROST
or threshold-signing design for finalizer keys is sourced anywhere in the
Crosslink corpus: a search for ``frost'' over tfl-book,
crosslink-protocol-guide, crosslink-spec, and zebra-crosslink returns
only upstream Zebra RedPallas/Orchard files unrelated to finalizers
(negative grep, 2026-07-15). Threshold custody is established practice
elsewhere in the Zcash stack --- FROST per RFC~9591 in the PCZT signing
flow (\emph{Wallet Guide}, \S``Partially created transactions (PCZT)'')
--- which makes its absence here a visible gap rather than an exotic ask.

The gap matters because of Definition~\ref{def:xl-nonhonest}'s caveat:
honest validator keys must remain secret \emph{indefinitely}, with secure
deletion on rotation (tfl-book,
\texttt{src/design/crosslink/construction.md} at \texttt{fe6e1d6f},
issue \#140) --- yet no key-rotation or revocation mechanism for
finalizer verification keys in the roster is specified in any source.
Bootstrap is equally unspecified: no source examined (tfl-book,
implementation book, either blog post) defines eligibility or formation
of the initial finalizer set. The prototype bootstraps ad hoc: the roster
starts empty; keys are deterministically derived from the config field
\texttt{insecure\_user\_name} (``temp seed for private/public key
pair''); test instrumentation can force-include roster members
(\texttt{TFInstr::ROSTER\_FORCE\_INCLUDE}); and a node adds itself with
\texttt{voting\_power} 0 when absent from the roster it hands to the BFT
engine (\texttt{zebra-crosslink/src/lib.rs},
\texttt{zebra-crosslink/src/test\_format.rs} at \texttt{6d02a1b8}).

\paragraph{Classification.} The ed25519 key and vote formats are
implementation-defined prototype facts. Key rotation, revocation,
threshold custody, and initial-roster formation (including any minimum
self-stake, and eligibility beyond top-$K$ by stake weight) are
\emph{open problems} in every design source.

\subsection{The prototype}\label{sec:xl-stake-prototype}

The zebra-crosslink prototype (@ \texttt{6d02a1b8}, 2026-04-21)
implements a reduced, implementation-defined variant of the above; we
record it as deployed-code fact about a prototype, not as design. The
placeholder activation height is \texttt{TFL\_ACTIVATION\_HEIGHT} $=$
\texttt{BlockHeight(2000)} (\texttt{zebra-crosslink/src/lib.rs}).

\emph{Roster mechanics.} Staking state changes take effect only when
their PoW block is finalized: function
\texttt{update\_roster\_for\_block()} runs while walking each newly
finalized block and applies \texttt{StakingAction} commands from
\texttt{Transaction::VCrosslink} transactions. The variants of
\texttt{StakingActionKind} are \texttt{Add}, \texttt{Sub},
\texttt{Clear}, \texttt{Move}, and \texttt{MoveClear} --- names diverging
from the book's stake/redelegate/unbond/claim --- with fields including
\texttt{insecure\_target\_name} and \texttt{insecure\_source\_name},
defined in the Shielded Labs librustzcash fork (git rev
\texttt{33dc74bf}; not present upstream)
(\texttt{zebra-crosslink/src/lib.rs},
\texttt{zebra-chain/src/transaction.rs}, \texttt{Cargo.toml} at
\texttt{6d02a1b8}). The BFT engine is Tenderlink --- Shielded Labs
``moved away from Malachite and now maintain our own lightweight version
of Tendermint, called Tenderlink'' (Updated Staking Design,
2025-10-27/11-05; dependency \texttt{tenderlink 0.1.0}, rev
\texttt{0e885095}) --- and the roster handed to it is sorted by
descending (stake, public key) --- ``Needs to uniquely \& stably
tie-break members with the same stake'' --- with cumulative stake
tracked ``to make weighted round-robin easy'' for proposer selection
(\texttt{zebra-crosslink/src/lib.rs}, \texttt{Cargo.toml} at
\texttt{6d02a1b8}).

\emph{Rewards.} The prototype ignores the coinbase design entirely: a
constant $R = \texttt{pos\_total\_reward} = 6000$~zats (``arbitrary
scale'') is apportioned per newly finalized PoW block as
$\mathsf{reward}_i = \lfloor v_i \cdot R / \sum_j v_j \rfloor$ over
finalizer voting powers $v_i$, with the integer-division dust going to
the largest-stake finalizer, and the result added directly to
\texttt{voting\_power} --- auto-compounding into stake weight rather
than paid via coinbase. The code acknowledges its artifacts: ``the
compounding is per PoW block'' and ``finalizers added in the same PoS
will get rewards for PoW blocks they couldn't have actually voted on''
(\texttt{zebra-crosslink/src/lib.rs} at \texttt{6d02a1b8}).

\emph{Gaps against the draft rules.} (1)~No power-of-10 quantization
check, no staking-epoch/staking-day/locked-phase gating, and no
unbonding-delay enforcement exist anywhere in the staking path. (2)~No
top-$K = 100$ active-set truncation: every roster member with non-zero
power participates. (3)~Staking transactions are stubs: a
\texttt{VCrosslink} transaction carrying a staking action bypasses the
\texttt{has\_inputs\_and\_outputs} check, so no value balancing against
the chain value pools occurs. (4)~The bonding lifecycle (unbond/claim)
is not implemented; \texttt{Sub}/\texttt{Clear} remove weight
immediately (\texttt{zebra-consensus/src/transaction/check.rs},
\texttt{zebra-crosslink/src/lib.rs} at \texttt{6d02a1b8}; negative grep
for quantiz/unbond/epoch, 2026-07-15).

\emph{Status.} Crosslink remains in feature-net validation: the
Incentivized Crosslink Feature Net launched with Milestone 5 (completed
2026-04-15; Season 1 from 2026-04-15; 500 ZEC allocated overall, 25 ZEC
for Season 1; ``Only cTAZ earned via block rewards counts''), following
Milestone 4's independent security review (concluded 2026-02-10). No NU8
activation date is stated (deployment-updates thread, fetched
2026-07-15).

\paragraph{Classification.} Implementation-defined throughout, and
self-describedly so; the book warns that the in-house BFT protocol's
``own security properties need to be more thoroughly verified''
(\texttt{book/src/design.md} at \texttt{6d02a1b8}). None of this
prototype behaviour constrains the eventual ZIP.

\subsection{What a specification must pin down}\label{sec:xl-stake-open}

We close with the open problems, each traceable to a gap documented
above.

\begin{enumerate}
\item Epoch parameterization: epoch length, staking-day length, and
  unbonding duration, in blocks
  (Table~\ref{tab:xl-stake-params}).
\item Redelegation latency: locked-phase exemption versus batched
  windows --- the timing of the sole protocol-level punishment
  (\S\ref{sec:xl-stake-params}).
\item The staking transaction format: how a revealed, quantized,
  ``transparent'' amount balances against Orchard value fields under
  protocol specification \S4.17 (\S\ref{sec:xl-stake-actions}).
\item Quantization bucket range and minimum stake: ``power of 10 ZEC
  (or ZAT)'' (\S\ref{sec:xl-stake-actions}).
\item Commission: protocol-fixed $c = 0.1$ versus per-finalizer market
  rates (\S\ref{sec:xl-stake-rewards}).
\item Issuance: the 40/40 split's remaining 20\%, pending-reward
  interaction with halvings and the 21M cap
  (\S\ref{sec:xl-stake-rewards}).
\item Reward mechanics: coinbase distribution versus compounding, and
  the GH \#158 goal against power-of-10 positions
  (\S\ref{sec:xl-stake-rewards}); reward eligibility across the
  unbond--claim lifecycle and for stake on non-active finalizers.
\item Fees and multiplicity of staking actions per transaction
  (\S\ref{sec:xl-stake-actions}).
\item Voting-unit mapping and the per-epoch roster snapshot
  (\S\ref{sec:xl-stake-abstract}).
\item Initial-roster formation, minimum self-stake, and finalizer
  eligibility beyond top-$K$ (\S\ref{sec:xl-stake-keys}).
\item Finalizer key rotation and revocation under the
  indefinite-secrecy requirement (\S\ref{sec:xl-stake-keys}).
\item Accountability records, and whether any future version adopts
  evidence-based slashing (\S\ref{sec:xl-stake-slashing}).
\item A formal model: nothing in the machine-checked corpus covers the
  roster, delegation, epochs, unbonding, or reward arithmetic
  (crosslink-spec @ \texttt{8edc3e94} models proposal validity only).
\end{enumerate}

%% RESOLVED 2026-07-15: quint-crosslink is empty upstream (git ls-remote
%% returns zero refs), so item 13 stands as stated. The mechanism-design
%% audit findings are published at nikete.com/crosslink-zebra-audit and
%% covered in \S\ref{sec:xl-econ}; no delegation/penalty design changes
%% are visible at 6d02a1b.

\section{The wallet-visible contract}\label{sec:xl-contract}

Crosslink changes what a Zcash node presents to everything above it. Today a
wallet, an exchange, or the voting stack consumes one object --- the node's
best chain --- and layers its own confirmation policy on top (\emph{Wallet
Guide}, \S``The proposal''). Under Crosslink~2 the node maintains two chains
and a per-block finality status, and every consumer must decide which chain
it reads and what it does when the two stop tracking each other. This
section fixes that contract as the sources state it: the two ledger views
and the theorems relating them (\S\ref{sec:xl-views}), finality as a
queryable status (\S\ref{sec:xl-finality-status}), the interface the
prototype implements today (\S\ref{sec:xl-node-interface}), and the four
wallet-facing decisions no source resolves --- anchor policy
(\S\ref{sec:xl-anchors}), confirmation policy (\S\ref{sec:xl-confpolicy}),
stall behaviour (\S\ref{sec:xl-stalls}), and the service-facing protocol
(\S\ref{sec:xl-services}).

Two texts carry the design (pinned, with the rest of the corpus, in
Table~\ref{tab:xl-sources}). The ECC Trailing Finality Layer book
(\texttt{zcash/tfl-book} @ \texttt{fe6e1d6f}, 2024-12-13; hereafter the
\emph{tfl-book}) defines the Crosslink~2 construction and has not changed
since; its own status page calls the design ``an early and incomplete
protocol design proposal'' (\texttt{src/introduction/status-and-next-steps.md}).
The Shielded Labs design book inside the implementation repository
(\texttt{zebra-crosslink} @ \texttt{6d02a1b8}, 2026-04-21; hereafter the
\emph{SL~book}) pastes the construction verbatim, declares divergences in
admonitions, and self-describes as a provisional precursor to ``a canonical
Zcash Improvement Proposal submission'' (\texttt{book/src/design/nutshell.md}).
No such ZIP exists: a repository-wide search of \texttt{zcash/zips} @
\texttt{69610984} (2026-07-09) finds no Crosslink deployment artifact and no
NU8 artifact of any kind (cf.\ \emph{Tachyon Guide}, \S``Relation to
Crosslink and the upgrade environment''). Shielded Labs' productionization
phase runs Q2~2026--Q2~2027 with no announced mainnet activation
(shieldedlabs.net/roadmap and forum thread t/49706, fetched 2026-07-15).
Nothing in this section is therefore \emph{specified}; claims are
\emph{designed but unspecified} at best, prototype code is cited as code,
and the open problems are named as such.

\subsection{Two ledger views}\label{sec:xl-views}

We write $\preceq_{\mathrm{bc}}$ for the prefix order on bc-valid chains,
$B \sim_{\mathrm{bc}} C$ when one of $B, C$ is a prefix of the other, and
$C^{\lceil k}$ for chain $C$ with its last $k$ blocks removed. The
bc-best-chain of node $i$ at time $t$ is $\mathsf{ch}_i^t$, notation the
tfl-book adopts from the ebb-and-flow model of Neu, Tas, and Tse (arXiv
2009.04987v3, 2021-02-04, still the latest revision as of 2026-07-15;
tfl-book \texttt{construction.md}:5).

\begin{definition}[Ledger views of Crosslink~2]\label{def:xl-ledger-views}
Fix the bc-confirmation-depth $\sigma \in \N^{+}$. At every time $t$, each
node $i$ derives two node-local chains from its bc-best-chain and its BFT
context:
\begin{enumerate}[label=(\roman*)]
\item the \emph{finalized chain} $\mathsf{localfin}_i^t$: a bc-valid chain
  obtained at the fixed confirmation depth $\sigma$ via the finalization
  update rule (Construction~\ref{def:xl-updatefin}, applied to the
  candidate of Definition~\ref{def:xl-lf-candidate});
\item the \emph{bounded-available chain}, for a per-node depth $\mu$ with
  $0 < \mu \le \sigma$:
  \[
    (\mathsf{localba}_{\mu})_i^t =
    \begin{cases}
      \mathsf{ch}_i^t{}^{\lceil\mu}
        & \text{if } \mathsf{localfin}_i^t \preceq_{\mathrm{bc}}
          \mathsf{ch}_i^t{}^{\lceil\mu},\\[2pt]
      \mathsf{localfin}_i^t & \text{otherwise.}
    \end{cases}
  \]
\end{enumerate}
(tfl-book \texttt{src/design/crosslink/construction.md}:391--399,
520--527.)
\end{definition}

Both views are plain block chains. Snap-and-Chat and Crosslink~1 instead
exposed \emph{sanitized transaction logs} --- flattened sequences with
invalid transactions filtered out --- and the tfl-book drops that machinery
``because we do not need sanitization, there is no need to express it as
a log of transactions rather than blocks''
(\texttt{construction.md}:393);
\S\ref{sec:xl-anchors} returns to why this matters to wallets. The
otherwise-case of the definition covers post-reorg difficulty situations in
which a new best chain outscores the old one in fewer than $\sigma$ blocks,
and it makes the following theorem trivial
(\texttt{construction.md}:547--555).

\begin{theorem}[Ledger prefix property]\label{thm:xl-ledger-prefix-wallet}
For any node $i$ that is honest at time $t$, and any confirmation depth
$\mu$: $\mathsf{localfin}_i^t \preceq_{\mathrm{bc}}
(\mathsf{localba}_{\mu})_i^t$.
\end{theorem}
\begin{proof}
By construction of $(\mathsf{localba}_{\mu})_i^t$: the otherwise-branch
returns $\mathsf{localfin}_i^t$ itself (tfl-book
\texttt{construction.md}:531--536).
\end{proof}

The depth $\mu$ is the node's own choice; the default ``should be'' $\mu =
\sigma$, and choosing $\mu < \sigma$ ``is at the node's own risk and may
increase the risk of rollback attacks against
$(\mathsf{localba}_{\mu})_i^t$ (it does not affect
$\mathsf{localfin}_i^t$)'' (\texttt{construction.md}:395--399, Security
caveat admonition).

\paragraph{Monotonicity and hazards.} The finalized view never moves
backward: the update rule advances $\mathsf{localfin}_i$ to the new
candidate only when that candidate descends from the current value, which
yields \emph{Local finalization linearity} --- the time series
$\mathsf{localfin}_i^{r \le t}$ is bc-linear. If the candidate instead
conflicts, the node keeps its old value and ``should record a finalization
safety hazard'', which ``can only happen if global safety assumptions are
violated''; the record should include $\mathsf{ch}_i^t$ and the history of
$\mathsf{localfin}$ updates (\texttt{construction.md}:462--488). Restricted
to a single node, Assured Finality
(Definition~\ref{def:xl-assured-finality-wallet})
is equivalent to this linearity; across nodes, linearity is necessary but
not sufficient (\texttt{construction.md}:466--473, 508).

\paragraph{The syncing gate.} Neither view is exposed to a client before
the node has synced: ``this chain state should not be exposed to clients of
the node until it has synced'', where synced means a baked-in checkpoint
bft-block $B_{\mathrm{checkpoint}}$ satisfies $B_{\mathrm{checkpoint}}
\le_{\mathrm{bft}} \mathsf{LF}(\mathsf{ch}_i^t)$,
$\mathsf{snapshot}(B_{\mathrm{checkpoint}}) \preceq_{\mathrm{bc}}
\mathsf{localfin}_i^t$, and the timestamp of $\mathsf{localfin}_i^t$ is
within a threshold of the current time
(\texttt{construction.md}:464, 529, 557--562).

\paragraph{Latency, and which view an application reads.} In the steady
state the finalized view trails the tip by about $\mu + 1 + \sigma$
bc-blocks --- at most ``slightly more than a factor of 2'' over
Snap-and-Chat's $\sigma_{\mathrm{sac}}$ when $\mu = \sigma =
\sigma_{\mathrm{sac}}$ (tfl-book \texttt{security-analysis.md}:265--269;
\texttt{questions.md}:27--31). The book states the application-facing
choice explicitly: since $\mathsf{localfin}$ lags
$\mathsf{localba}_{\sigma}$ by only a few blocks outside Stalled Mode,
``the main factor influencing the choice of a given application to use
$\mathsf{localfin}_i^t$ or $(\mathsf{localba}_{\mu})_i^t$ is not the
average latency, but rather the desired behaviour in the case of a
finalization stall: i.e.\ stall immediately, or keep processing user
transactions until $L$ blocks have passed''
(\texttt{construction.md}:415). Stalls are \S\ref{sec:xl-stalls}.

\paragraph{Parameters.} Production values of $\sigma$ and of the
finalization gap bound $L$ are unchosen. The tfl-book requires $L$
``significantly greater than $\sigma$'' and ``at least $2\sigma$'' in
practice (\texttt{construction.md}:391, 405); the prototype ships
$\{\sigma = 3,\ L = 7\}$ with the caveat ``No verification has been done on
the security or performance of these parameters''
(\texttt{zebra-crosslink/src/chain.rs}:206--222). Wall-clock statements
are additionally provisional: block target spacing is 75\,s today, and ZIP
218 (Status: Draft) proposes 25\,s at NU7, describing itself as
``complementary to \ldots{} finality mechanisms such as Crosslink''
(\texttt{zips/zip-0218.md}:3, 76--78).

\paragraph{Classification.} \emph{Designed but unspecified} throughout:
Definition~\ref{def:xl-ledger-views} and
Theorem~\ref{thm:xl-ledger-prefix-wallet} come from a design book at pre-ZIP
rigor, and the constants are prototype code, not consensus.

\subsection{Finality as a status}\label{sec:xl-finality-status}

The tfl-book's term of art is \emph{assured finality}: a protocol property
that assures transactions cannot be reverted \emph{by that protocol}. It
eschews ``absolute finality'' because ``it is not feasible for any protocol
to prevent reversing final transactions `out of band'\,''; ``final'' in the
book always means assured finality, in contrast to proof-of-work's
probabilistic finality (tfl-book \texttt{src/terminology.md}:11--13, 19).

\begin{definition}[Assured Finality]\label{def:xl-assured-finality-wallet}
An execution of Crosslink~2 has \emph{Assured Finality} iff for all times
$t, u$ and all nodes $i, j$ (potentially the same) such that $i$ is honest
at time $t$ and $j$ is honest at time $u$: $\mathsf{localfin}_i^t
\sim_{\mathrm{bc}} \mathsf{localfin}_j^u$ (tfl-book
\texttt{construction.md}:501--506).
\end{definition}

\begin{theorem}[Two-sided safety]\label{thm:xl-af-two-sided}
An execution of Crosslink~2 has Assured Finality if \emph{either} of the
following holds:
\begin{enumerate}[label=(\roman*)]
\item the $\Pi_{\mathrm{bc}}$ subprotocol has Prefix Agreement at
  confirmation depth $\sigma$ --- honest best-chains, truncated by
  $\sigma$, are pairwise compatible across all nodes and times; or
\item the $\Pi_{\mathrm{bft}}$ subprotocol has Final Agreement --- the
  last-final blocks of any two bft-valid blocks in honest view are
  pairwise compatible.
\end{enumerate}
(tfl-book \texttt{security-analysis.md}:72--77, 86--91;
hypothesis definitions at \texttt{construction.md}:337--340, 203--206.)
\end{theorem}

Assured Finality therefore survives the failure of either subprotocol's
safety assumption alone; it fails only under the weaker of the two, which
is the keystone guarantee a wallet's ``final'' badge would rest on. We
defer the proofs, which run through the Local fin-depth lemma and the
Linear prefix lemma, to the security analysis
(\S\ref{sec:xl-safety-theorems}); Prefix Agreement and Final Agreement
are stated in full as Definitions~\ref{def:xl-prefix-agreement}
and~\ref{def:xl-final-agreement}.

One caveat bounds what is actually proved. Below the marker ``\$TODO\$ Not
updated for Crosslink 2'' the tfl-book's safety chapter still argues in the
Snap-and-Chat sanitized-log formalism ($\mathrm{LOG}_{\mathrm{fin}}$,
$\mathrm{LOG}_{\mathrm{ba}}$, $\mathsf{san\text{-}ctx}$), and the
$\mathrm{LOG}_{\mathrm{ba}}$ agreement proof trails off unfinished (``What
we want to show is that, under some conditions on executions, \ldots'')
(\texttt{security-analysis.md}:54, 137--160, 253--257); the SL~book copies
this state verbatim. The bin split for this section:
Theorem~\ref{thm:xl-af-two-sided} is complete as stated;
transaction-log-level safety for the bounded-available view is
\emph{designed but unspecified} with an incomplete proof.

\paragraph{The Shielded Labs definition.} The SL~book defines ``crosslink
finality'' by five properties: \emph{accountable} (on-chain voting
records), \emph{irreversible within protocol scope} (undoing requires
software modification or human intervention), \emph{objectively
verifiable} (determinable from local block history alone), \emph{globally
consistent} (all sufficiently-synced nodes agree on finality status), and
an \emph{asymmetric cost-of-attack defense} that is a literal TODO in the
source (\texttt{book/src/design/terminology.md}, ``Crosslink Finality'').
Its ledger-state framing makes finality the sole ``height-eventual Ledger
State'': a property of height $H$ computable only from a later attesting
block, on which all observers of that attesting block agree ``even though
they may observe the attesting block rollback. In the event of such a
rollback, the protocol guarantees that an alternative block will
eventually attest to the finality status of the same candidate block'';
the General Ledger State Design Invariant requires all ledger state,
finality status and the proof-of-stake roster included, to be computable
entirely from PoW blocks and transactions
(\texttt{book/src/design/nutshell.md}:106--128). For light clients the
SL~book adds a negative constraint: its design drops the tfl-book's
``outer signature'', so ``a direct hash or copy of the most recent BFT
finality certificate is insufficient evidence'' that a given bitwise copy
is or will become canonical
(\texttt{book/src/design/cl2-construction.md}:3--16) --- any wire-format
finality proof must be designed around that.

\paragraph{Classification.}
Definition~\ref{def:xl-assured-finality-wallet} and
Theorem~\ref{thm:xl-af-two-sided} are \emph{designed but unspecified}
(complete statements and proofs, pre-ZIP rigor); the SL five-property
definition is \emph{designed but unspecified} with its fifth property an
acknowledged gap; ledger-level $\mathsf{localba}$ safety is an
\emph{open problem} in the sources' own accounting.

\subsection{The interface implemented today}\label{sec:xl-node-interface}

The prototype (\texttt{zebra-crosslink} @ \texttt{6d02a1b8}) exposes
finality through a typed internal service and a JSON-RPC surface. The
internal state service defines
\begin{center}
\texttt{enum TFLBlockFinality \{ NotYetFinalized, Finalized,
CantBeFinalized \}}
\end{center}
with request variants
\begin{center}
\texttt{IsTFLActivated}, \texttt{FinalBlockHeightHash},
\texttt{FinalBlockRx}, \texttt{SetFinalBlockHash},\\
\texttt{BlockFinalityStatus(height, hash)},
\texttt{TxFinalityStatus(txid)}, \texttt{Roster},\\
\texttt{FatPointerToBFTChainTip}, \texttt{StakingCmd}
\end{center}
(\texttt{zebra-state/src/crosslink.rs}:11--49; variant
\texttt{FinalBlockRx} is a broadcast receiver). The RPC methods, every
one doc-commented as ``Placeholder''
(\texttt{zebra-rpc/src/methods.rs}:246--420), are:
\begin{itemize}[itemsep=1pt]
\item \texttt{is\_tfl\_activated};
\item \texttt{get\_tfl\_roster\_zec} and \texttt{get\_tfl\_roster\_zats};
\item \texttt{get\_tfl\_fat\_pointer\_to\_bft\_chain\_tip};
\item \texttt{get\_tfl\_final\_block\_hash} and
  \texttt{get\_tfl\_final\_block\_height\_and\_hash};
\item \texttt{get\_tfl\_block\_finality\_from\_hash} and
  \texttt{get\_tfl\_tx\_finality\_from\_hash};
\item \texttt{set\_tfl\_finality\_by\_hash};
\item \texttt{subscribe\_tfl\_new\_final\_block\_hash}.
\end{itemize}
The docs concede that the final-block RPCs ``for the sake of testing
\ldots{} currently treat pre-reorg block as final''.

Block-level finality is computed as a tri-state: a height above the final
height returns \texttt{NotYetFinalized} (``N.B. this may be invalidated by
the time it is received''); otherwise the status is \texttt{Finalized} iff
the best-chain hash at that height equals the queried hash, else
\texttt{CantBeFinalized} (\texttt{zebra-crosslink/src/lib.rs}:1327--1374).
This matches the SL~book's globally-consistent claim only below the final
height, as it must: above it, the answer is a moving target.

Two behaviours contradict the SL~book's own definition of crosslink
finality, and they sit exactly on the surface a wallet would consume:
\begin{enumerate}[label=(\roman*)]
\item Transaction finality omits the fork check. Request
  \texttt{TxFinalityStatus} returns \texttt{Finalized} iff
  $\texttt{tx.height} \le \texttt{final\_height}$, with a literal
  \texttt{// TODO: CantBeFinalized} --- it never checks best-chain
  membership of the containing block, so a transaction on a losing fork
  below the final height is reported \texttt{Finalized}
  (\texttt{zebra-crosslink/src/lib.rs}:1412--1431). A blocking wait built
  on \texttt{FinalBlockRx} loops until some status other than
  \texttt{NotYetFinalized} is returned
  (\texttt{zebra-rpc/src/methods.rs}:2150--2195).
\item Method \texttt{set\_tfl\_finality\_by\_hash} is intended to be
  regtest-only but currently ships \texttt{regtest\_override = true}
  (\texttt{zebra-rpc/src/methods.rs}:1975--2000).
\end{enumerate}
Both are prototype-status defects rather than design statements, but a
finality status that is neither objectively verifiable nor globally
consistent is the current observable behaviour.

\paragraph{Before activation: the fallback depth.} Whenever no BFT-final
block is known (including before TFL activation), the implementation
falls back to the block-locator's last hash as its final block: function
\texttt{tfl\_reorg\_final\_block\_height\_hash} yields, in effect, the
block at depth \texttt{MAX\_BLOCK\_REORG\_HEIGHT}, which the fork pins
at 99, one less than the transparent coinbase maturity of 100; an
internal \texttt{latest\_final\_block}, once set by BFT, takes
precedence (\texttt{zebra-crosslink/src/lib.rs}:235--332;
\texttt{zebra-state/src/constants.rs}:14--31). Upstream Zebra meanwhile
moved the same constant to 1000 (commit \texttt{8767a91ba}, 2026-06-17),
documenting it as ``a local-only node policy; it is not part of
consensus'' sized ``as a defence-in-depth measure against sustained
consensus splits'' (\texttt{zebra} @ \texttt{9f3166b9},
\texttt{zebra-chain/src/parameters/constants.rs}:20--30).
The lesson for prose about pre-Crosslink ``de facto finality'': the
rollback depth is node policy, not consensus, and it currently differs by
a factor of ten between upstream and the very fork implementing
Crosslink. The \emph{Ironwood Guide} (\S``Proof of ownership of \emph{some}
valid note: the commitment tree'', paragraph ``Anchor choice as the tree
grows'') states that anchoring ``at least 100 blocks deep (Zcash's reorg
limit) avoids this entirely''; against upstream's 1000-block window that
sentence overstates a policy as a limit.
%% TODO: file a cross-volume erratum against the Ironwood Guide sentence
%% (ironwood-guide.tex:1005-1034) once the volume owner confirms wording.

\paragraph{Classification.} \emph{Specified as code} at prototype rigor
only --- every RPC is a self-described placeholder, and nothing in this
subsection is consensus.

\subsection{Anchor semantics: the undecided keystone}
\label{sec:xl-anchors}

Crosslink~2 leaves transaction validity alone: ``The consensus rules
applying to bc-transactions are entirely unchanged, including any rules
that depend on bc-block height or previous bc-blocks'', because the
construction never reorders or sanitizes transactions; contextual validity
in $\Pi_{\mathrm{bc}}$ is that of $\Pi_{\mathrm{origbc}}$ modulo the
non-contextual new block rules (tfl-book \texttt{construction.md}:267,
701--703). The Orchard anchor rule is therefore inherited as-is: a spend
may reference the treestate of \emph{any} main-chain block from Orchard
activation onward, with no sliding window (\emph{Ironwood Guide},
\S``Proof of ownership of \emph{some} valid note: the commitment tree'',
paragraph ``Anchor choice as the tree grows''; mainnet floor 1{,}687{,}104)
--- where ``main chain'' now means the $\mathsf{localba}$ history. No
source restricts spend anchors to the finalized prefix: not the tfl-book,
not the SL~book, not the implementation (grep over
\texttt{zebra-crosslink} @ \texttt{6d02a1b8}: no anchor-rule changes), and
not the ZIP corpus (@ \texttt{69610984}).

The wallet-level anchor policy under two ledgers is thus the single most
consequential undecided item in the contract. A future ZIP must decide
two things: (a)~whether consensus adds any anchor restriction --- nothing
is proposed anywhere --- and (b)~what a wallet's default anchor and
note-selection policy becomes relative to the finalized height. The
branch costs are concrete:
\begin{itemize}
\item \emph{Anchor only in $\mathsf{localfin}$.} Every spend inherits
  finalization latency --- about $\mu + 1 + \sigma$ blocks in the steady
  state (\S\ref{sec:xl-views}), up to $L$ during a stall --- and recently
  received notes cannot be spent until their treestate finalizes. In
  exchange, anchor invalidation becomes impossible whenever Assured
  Finality holds, eliminating the surviving-reorg caveat the \emph{Wallet
  Guide} records (\S``The proposal'': a rollback past the anchor
  invalidates the transaction while its revealed nullifiers link any
  retry).
\item \emph{Keep the tip-relative anchor.} ZIP 315's rule --- ``if the
  current block is at height H, the anchor SHOULD reflect the final
  treestate of the block at height H-3'' (\texttt{zips/zip-0315.rst}:586--587)
  --- preserves latency and the anonymity-set size and uniformity that
  motivate near-tip anchoring (\emph{Ironwood Guide}, ibid.), and retains
  the reorg caveat for the non-final suffix.
\end{itemize}

The bounded-availability argument already couples anchors to guaranteed
balances: a party's guaranteed minimum balance is ``the minimum balance
taken over all possible transaction histories that extend the finalized
chain'', accounting for republication of its transactions without consent;
``we know that shielded transactions cannot be reapplied after their
anchors have been invalidated. It may be desirable to further constrain
re-use in order to make guaranteed minimum balances easier to compute''
(tfl-book
\texttt{the-arguments-for-bounded-availability-and-finality-overrides.md}:50).
That sentence is the closest any source comes to proposing an anchor
restriction, and it proposes nothing concrete.

Why plain chains matter here: under Snap-and-Chat sanitization, expiry
preserved the idempotence argument (an expired transaction stays expired
on recheck), but anchors broke it --- ``it \emph{is} technically possible
for a duplicate transaction that was invalid because of a missing anchor
to \emph{become} valid in a subsequent block \ldots{} the same commitment
treestate can be produced by two unrelated transactions'' --- and some
finalized outputs could become unspendable outright, transactions in
$\mathrm{LOG}_{\mathrm{fin}}$ stranded on a fork no longer in any best
chain (tfl-book \texttt{notes-on-snap-and-chat.md}:67--70, 125 and
``Spending finalized outputs''). Crosslink~2 fixes both by construction,
via the Linearity and Last Final Snapshot rules, and avoiding sanitization
``would avoid breaking assumptions that may be made by existing Zcash node
implementations and other consumers of the Zcash block chain''
(\texttt{security-analysis.md}:263; \texttt{potential-changes.md}:348--363).
The wallet-visible consequence: under Crosslink~2, exactly the anchor and
expiry semantics documented in the \emph{Ironwood Guide} and \emph{Wallet
Guide} carry over --- on each of the two chains separately.

\paragraph{Classification.} The inherited anchor rule is \emph{designed
but unspecified} (construction text, both books). The anchor policy under
two ledgers --- both the consensus question and the wallet default --- is
an \emph{open problem}, and this volume flags the branch-cost analysis
above as its own synthesis.

\subsection{Confirmation policy under two ledgers}
\label{sec:xl-confpolicy}

The deployed baseline is ZIP 315 (Status: Draft), which the \emph{Wallet
Guide} develops in \S``The proposal'': a \textsf{ConfirmationsPolicy} with
trusted depth 3, untrusted depth 10, zero-conf shielding allowed by
default, an advisory floor of 1, and the $H-3$ anchor rule quoted above.
ZIP 315 predates Crosslink and never mentions it.

How that policy maps onto a finality status is unresolved everywhere. The
candidate shapes: the trusted/untrusted distinction collapses to
final/not-final; it becomes a three-state policy (final /
trusted-unfinalized / untrusted-unfinalized); or it stays unchanged with
finality surfaced only in UX. Shielded Labs commits to exposing the
status --- ``augmentation of existing RPC APIs for blocks and transactions
to include finality status'' (\texttt{book/src/design/deliverables.md},
GH~\#105) --- but no mapping into any confirmations policy exists in any
source. A related open point is whether displayed confirmation counts
should saturate at ``final'': the SL~book argues that when users' rollback
thresholds diverge during a stall, ``those users who require finality and
those users who do not \ldots{} diverge in their expectations and
behavior, which is a growing economic and governance risk''
(\texttt{book/src/design/security-properties.md}:17--19) --- finality
removes that fragmentation only if wallets adopt it uniformly.

The ECC design goals pull in two directions at once: wallet developers
``should not be required to make any changes through protocol transitions
as long as they rely solely on the lightwalletd protocol or a full node
API''; yet ``the Crosslink construction may require wallets to alter their
context of valid transactions differently from the current NU5/NU6
consensus protocol''; and there must be ``no security or safety
degradation due to wallet user behavior \ldots{} assuming users follow
their current behaviors unchanged'' (tfl-book
\texttt{src/design/goals.md}:9, 15--16, 28). The first goal is untestable
today: no lightwalletd or zaino wire protocol for finality status exists
in any source.

\paragraph{Classification.} \emph{Open problem} in every direction that
matters to a wallet author; the only \emph{specified}-adjacent artifact is
ZIP 315 itself, a Draft that predates the construction.

\subsection{Stalls: the divergence that decides wallet behaviour}
\label{sec:xl-stalls}

The finality gap at time $t$ is, roughly, the number of bc-blocks not yet
finalized; when roughly constant it corresponds to the time a transaction
takes to finalize after inclusion. A gap exceeding $L$ ``likely signals an
exceptional or emergency condition, in which it is undesirable to keep
accepting user transactions that spend funds'' (tfl-book
\texttt{construction.md}:403--405). The construction enforces this as a
block-validity rule: with $\mathsf{finality\text{-}depth}(H) :=
\mathsf{height}(H) - \mathsf{height}(\mathsf{snapshot}(\mathsf{LF}(H)))$,
every bc-block must satisfy $\mathsf{finality\text{-}depth}(H) \le L$ or
be a stalled block (\texttt{construction.md}:680--693).

\paragraph{Stalled Mode, as ECC designed it.} The proposed Zcash
instantiation of \textsf{is-stalled-block} is coinbase-only blocks ---
``Since funds cannot be spent in coinbase-only blocks, the vast majority
of attacks that we are worried about would not be exploitable'' ---
possibly with shielded coinbase disabled; mining pools cannot distribute
rewards during the stall (a rug-pull risk the book notes); and Stalled
Mode ``can only be entered by consensus of a larger validator set, or if
there is an availability failure of the finalization protocol'' (tfl-book
\texttt{the-arguments-for-bounded-availability-and-finality-overrides.md}:24,
103--111, 187).

\paragraph{The divergence.} The SL~book pastes the construction verbatim
but disclaims exactly this rule: ``The zebra-crosslink design in this book
diverges from this text by \emph{not} including `Stalled Mode' rules. This
simplification to the protocol rules follows the rationale that we are
already relying heavily on the stakers to guard safe operation of the
protocol by redelegating appropriately to prevent hazards like BFT
stalls'' (\texttt{book/src/design/cl2-construction.md}:9--12). No
stalled-mode or coinbase-only enforcement exists anywhere in the
implementation (grep over \texttt{zebra-crosslink} @ \texttt{6d02a1b8}
sources). The paste disclaimer adds that Daira-Emma Hopwood and Jack
Grigg ``have not reviewed and do not necessarily endorse its content or
design changes relative to Crosslink 2''; apart from the admonitions the
construction text is byte-identical to the tfl-book's.

\begin{remark}[Expiry as stall cleanup --- synthesis]
\label{rem:xl-stall-expiry}
Neither book discusses the interaction with transaction expiry; the
following is this volume's synthesis. Under the ECC branch a stall is a
hard availability stop with automatic cleanup: coinbase-only blocks keep
heights advancing, so every pending spend reaches its expiry height ---
target plus \texttt{DEFAULT\_TX\_EXPIRY\_DELTA} $= 40$ blocks (ZIP 203)
--- within about 40 blocks of the stall, and the wallet's locked notes
release under the unexpired-spend condition (\emph{Wallet Guide},
\S``Expiry (ZIP-203)'' and \S``Expiry and fund availability''). Under the
SL branch there is no such cleanup: spends keep confirming into an
unbounded finalization gap --- precisely the condition the ECC arguments
chapter deems unacceptable --- and a service's
\texttt{NotYetFinalized} exposure grows without bound. Whether a
Crosslink ZIP should couple the expiry delta to $L$ (for instance
$\Delta < L$, or $\Delta$ on the order of the finality latency) is
unaddressed in any source.
\end{remark}

\paragraph{Overrides.} Long rollbacks are argued to need a formalized,
governance-specified procedure, with the Bitcoin 2010 value-overflow
rollback (53 blocks, about 9 hours) and the Ethereum DAO fork as
precedents; bounded availability is complementary because it bounds the
period of spend transactions an intentional rollback could affect, and
the envisioned procedures preserve coinbase-only mining rewards across
the rollback (tfl-book \texttt{the-arguments\ldots{}md}:29, 148--198).
On the governance side, the SL draft ZIP on user-led emergency hardforks
``explicitly rejects the use of an emergency user-led hard fork to
mitigate harm arising from application-layer failures'' --- not wallet
bugs, bridge exploits, custodial failures, nor DAO-style reversals
(\texttt{book/src/design/userled-emergency-hardforks.md}:13) --- bounding
what the escape hatch may be used for.

\paragraph{Classification.} The Finality depth rule is \emph{designed but
unspecified} (tfl-book normative text); its omission is likewise
\emph{designed but unspecified} (SL~book admonition plus absent code).
The conflict itself is a finding: a ZIP must pick one branch, and each
branch changes what wallets and exchanges must handle during a stall ---
a hard availability stop with expiry churn, or indefinitely growing
unfinalized exposure. Remark~\ref{rem:xl-stall-expiry} is synthesis.

\subsection{Exchanges, bridges, and downstream volumes}
\label{sec:xl-services}

The stated service-facing goal is uniformity: ``CEXes and decentralized
services like DEXes/bridges are willing to rely on Zcash transaction
finality. E.g.\ Coinbase reduces its required confirmations to no longer
than Crosslink finality. All or most services rely on the same canonical
(protocol-provided) finality instead of enforcing their own additional
delays or conditions, so users get predictable transaction finality
across services'' (\texttt{book/src/design/scoping.md}:10, 30; GH~\#131,
\#128). The committed interface deliverables are finality status on block
and transaction RPCs (GH~\#105), proof-of-stake RPCs for finalizer
identities, staking weights, and delegation info (GH~\#104), and
``contingency mode, finality traversal'' RPCs (GH~\#172); external
validation partnerships prioritize wallet developers, finalizers, and
exchanges (\texttt{book/src/design/deliverables.md}).

The risk model a service would price is the SL five-component
allocation: miners ``cannot unilaterally violate finality, even if an
attacker controls $>1/2$ of the hashpower''; finalizers cannot
unilaterally compromise liveness or censorship-resistance even with
$>2/3$ of stake-weighted votes, and cannot execute rollbacks though they
can prevent them; more than $1/3$ of stake can stall; and a collusion of
$>1/2$ hashpower with $>2/3$ stake can violate finality, execute
rollbacks, and mount unavailability attacks
(\texttt{book/src/design/five-component-model.md}:24--46). The SL~book
separates Finality Progress from Liveness precisely because the
anticipated failure is a stall under continued block production, and part
of the design intent is to ``convert censorship attacks into stall
attacks'' (\texttt{book/src/design/security-properties.md}:17--19, 35).

What remains open on this surface: deposit semantics during a
finalization stall; whether \texttt{CantBeFinalized} can ever surface for
a transaction a service already credited; the contingency-mode and
finality-traversal RPC semantics (GH~\#172 is open); and the light-client
wire protocol of \S\ref{sec:xl-confpolicy}. Each is an \emph{open
problem}: the slogan is designed, the contract behind it is not.

\paragraph{Downstream volumes.} For coin-voting, Crosslink finality bears
on the snapshot pin: a round pins $(\textsf{snapshot\_height},
\textsf{snapshot\_blockhash})$ into its round identifier (\emph{Voting
Guide}, \S``The five-message round''), and requiring the pinned height to
be at or below the finalized height would make the pin irreversible under
Assured Finality; the orthogonal trust gap --- snapshot roots ``posted,
never verified'' (\emph{Voting Guide}, \S``The snapshot roots: posted,
never verified'') --- is untouched by finality and remains open. This
resolution analysis is this volume's synthesis.
%% TODO: confirm with the Voting Guide owner that the caveat is framed as
%% partially resolvable by Crosslink, not resolved.
Project Tachyon is formally independent of Crosslink; the \emph{Tachyon
Guide} (\S``Relation to Crosslink and the upgrade environment'') records
the independence as asserted rather than examined, and the NU8 vacuum it
documents is unchanged at \texttt{zcash/zips} @ \texttt{69610984}
(2026-07-09).

\paragraph{Classification.} The service goals and deliverables are
\emph{designed but unspecified} (scoping and deliverables documents with
open issue numbers); the five-component allocations are \emph{designed
but unspecified} security claims pending the volume's security analysis;
every concrete service-facing behaviour during a stall is an \emph{open
problem}.

\section{Security posture and open problems}\label{sec:xl-security}

The security posture of Crosslink~2 is best stated as an inventory of four
artifacts of unequal maturity. The design book proves its central safety
theorem twice over, from two independent assumptions, and then breaks off:
everything below a stated line of its analysis is marked ``Not updated for
Crosslink~2'', and its last proof sketch ends mid-sentence (tfl-book @
\texttt{fe6e1d6}, 2024-12-13, \texttt{src/design/crosslink/}%
\texttt{security-analysis.md}, lines 54 and 253--257). A formal model
machine-checks exactly one validity predicate of the BFT side and nothing
else (crosslink-spec @ \texttt{8edc3e9}, 2025-07-23). The implementation
departs from the analyzed design in two load-bearing ways that it documents
itself (zebra-crosslink @ \texttt{6d02a1b}, 2026-04-21). An independent
mechanism-design review finds the construction's safety ``structurally
strong once finality is reached'' and its finality progress ``economically
fragile in the baseline design'' (nikete.com/crosslink-zebra-audit,
completion announced 2026-02-05). We take the four in turn, keep the bin of
every claim visible --- \emph{specified} here can only mean verifiable
source at a pinned commit, since no Crosslink ZIP exists (zips @
\texttt{6961098}, 2026-07-09) --- and close with the open-problem ledger
and its revisit triggers. The rules and properties named below are
stated in \S\ref{sec:xl-bft}, \S\ref{sec:xl-bc}, and
\S\ref{sec:xl-ledgers}; the source pins are those of
Table~\ref{tab:xl-sources}.

\subsection{The safety theorems and their assumptions}
\label{sec:xl-safety-theorems}

The book's central safety goal is \emph{Assured Finality}
(Definition~\ref{def:xl-assured-finality}): in a given
execution, for all times $t, u$ and all nodes $i, j$ such that $i$ is
honest at time $t$ and $j$ is honest at time $u$, the finalized views
agree, $\mathsf{fin}_i^t \asymp_{\mathrm{bc}} \mathsf{fin}_j^u$ --- either
is a prefix of the other (tfl-book @ \texttt{fe6e1d6},
\texttt{construction.md}, lines 504--506). The name is chosen against
``absolute finality'': the property binds \emph{the protocol}, and the
book states explicitly that out-of-band reversal by community fork remains
possible (\texttt{terminology.md}, lines 11--13). The implementation's
gloss, ``crosslink finality'', strengthens the wording: accountable via
explicit on-chain voting records, irreversible within protocol scope,
objectively verifiable from local block history, globally consistent
(zebra-crosslink @ \texttt{6d02a1b},
\texttt{book/src/design/terminology.md}, lines 7--13). Two theorems each
imply Assured Finality from a safety property of one subprotocol alone.

\begin{theorem}[Prefix Agreement implies Assured Finality]
\label{thm:xl-safety-prefix}
In an execution of Crosslink~2 in which subprotocol $\Pi_{\mathrm{bc}}$
has Prefix Agreement (Definition~\ref{def:xl-prefix-agreement}) at
confirmation depth $\sigma$ --- all honest nodes'
bc-best-chains, truncated by $\sigma$ blocks, pairwise agree across all
times --- that execution has Assured Finality (tfl-book @
\texttt{fe6e1d6}, \texttt{security-analysis.md}, lines 72--77;
Prefix Agreement at \texttt{construction.md}, lines 96--99 and 337--340).
\end{theorem}

The written proof is complete and short: the Local fin-depth lemma places
each $\mathsf{fin}_i^t$ inside some earlier $\sigma$-truncated best chain
of node $i$ (\texttt{construction.md}, lines 491--494), Prefix Agreement
relates any two such chains, and the Linear prefix lemma --- two prefixes
of a common chain agree --- closes the argument (\texttt{construction.md},
lines 77--82). The book contrasts this with the corresponding result for
Snap-and-Chat: ``[NTT2020, Theorem~2] is a much more complicated proof
that takes nearly 3 pages, not counting the reliance on results from
[PS2017]'' (\texttt{security-analysis.md}, line 162). The comparison is
fair in kind but not in coverage: NTT2020's Theorem~2 is a probabilistic
security statement about the longest-chain protocol itself (failure
probability $e^{-\Omega(\sqrt{\sigma})}$ for block-production rate
$p < (n-2f)/(2 \Delta n (n-f))$; arXiv:2009.04987v3, Appendix~C), whereas
Theorem~\ref{thm:xl-safety-prefix} \emph{assumes} the corresponding
property of $\Pi_{\mathrm{bc}}$ as a black box. That division of labor is
this series' scope rule in protocol form: Crosslink consumes
$k$-deep-confirmation guarantees of the PoW layer; it does not re-derive
them.

\begin{theorem}[Final Agreement implies Assured Finality]
\label{thm:xl-safety-final}
In an execution of Crosslink~2 in which subprotocol $\Pi_{\mathrm{bft}}$
has Final Agreement (Definition~\ref{def:xl-final-agreement}) --- for
all bft-valid blocks $C, C'$ in honest view at
any times, $\textsf{bft-last-final}(C) \asymp_{\mathrm{bft}}
\textsf{bft-last-final}(C')$ --- that execution has Assured Finality
(tfl-book @ \texttt{fe6e1d6}, \texttt{security-analysis.md}, lines 86--91;
Final Agreement at \texttt{construction.md}, lines 203--206).
\end{theorem}

\begin{remark}[The written proof of Theorem~\ref{thm:xl-safety-final} is
defective]\label{rem:xl-final-agreement-proof}
The proof text pasted under the theorem in the book is a verbatim
duplicate of the proof of Theorem~\ref{thm:xl-safety-prefix}: it invokes
Prefix Agreement of $\Pi_{\mathrm{bc}}$, not Final Agreement of
$\Pi_{\mathrm{bft}}$ (\texttt{security-analysis.md}, lines 86--91). The
intended route presumably runs through the Linearity rule, as sketched at
\texttt{construction.md}, line 592, but no correct written proof exists in
the sources we have read. Bin: \emph{open problem} --- the theorem
statement is the design's load-bearing claim for the case of a subverted
PoW layer, and it currently has no proof.
\end{remark}

The pair of theorems, where both hold, gives the design's headline
property: safety of the finalized ledger requires only \emph{one} of
$\{\Pi_{\mathrm{bc}}$ safety, $\Pi_{\mathrm{bft}}$ safety$\}$ --- the
stronger of the two. The book carries the same two-route structure down to
the ledger level: in an execution where $\Pi_{\mathrm{bc}}$ has Prefix
Agreement at depth $\sigma$ \emph{or} $\Pi_{\mathrm{bft}}$ has Final
Agreement, any two honest nodes' finalized ledgers
$\mathsf{LOG}_{\mathrm{fin}}$ agree at all times
(\texttt{security-analysis.md}, lines 117--176). That section, however,
sits below the ``Not updated for Crosslink~2'' line (line 54), so its
statements are for Crosslink~1 and their transfer is unverified.

Two further results are proved, both assuming the \emph{one-third bound
on non-honest voting} (Definition~\ref{def:xl-one-third}): strictly
fewer than one third of each epoch's voting
units are ever cast non-honestly, counted over the entire execution
(\texttt{construction.md}, lines 131--141). First, uniqueness: any
bft-valid block for a given epoch in honest view commits to the same
proposal, by quorum intersection of two two-thirds vote sets, adapted from
[CS2020, Lemma~1] (\texttt{security-analysis.md}, lines 208--213). Second,
snapshot validity: with Prefix Consistency of $\Pi_{\mathrm{bc}}$ at depth
$\sigma$ added, every bc-chain snapshot contributing to
$\mathsf{LOG}_{\mathrm{fin}}$ eventually sits in every honest node's
bc-best-chain (\texttt{security-analysis.md}, lines 232--251).

\paragraph{Assumption transfer.} The theorems assume properties of the
\emph{modified} subprotocols $\Pi_{\mathrm{bc}}$ and $\Pi_{\mathrm{bft}}$,
while the underlying literature supports the \emph{original} protocols
$\Pi_{\mathrm{origbc}}$ and $\Pi_{\mathrm{origbft}}$. In the BFT
direction the transfer is argued: $\Pi_{\mathrm{bft}}$ only adds
structural components and tightens validity rules, so every
$\Pi_{\mathrm{bft}}$ execution maps to a $\Pi_{\mathrm{origbft}}$
execution with the additions erased. In the PoW direction it is
acknowledged as unproven: the modifications ``seem unlikely to have any
significant effect on this property'', followed by ``TODO: that is a
handwave; we should be able to do better, as we do for
$\Pi_{\mathrm{bft}}$ below'' (\texttt{security-analysis.md}, lines 119 and
170). Bin: \emph{open problem}.

\paragraph{The residual attack the design accepts.} If the network is
partitioned \emph{and} $\Pi_{\mathrm{bft}}$ is subverted by more than one
third of validators, the adversary can vote with an additional one third
on each side of the fork, so that each side's last-final bft-block is
consistent with its own fork. The book calls this unavoidable:
``$\Pi_{\mathrm{bc}}$ doesn't include the mechanisms needed to maintain
finality under partition; it takes a different position on the CAP
trilemma.'' Snap-and-Chat, by contrast, loses safety under BFT subversion
even \emph{without} a partition; Crosslink's stated aim is to limit safety
loss to attacks ``like'' the unavoidable one
(\texttt{security-analysis.md}, lines 109--115). The corresponding goal
exceeds NTT2020's own argument: the paper's Figure-9c reasoning --- one
honest vote is needed to finalize a snapshot, hence only valid snapshots
finalize --- applies only for adversarial fractions $n/3 < f < n/2$,
whereas the book wants the conclusion to survive an adversary holding
100\% of validators but under 50\% of $\Pi_{\mathrm{lc}}$ hash rate
(\texttt{notes-on-snap-and-chat.md}, lines 205--225).

\paragraph{Classification.} Every statement above is \emph{designed but
unspecified}: the proofs live in a design book at a pinned commit, not in
a ZIP, a peer-reviewed paper, or machine-checked form. Within that bin,
Theorem~\ref{thm:xl-safety-prefix} is complete on its stated assumptions;
Theorem~\ref{thm:xl-safety-final} lacks a correct written proof
(Remark~\ref{rem:xl-final-agreement-proof}); the ledger-level and
snapshot-validity theorems predate Crosslink~2; and the
$\Pi_{\mathrm{origbc}} \to \Pi_{\mathrm{bc}}$ assumption transfer is an
acknowledged gap.
\subsection{Liveness and the unfinished analysis}\label{sec:xl-liveness}

The liveness side of the analysis is prose, and the book says so. Two
conclusions are argued but not proven: subprotocol $\Pi_{\mathrm{bc}}$
remains live under the same conditions as $\Pi_{\mathrm{origbc}}$, because
the added consensus rules never obstruct producing a coinbase-only stalled
block on the current best chain; and $\Pi_{\mathrm{bft}}$ remains live
under the same conditions as $\Pi_{\mathrm{origbft}}$, because the
Linearity and Tail Confirmation rules always permit an honest proposer to
re-propose its parent's \texttt{headers\_bc} (tfl-book @ \texttt{fe6e1d6},
\texttt{security-analysis.md}, lines 25--42). Two open TODOs sit inside
the argument: that a new higher-scoring snapshot eventually becomes final
(``make that more precise''), and that Stalled Mode triggers only when the
BFT liveness assumptions are violated --- ``more detailed argument needed,
especially for the dependence on $L$'', where the dependence involves
$\sigma$, network delay, honest hash rate, and honest proposers and voting
units (lines 46 and 50). The quantitative relationship among these
parameters was never worked out anywhere in the sources we have read; the
book's design guidance is only that $L$ be ``significantly greater than
$\sigma$'' and ``in practice \ldots at least $2\sigma$''
(\texttt{construction.md}, lines 391 and 405). Bin: \emph{open problem}.

Three further admissions bound what the analysis currently delivers.
First, the rollback exposure of the available ledger: the book proves that
$\mathsf{LOG}_{\mathrm{ba}}$ never rolls back past the agreed
$\mathsf{LOG}_{\mathrm{fin}}$, but notes the result is weaker than
desired, since rollbacks of $\mathsf{LOG}_{\mathrm{ba}}$ down to the
finalization point ``may be of unbounded length'' --- bounding them by $L$
is asserted loosely and ``we have also not proven that so far''
(\texttt{security-analysis.md}, lines 180--190). Second, rule necessity:
the claim that every Crosslink~2 rule is individually necessary is
qualified --- ``some rules, e.g.\ the Linearity rule, only contribute
heuristically to security in the analysis so far'' (lines 275--281).
Third, validation cost: honestly voting on a proposal may recursively
require validating the referenced bft-chain, each bft-block's
\texttt{headers\_bc} chain, and so on --- the book's own gloss is ``Wait
what, how much validation is that?'' --- with denial-of-service bounded
only by validation ordering (check the supermajority and PoW first) and by
the Linearity rule limiting live bc-chains per bft-chain to one
(\texttt{construction.md}, lines 641--652).

\begin{remark}[Fork-choice independence]\label{rem:xl-fork-choice}
Crosslink~2 deliberately does not modify the fork-choice rule of
$\Pi_{\mathrm{bc}}$, and the book identifies this as the reason the
bidirectional-dependence liveness failure of Casper-FFG-style designs
(NTT2020, Appendix~E, ``Bouncing Attack on Casper FFG'') does not arise:
an honest bc-block-producer ``must not use information from the BFT
protocol, other than the specified consensus rules, to decide which
bc-valid-chain to follow'' (\texttt{security-analysis.md}, lines 15--23;
\texttt{construction.md}, lines 715--717). The same modularity is why
$\mathsf{snapshot}(\mathsf{LF}(H)) \preceq_{\mathrm{bc}} H$ was rejected
as a fork-choice constraint: under a potentially subverted
$\Pi_{\mathrm{bft}}$, ``we cannot say anything useful about
$\mathsf{snapshot}(B)$'', so the constraint would break the modular
safety and liveness arguments. Honest validators may instead simply refuse
to vote for proposals embedding long rollbacks --- the voting rule is
``only if'', not ``iff'' --- at an accepted cost in BFT liveness (``and
that's okay''; \texttt{questions.md}, lines 44--58).
\end{remark}

\begin{remark}[The model is not NTT2020's model]\label{rem:xl-model}
The book's communication model deliberately avoids global-stabilization-%
time partial synchrony as a base assumption --- messages may be arbitrarily
delayed or dropped --- with a detailed critique of applying the DLS1988
GST formalization to liveness of non-terminating blockchain protocols,
including in NTT2020 and the Streamlet paper; safety properties are stated
as unconditional properties of executions, with probabilistic reasoning
factored out separately (\texttt{construction.md}, lines 13--41 and
342--366). A consequence worth recording: NTT2020 defines a
$(\beta_1, \beta_2)$-secure ebb-and-flow protocol by exactly such
GST/GAT-parameterized environments (Definition~3) and proves Snap-and-Chat
optimally $(1/3, 1/2)$-secure (Theorem~1; arXiv:2009.04987v3), and no
theorem of that shape appears for Crosslink~2 in any source we have read.
The two analyses are therefore not directly comparable, by the book's own
choice of model. Bin: \emph{open problem} --- no
$(\beta_1, \beta_2)$-style security statement for Crosslink~2 exists in
any pinned source.
\end{remark}

The book's status chapter has not moved since 2024: ``This is an early and
incomplete protocol design proposal. It has not been well vetted for
feasibility and safety'', with ``Security and safety analyses'' still on
the missing-components list beside subprotocol selection, issuance, a
transition plan, and ZIPs (\texttt{status-and-next-steps.md}, lines 3 and
19--27). The book has been untouched since 2024-12-13 while the
implementation advanced through 2026-04; which document is the
authoritative security argument for the protocol actually being built is
itself an open problem (\S\ref{sec:xl-open}).

\subsection{Bounded availability, Stalled Mode, and overrides}
\label{sec:xl-bounded}

Crosslink's finality is \emph{conditional} by construction. An honest
validator votes for a proposal only if it is bft-proposal-valid and its
snapshot is $\sigma$-confirmed in the validator's own bc-best-chain
(\texttt{construction.md}, lines 654--657), so a coalition holding more
than one third of an epoch's voting units can stall finality
indefinitely: finality arrives only while a two-thirds quorum remains
live and honest enough. The design's answer is to make the stall visible
and bounded in its consequences rather than impossible: once the gap
between the tip and the finalized snapshot exceeds $L$ blocks, the
Finality depth rule forces every new block to be a coinbase-only
``stalled'' block
(\texttt{the-arguments-for-bounded-availability-and-finality-%
overrides.md}, lines 180--188), and Stalled Mode ``can only be entered by
consensus of a larger validator set, or if there is an availability
failure of the finalization protocol'' (tfl-book @ \texttt{fe6e1d6},
\texttt{the-arguments-\ldots.md},
line 187). The same machinery is a censorship defense: part of the design
intent is ``to `convert censorship attacks into stall attacks'\,'' --- an
attacker cannot censor targeted transactions without stalling the whole
service, and a stall is immediately and consensually visible while
censorship tends to be visible only to its victims
(\texttt{book/src/design/security-properties.md}, line 35).

The argument for bounding availability at all is the CAP position:
finality plus dynamic availability plus the possibility of partition
implies an unbounded finalization gap (made precise by [LR2020]), and
allowing spends into an unbounded gap risks a crisis of confidence; the
book cites Ethereum's May 2023 double finality stall ($\approx$25 and
$\approx$64 minutes, a Prysm resource-exhaustion bug) as evidence that
stalls happen, that short ones self-heal, and that long ones need human
intervention (tfl-book @ \texttt{fe6e1d6},
\texttt{the-arguments-\ldots.md}, lines 31--41 and 135--139). The reason
to prefer coinbase-only blocks over halting outright is the
tail-thrashing attack: during a stall of a halted chain, an adversary with
10\% of hash power builds a 10-block fork in 100 block times because the
honest chain is not advancing, letting the $k$-block tail thrash between
chains --- and hash rate is expected to \emph{fall} during a stall as
miners anticipate rollback (lines 200--229). Two caveats attach. The
constraint binds only while a producer's view of the finalization point is
not advancing: ``an adversary that has subverted the BFT protocol in a way
that does not keep the finalization point from advancing, can always
avoid being constrained by Stalled Mode'' (\texttt{construction.md}, lines
409--411). And finality can be overridden deliberately: the book argues
override must remain possible for exploited flaws (Bitcoin's 2010
value-overflow rollback of 53 blocks over roughly nine hours; Ethereum's
DAO fork), should be \emph{formalized} under a well-defined governance
process, and composes with bounded availability --- validators asked to
stop cause a stall, the chain enters Stalled Mode after $L$ blocks, only a
bounded window of user spends is exposed to the rollback, and miners'
coinbase-only rewards would be preserved across it
(\texttt{the-arguments-\ldots.md}, lines 148--198).

\begin{remark}[Long-range attacks and key hygiene]\label{rem:xl-longrange}
The one-third bound counts ballots cast at \emph{any} time in the
execution, so a compromised old validator key used to vote for a
long-finished epoch violates it retroactively. The book's consequence is
operational: ``validator keys of honest nodes must remain secret
indefinitely. Whenever a key is rotated, the old key must be securely
deleted.'' The recommended mitigation is a checkpoint bft-block baked into
each released node version, with local finality exposed to clients only
once synced past it; improvements are tracked as tfl-book issue \#140
(\texttt{construction.md}, lines 149--154 and 557--562). Bin:
\emph{designed but unspecified}, with the issue open.
\end{remark}

\paragraph{Classification.} The Finality depth rule and Stalled Mode are
\emph{designed but unspecified} consensus rules of the book's Crosslink~2;
the implementation has dropped them (\S\ref{sec:xl-impl-divergences}). The
governance process that a formalized finality override presupposes does
not exist in any pinned source; the accountability machinery that would
let a bc-transaction carry objective evidence of a BFT safety violation
(two final bft-blocks with the same parent, a ``bft-final-vee'') is a
change marked [Recommended] but not adopted into Crosslink~2, and remains
unimplemented (\texttt{potential-changes.md}, lines 12--24;
\texttt{questions.md}, lines 41--42). Both are \emph{open problems}.

\subsection{Machine-checked coverage: the Quint model}
\label{sec:xl-quint-coverage}

The only machine-checked artifact in the corpus is the Informal Systems
Quint specification, and its own README bounds its scope: it is a
self-described proof of concept that encodes ``a tiny part of the system,
namely the validity check, of a crosslink proposal within the BFT
consensus mechanism'', with most of the required information ``inferred
\ldots from an insightful Youtube talk'', a warning that ``the encoded
logic might not be up-to-date with the current protocol design ideas'',
and a closing question to the community ``whether we should move forward
with applying for a grant'' (crosslink-spec @ \texttt{8edc3e9},
\texttt{README.md}, lines 27 and 93).

What the model checks is nonetheless crisp. The validity predicate is
\[
\begin{aligned}
\mathsf{isValid}(\mathit{bft},\mathit{pow},p) = {}&
  \mathit{pow}.\mathsf{contains}(p.\mathsf{block})
  \;\wedge\; p.\mathsf{bftContext} = \mathit{bft}.\mathsf{last}()\\
&\wedge\; \mathsf{sigmaConfirmed}(p.\mathsf{block},
    p.\mathsf{confirmation\_tail}, \sigma)\\
&\wedge\; \mathsf{linear}(p.\mathsf{block}, p.\mathsf{bftContext},
    \mathit{pow}),
\end{aligned}
\]
over a single node's \texttt{pow\_store} and \texttt{bft\_store}
(\texttt{validity.md}, lines 131--176). Against it the spec generates a
finalization witness (reachability of three BFT blocks) and checks the
invariant \texttt{prefix\_inv} --- ``the PoW history defined by BFT
finalization is a path in the PoW block store'' --- by random simulation
($10{,}000$ traces of $20$ steps, roughly $132$ seconds, no violation
found) and by TLC model checking after transpilation to TLA\textsuperscript{+},
on stores of at most $8$ blocks at confirmation depth $1$, in about three
minutes (\texttt{README.md}, lines 37--89; \texttt{validity.md}, lines
405--410 and 463--482; \texttt{src/validity.cfg}). The model contains no
network, no adversary, no voting, no epochs, no signatures, none of the
bc-side rules (Extension, Last Final Snapshot, Finality depth), no Stalled
Mode, and no liveness property.

\begin{remark}[A divergence the spec itself flags]\label{rem:xl-quint-div}
The Quint operator $\mathsf{linear}(b, \mathit{ctx}, \mathit{store})$
conjoins descent with
$\neg(\mathit{ctx}.\mathsf{pow\_hash} = b.\mathsf{hash})$: it
\emph{rejects} a proposal whose snapshot equals the parent's snapshot. The
book's Linearity rule uses $\preceq_{\mathrm{bc}}$, which includes
equality --- and equality is required for $\Pi_{\mathrm{bft}}$ liveness,
since an honest proposer repeats its parent's \texttt{headers\_bc} when it
cannot extend. The spec's own comment reads ``this actually seems allowed.
TODO: perhaps remove'' (crosslink-spec @ \texttt{8edc3e9},
\texttt{validity.md}, lines 171--176; tfl-book @ \texttt{fe6e1d6},
\texttt{construction.md}, lines 71, 573, 611--616, and 627).
\end{remark}

\paragraph{Classification.} The Quint model and its check results are
\emph{specified} in the only sense available --- runnable, pinned source
--- but of deliberately narrow scope. Machine-checking has not progressed
since: the spec's last commit is 2025-07-23, the companion
\texttt{quint-crosslink} repository contains no commits at all, and a web
search on 2026-07-15 found no public evidence of a funded extension.
Coverage of the bc-side rules, of any multi-node or adversarial property,
and of liveness is an \emph{open problem}.
\subsection{The implementation diverges from the analyzed design}
\label{sec:xl-impl-divergences}

The zebra-crosslink book embeds verbatim copies of the tfl-book's
construction and security-analysis chapters, made ``for the purposes of
precisely committing to a (potentially incomplete, confusing, or
inconsistent) precise set of book contents for a security design
review''; diffing them against tfl-book @ \texttt{fe6e1d6} shows the only
changes are three warning admonitions, and the copies state that
Daira-Emma Hopwood and Jack Grigg ``have not reviewed and do not
necessarily endorse'' them (zebra-crosslink @ \texttt{6d02a1b},
\texttt{book/src/design/cl2-construction.md} and
\texttt{design/analysis/cl2-security-analysis.md}). The project's own
design chapter is franker still: ``These extensions and changes from the
Crosslink 2 construction may violate some of the assumptions in the
ssecurity [sic] proofs. This still needs to be reviewed after this design
is more complete'', and ``We are using a custom-fit in-house BFT protocol
for prototyping. It's [sic] own security properties need to be more
thoroughly verified'' (\texttt{book/src/design.md}, lines 11--12). Two of
the
admonitions record load-bearing divergences.

\paragraph{Divergence 1: Stalled Mode is dropped.} ``The zebra-crosslink
design in this book diverges from this text by \emph{not} including
`Stalled Mode' rules. This simplification to the protocol rules follows
the rationale that we are already relying heavily on the stakers to guard
safe operation of the protocol by redelegating appropriately to prevent
hazards like BFT stalls'' (\texttt{book/src/design/cl2-construction.md},
lines 10--12). The implementation as designed therefore has an unbounded
finalization gap with spends continuing --- precisely the condition the
tfl-book's bounded-availability chapter argues is unacceptable, and the
regime whose tail-thrashing analysis motivated Stalled Mode in the first
place (tfl-book @ \texttt{fe6e1d6}, \texttt{the-arguments-\ldots.md},
lines 17--27; \S\ref{sec:xl-bounded}). No pinned source reconciles the two
positions.

\paragraph{Divergence 2: the outer signature is dropped.} ``The
zebra-crosslink design does not use the `outer signature' referred to in
this chapter\ldots\ This means a direct hash or copy of the most recent
BFT finality certificate is insufficient evidence to ensure that specific
bitwise copy is canonical or will become canonical in the chain''
(\texttt{book/src/design/cl2-construction.md}, lines 14--16). The book's
model has the proposer re-sign $(P, \mathsf{proof}_P)$ with a strongly
unforgeable signature (tfl-book @ \texttt{fe6e1d6},
\texttt{construction.md}, lines 156--159); the implementation's
replacement mechanism and its security argument are undocumented in the
sources we have read.
%% TODO: revisit when zebra-crosslink documents the outer-signature
%% replacement; no security argument exists at 6d02a1b.

\paragraph{Which BFT protocol, exactly.} The book's analysis instantiates
$\Pi_{\mathrm{origbft}}$ as adapted Streamlet: notarization at two-thirds
of voting units, finality at the second of three adjacent notarized blocks
with consecutive epochs, liveness after five consecutive honest-leader
epochs (tfl-book @ \texttt{fe6e1d6}, \texttt{construction.md}, lines
218--241 and 614--618). The implementation at HEAD instead uses
``tenderlink'' 0.1.0, an in-house Tendermint-style engine pinned from
\texttt{git.sr.ht/\textasciitilde azmr/stuff} at revision
\texttt{0e885095\ldots}; an earlier Malachite (Informal Systems)
integration survives only as dead code --- \texttt{src/malctx.rs} still
imports \texttt{malachitebft\_*} crates but is not declared as a module,
and \texttt{Cargo.lock} contains no malachitebft packages
(zebra-crosslink @ \texttt{6d02a1b}, \texttt{zebra-crosslink/Cargo.toml};
\texttt{Cargo.lock}, lines 5445--5459; \texttt{src/lib.rs}). No published
analysis shows tenderlink satisfies the properties the Crosslink safety
arguments assume of $\Pi_{\mathrm{origbft}}$ --- Final Agreement,
objective block validity, sound two-thirds notarization. Bin: \emph{open
problem}, in the project's own words quoted above.

\paragraph{Parameters, and the black box today.} The prototype ships
$\sigma = 3$ and $L = 7$ under
\texttt{PROTOTYPE\_PARAMETERS}, annotated ``No verification has been done
on the security or performance of these parameters''
(\texttt{zebra-crosslink/src/chain.rs}, lines 219--222); production values
are unspecified everywhere. For comparison, the $\Pi_{\mathrm{bc}}$
interface that Crosslink would consume today is the fork's hard rollback
limit $\texttt{MAX\_BLOCK\_REORG\_HEIGHT} = 99$ --- the non-finalized
state never reorganizes deeper, and zebra-crosslink reuses the constant
as its pre-BFT fallback finality depth (zebra-crosslink @
\texttt{6d02a1b}, \texttt{zebra-state/src/constants.rs}, line 31;
\texttt{zebra-crosslink/src/lib.rs}, lines 260--292); upstream Zebra has
since raised the same constant to $1000$ (zebra @ \texttt{9f3166b},
\texttt{zebra-chain/src/parameters/constants.rs}, line 30), and the
\emph{Ironwood Guide}'s 100-block ``reorg horizon'' for anchor selection
is flagged as a cross-volume erratum in \S\ref{sec:xl-node-interface}.
We cite the bound as interface, per this series' scope rule; its
derivation belongs to the PoW layer.

\paragraph{Implementation evidence of reorg fragility.} Deriving the
pre-Crosslink final block from Zebra's block locator proved subtle: the
node wrote three alternative implementations of
\texttt{tfl\_reorg\_final\_}\allowbreak\texttt{block\_height\_hash()} and
cross-asserted them, and the retained comment reads ``NOTE: possible race
condition: only for testing / Sam: YES! Indeed there were race
conditions.'' The node also keeps a stateful
\texttt{latest\_final\_block} that shadows the state-service view,
mirroring the book's node-local $\mathsf{fin}_i$ state machine, whose
``record finalization safety hazard'' branch arises exactly when the
global assumptions fail (zebra-crosslink @ \texttt{6d02a1b},
\texttt{zebra-crosslink/src/lib.rs}, lines 235--332; tfl-book @
\texttt{fe6e1d6}, \texttt{construction.md}, lines 471--488). The
interleaving of BFT finality with PoW reorg handling is where the
implementation has already met race conditions; it deserves adversarial
review once the design settles.

\subsection{Economic security: the mechanism-design review}
\label{sec:xl-econ}

The zebra-crosslink design decomposes security across five components:
wallets (privacy and censorship defense), miners (liveness and
censorship resistance; unable to violate finality even with more than
half the hash power), finalizers (more than one third of stake-weighted
votes can stall but cannot unilaterally roll back), stakers (holding
finalizers accountable solely by \emph{redelegating} --- there is no
slashing), and users (last-resort accountability via a user-led emergency
hard fork, with STEEM/HIVE cited as precedent). Its own attack scenarios:
a collusion of more than half the hash power \emph{and} more than two
thirds of stake can violate finality and execute rollback or
unavailability attacks; and because staking and redelegation transactions
themselves need censorship resistance, an attacker with more than half
the hash power can censor redelegations and thereby disable the sole
staker-accountability mechanism (zebra-crosslink @ \texttt{6d02a1b},
\texttt{book/src/design/five-component-model.md}, lines 5--52). The first
deployment concentrates the finalizer set deliberately: the active set is
the top $K = 100$ finalizers by stake weight, and support for
finalizers run by arbitrary or casual users is explicitly out of scope
(\texttt{book/src/design/scoping.md}, lines 53--57;
\texttt{nutshell.md}, line 104).

The independent review of this mechanism is the nikete (Nicolas Della
Penna) audit: pre-registered 2025-12-11, completion announced by Shielded
Labs on 2026-02-05, report at \texttt{nikete.com/crosslink-zebra-audit}.
Its verdict: ``safety (irreversibility) is structurally strong once
finality is reached'', but ``finality progress (liveness) is economically
fragile in the baseline design and requires explicit incentive
constraints''. Three failure modes: \emph{hostage holdout} --- pending
rewards accumulate until a finality-updating block, so a blocking
coalition rationally delays finality to grow its future payout;
\emph{zombie set} --- the finality gadget becomes permanently inoperable
through validator attrition during a freeze; and \emph{phantom security}
--- if stake can exit economically while a frozen validator set retains
voting weight, the effective cost to corrupt the frozen set decays over
time. Its recommendations: signer-conditioned commission, first-inclusion
miner bounties tied to finalized height, finality-gap telemetry,
anti-jackpot constraints on pending rewards during stalls, and
governance-level liveness resets (nikete.com/crosslink-zebra-audit,
accessed 2026-07-15; pre-registration at
\texttt{book/src/design/analysis/history/2025-12-11-mech-design.md} and
\texttt{\ldots/2025-12-11-mech-design/proposal.md} @ \texttt{6d02a1b8}).
The
audit's ground rules bind its scope: no protocol-level slashing in V1 is
a design constraint, the assumed economics were a 40/40 issuance split
and 5-day epochs with 5--10-day unbonding, and privacy impacts on stakers
or finalizers were out of scope
(\texttt{\ldots/2025-12-11-mech-design/proposal.md}, lines 33, 58, and
64--66). Whether any recommendation has been incorporated is not visible
in the repositories read.
%% TODO: recheck zebra-crosslink after the next milestone for adoption of
%% the audit's recommendations (none visible at 6d02a1b).

The reward mechanics the audit targets are design-stage: pending rewards
$R$ distribute at a finality-updating block as
$S_{\mathrm{finalizer}} = c \cdot W_{\mathrm{finalizer}} /
W_{\mathrm{total}}$ and
$S_{\mathrm{staker}} = (1 - c) \cdot W_{\mathrm{staker}} /
W_{\mathrm{total}}$ with commission rate $c = 0.1$, totals possibly
summing below $S_{\mathrm{total}}$ because stake assigned outside the
active set earns nothing (\texttt{book/src/design/nutshell.md}, lines
19--37). The hostage-holdout finding is a direct consequence of that
``until a finality-updating block'' accumulation. One privacy cost is
already fixed at design stage: staking positions are ledger state with
on-chain amount, target-finalizer key, and per-position redelegation key
--- ``staked amounts are revealed so the community can verify the total
ZEC staked and its distribution across finalizers'' --- mitigated only by
quantizing stakes to powers of ten ZEC; the Crosslink Staking Orchard
Restriction confines staking transactions to Orchard, fees, and burns,
and Penumbra-style shielded delegation is deferred to ``continued
exploration in future versions'' (\texttt{nutshell.md}, lines 39--104;
shieldedlabs.net/crosslink-updated-staking-design/, 2025-10-27, accessed
2026-07-15).

\paragraph{Classification.} The five-component model, the reward split,
and the staking design are \emph{designed but unspecified}; the audit
findings are \emph{specified} as a pinned report but normative for
nothing. Whether the no-slashing, redelegation-only accountability model
survives the phantom-security finding is an \emph{open problem} --- the
audit proposal itself lists ``unbounded divergence between PoW progress
and a stalled BFT subprotocol'' among the failure modes to examine, which
is exactly the regime the implementation entered by dropping Stalled Mode
(\S\ref{sec:xl-impl-divergences}).

\subsection{Standing, siblings, and the open-problem ledger}
\label{sec:xl-open}

\paragraph{Normative standing.} No Crosslink ZIP exists: the zips
repository at \texttt{6961098} (2026-07-09) contains Crosslink only as a
citation footnote in draft ZIP~218 (25-second block target spacing) and as
the source of the Assured Finality definition cited by
\texttt{draft-ecc-}\allowbreak\texttt{onchain-accountable-voting}; the
zebra-crosslink book
says the ZIP draft exists as a Google Doc, to be refined ``as the
prototype approaches maturity'' (\texttt{book/src/design.md}, line 7).
NU8 scoping is an open governance discussion (``NU8 Decision Framework'',
forum thread opened 2026-02-11) in which Crosslink competes with the
25-second-blocks proposal and with Tachyon; the \emph{Tachyon Guide}
records the same environment from the other side (\S``Relation to
Crosslink and the upgrade environment''). Shielded Labs' roadmap shows
Milestones 1--5 complete (2025-03-21 through 2026-04-15) and a
productionization phase running Q2~2026 -- Q2~2027 --- testnet
deployment, design iteration, potential integration
(shieldedlabs.net/roadmap/, accessed 2026-07-15). No confirmed mainnet
activation date exists.

\paragraph{What Crosslink would change for the siblings.} Today the stack
expresses transaction confirmation entirely as depth policy on the
black-box PoW chain: the \emph{Wallet Guide} documents the ZIP-315
\texttt{ConfirmationsPolicy} with trusted depth $3$ and untrusted depth
$10$ (\S``Transaction construction'', Definition ``Confirmations
policy''), and the \emph{Ironwood Guide} records the 100-block reorg
horizon governing anchor robustness (\S``Proof of ownership of
\emph{some} valid note: the commitment tree''). Crosslink would add a
protocol-level notion beneath both: finality status is the sole
``height-eventual'' ledger state it introduces --- computable for height
$H$ only from a later attesting block, with the guarantee that if the
attesting block rolls back, ``an alternative block will eventually attest
to the finality status of the same candidate block'' (zebra-crosslink @
\texttt{6d02a1b}, \texttt{book/src/design/nutshell.md}, lines 116--124).
How wallet policy would consume assured finality is unspecified in every
source read. Separately, to prevent a terminology collision: the
\emph{Voting Guide}'s ``finalization'' is an application-level tally
event with its own caveat --- a round can finalize with empty results
(\S``Trust model and verdict'', ``Tally liveness admits finalization
without results'') --- and has no relation to consensus finality.

\paragraph{The open-problem ledger.} We close with the ledger, each entry
carrying its bin and, where one exists, its revisit trigger.

\begin{enumerate}
\item \emph{The authoritative security argument.} The analyzed design
  (tfl-book, frozen 2024-12-13, analysis incomplete below
  \texttt{security-analysis.md}, line 54) and the built design
  (zebra-crosslink, advancing through 2026-04, Stalled Mode and outer
  signature dropped) have diverged with no document proving the latter.
  Trigger: completion or supersession of the tfl-book analysis.
\item \emph{The Final Agreement proof defect}
  (Remark~\ref{rem:xl-final-agreement-proof}) and the unproven
  $\Pi_{\mathrm{origbc}} \to \Pi_{\mathrm{bc}}$ assumption transfer
  (\S\ref{sec:xl-safety-theorems}).
\item \emph{The parameter gap.} No quantitative relationship among $L$,
  $\sigma$, network delay, and honest hash and stake exists; the
  prototype's $\sigma = 3$, $L = 7$ carry an explicit no-verification
  warning, against a rollback bound that is itself only node policy ---
  1000 blocks in mainline Zebra at \texttt{9f3166b}, 99 in the
  zebra-crosslink fork (\S\ref{sec:xl-impl-divergences}).
\item \emph{The Stalled Mode question.} Whether the eventual ZIP adopts
  the book's bounded availability or the implementation's
  redelegation-only stance, and whether dropping the rule reopens the
  unbounded-gap and tail-thrashing hazards the book used to justify it
  (\S\ref{sec:xl-impl-divergences}). Trigger: a filed Crosslink ZIP.
\item \emph{The BFT engine.} No analysis shows tenderlink --- or any
  concrete $\Pi_{\mathrm{origbft}}$ candidate --- satisfies Final
  Agreement, objective validity, and sound notarization
  (\S\ref{sec:xl-impl-divergences}).
\item \emph{Machine-checking.} Coverage stops at
  $\mathsf{isValid}$, with one flagged divergence from the book
  (Remark~\ref{rem:xl-quint-div}); the bc-side rules, adversarial
  settings, and liveness are unchecked. Trigger: any extension of the
  Quint spec or content in \texttt{quint-crosslink}.
\item \emph{Accountability machinery.} On-chain evidence of BFT safety
  violations remains a non-adopted [Recommended] change
  (\S\ref{sec:xl-bounded}); long-range key hygiene improvements remain
  tfl-book issue \#140 (Remark~\ref{rem:xl-longrange}).
\item \emph{Audit follow-through.} Adoption of the mechanism-design
  recommendations, and survival of the no-slashing model against phantom
  security, are unresolved (\S\ref{sec:xl-econ}). Shielded Labs states
  further reviews by established firms precede activation.
\item \emph{Design-source discrepancies.} The staking epoch is roughly
  seven days in the nutshell (a one-day staking day plus a six-day locked
  phase) but five days in the 2025-10-27 blog post and the audit
  proposal; which is current is unrecorded
  (\texttt{nutshell.md}, lines 64--80; \texttt{proposal.md}, line 58).
\end{enumerate}

%% TODO: verify whether simtfl (github.com/zcash/simtfl), the simulator
%% the tfl-book references, produced any published analysis; not checked
%% in this pass.
%% TODO(assembly): the tfl-book's Crosslink 1 rule names in
%% potential-changes.md (Increasing Score rule etc.) predate Crosslink 2
%% and the chapter truncates mid-sentence; do not cite them as current
%% design anywhere in this volume (potential-changes.md, lines 5-8 and
%% 371-372).

\end{document}
